\documentclass[11pt]{amsart}

\usepackage[margin=1.1in]{geometry}
\usepackage{amsmath,amssymb,amsthm,mathtools}
\usepackage{booktabs,longtable,array}
\usepackage[shortlabels]{enumitem}
\usepackage{microtype}
\usepackage{xcolor}
\usepackage{colortbl}
\usepackage{tikz}
\usetikzlibrary{arrows.meta,calc}
\definecolor{ourcolumn}{gray}{0.93}
\usepackage[colorlinks=true,linkcolor=blue!50!black,citecolor=green!35!black,urlcolor=blue!60!black]{hyperref}
\usepackage{aliascnt}
\usepackage[capitalise,nameinlink]{cleveref}
\hypersetup{pdftitle={Linnik's constant is at most 3.99},pdfauthor={Eric Naslund}}

\numberwithin{equation}{section}

\theoremstyle{plain}
\newtheorem{theorem}{Theorem}[section]
\newaliascnt{proposition}{theorem}\newtheorem{proposition}[proposition]{Proposition}\aliascntresetthe{proposition}
\newaliascnt{lemma}{theorem}\newtheorem{lemma}[lemma]{Lemma}\aliascntresetthe{lemma}
\newaliascnt{corollary}{theorem}\newtheorem{corollary}[corollary]{Corollary}\aliascntresetthe{corollary}
\newaliascnt{inp}{theorem}\newtheorem{inp}[inp]{Input}\aliascntresetthe{inp}
\theoremstyle{definition}
\newaliascnt{definition}{theorem}\newtheorem{definition}[definition]{Definition}\aliascntresetthe{definition}
\newaliascnt{assumption}{theorem}\newtheorem{assumption}[assumption]{Assumption}\aliascntresetthe{assumption}
\theoremstyle{remark}
\newaliascnt{remark}{theorem}\newtheorem{remark}[remark]{Remark}\aliascntresetthe{remark}
\newaliascnt{example}{theorem}\newtheorem{example}[example]{Example}\aliascntresetthe{example}

\crefname{theorem}{Theorem}{Theorems}
\crefname{proposition}{Proposition}{Propositions}
\crefname{lemma}{Lemma}{Lemmas}
\crefname{corollary}{Corollary}{Corollaries}
\crefname{inp}{Input}{Inputs}
\crefname{definition}{Definition}{Definitions}
\crefname{assumption}{Assumption}{Assumptions}
\crefname{remark}{Remark}{Remarks}
\crefname{example}{Example}{Examples}

% ---------------------------------------------------------------- notation
\newcommand{\LL}{\mathcal{L}}
\newcommand{\RR}{\mathbb{R}}
\newcommand{\CC}{\mathbb{C}}
\newcommand{\NN}{\mathbb{N}}
\newcommand{\ZZ}{\mathbb{Z}}
\newcommand{\QQ}{\mathbb{Q}}
\renewcommand{\Re}{\operatorname{Re}}
\renewcommand{\Im}{\operatorname{Im}}
\newcommand{\cB}{\mathcal{B}}
\newcommand{\cR}{\mathcal{R}}
\newcommand{\cS}{\mathcal{S}}
\newcommand{\eps}{\varepsilon}
\newcommand{\lo}{\mathrm{lo}}
\newcommand{\hi}{\mathrm{hi}}
\newcommand{\1}{\mathbf{1}}
\newcommand{\bchi}{\overline{\chi}}
\DeclareMathOperator{\ord}{ord}
\DeclareMathOperator{\supp}{supp}

\begin{document}

\title{Linnik's constant is at most $3.99$}

\author[Eric Naslund]{Eric Naslund\textsuperscript{$*$}}
\email{naslund.math@gmail.com}

\date{September 30, 2026}

\subjclass[2020]{11N13 (primary); 11M20, 11M26, 11N35, 11Y35, 68V15, 90C05 (secondary)}
\keywords{Least prime in an arithmetic progression, Linnik's theorem, zeros of Dirichlet
$L$-functions, zero density, Selberg sieve, computer-assisted proof}

\begin{abstract}
Let $P(a,q)$ denote the least prime congruent to $a$ modulo $q$, where $(a,q)=1$.
We give a computer-assisted proof that $P(a,q)\ll q^{3.99}$, improving the exponent $5$
of Xylouris. More precisely, $P(a,q)<q^{3.99}$ for every sufficiently large $q$, uniformly
in $a$.

The main new ingredient is a graded near-density estimate for zeros of Dirichlet
$L$-functions, a weighted form of Heath-Brown's Lemma~12.1 of the kind he asked for in 1992. We
combine it with the zero-location estimates and far-density method of Heath-Brown and Xylouris.

The second novelty is the scale of the case analysis. Heath-Brown and Xylouris closed their final
case analyses with $14$ and $21$ main cases, each by a chain of inequalities evaluated in floating
point. In this paper, with the help of advanced AI models, we can push this much further. The middle range of the first zero is divided into $4453$ root cases and $4788$
terminal cases, each closed by its own linear program, and exact integer dual certificates on
$4{,}196{,}879$ threshold boxes bound the normalized zero sum strictly below~$1$. Exceptional
zeros and the remaining exterior range are treated separately. This case analysis is a kind of
systematic brute force enabled by AI: an analysis of this size, with parameters tuned to each case,
would be very laborious or nearly impossible to carry out by hand, and it lets the new estimate be
applied separately in each case.

The near-density lemma, certificate soundness, and a conditional passage from a certified
case to a prime are formalized in Lean, assuming published analytic inputs and specified
facts about zeros.
\end{abstract}

\maketitle

\begin{quote}
\small\noindent\textsuperscript{$*$} This paper was written by Claude Opus 5.5 and GPT Astra 6 under my prompting and supervision. The research was carried out by the same models under my prompting and supervision, with some of the formalization work also done by xAI's Grok. My recommendation is for the reader to explore this paper with AI as well, and use their favorite model to ask the paper questions. The certificates are checked by three separate programs, one of them formally verified, and
\cref{sec:verification} states precisely which parts of the proof are formally verified in Lean and which have conventional proofs only.
\end{quote}

\clearpage
\setcounter{tocdepth}{1}
\begingroup
\small
\tableofcontents
\endgroup

% !TEX root = ../linnik399.tex
\section{Introduction}\label{sec:intro}

\subsection{The result}
For an integer $q\ge2$ and an integer $a$ with $(a,q)=1$, let $P(a,q)$ be the least prime
congruent to $a$ modulo $q$. Dirichlet's theorem \cite{Dirichlet1889beweis} ensures that this
prime exists. Linnik \cite{Linnik1944least,Linnik1944deuring} proved that there are absolute
constants $C$ and $L$ such that
\begin{equation}\label{eq:linnik}
  P(a,q)\le Cq^L
\end{equation}
for every reduced residue class. An exponent $L$ for which such a constant $C$ exists is
called \emph{admissible}; Linnik's constant is the infimum of the admissible exponents.

\begin{theorem}[Least prime in a reduced residue class]\label{thm:main}
There is an absolute constant $q_0\ge2$ such that, for every integer $q\ge q_0$ and every
integer $a$ with $(a,q)=1$, there is a prime $p$ satisfying
\[
  p\equiv a\pmod q,\qquad p<q^{3.99}.
\]
Consequently, there is an absolute constant $C>0$ such that
$P(a,q)\le Cq^{3.99}$ for all $q\ge2$ and $(a,q)=1$. In particular, Linnik's constant is at
most $3.99$.
\end{theorem}

We do not compute $q_0$ or $C$; the issue of effectivity is discussed at the end of
\cref{sec:join}. Except in the presence of an extremely close real zero, the proof finds a
prime in the smaller interval
\[
  q^{3.156342}<p<q^{3.99}.
\]
The lower endpoint comes from the support of the weight used to detect primes. When a
real zero is extremely close to $1$, Heath-Brown's theorem on Siegel zeros
\cite{HeathBrown1990siegel} instead gives $P(a,q)\le q^{3.5}$.

The argument combines published analytic estimates, new lemmas proved here, and exact
finite computations. The theorem requires all three parts. In particular, checking a
numerical certificate proves a bound for its specified configuration of zeros; the
analytic argument must also show that every modulus gives a configuration covered by
one of the certificates. We make this passage explicit in \cref{sec:casetree,sec:join}.

\subsection{Earlier work}
The admissible exponents in \cref{tab:history} record the main numerical developments.
Heath-Brown \cite{HeathBrown1992zero} obtained $5.5$, and Xylouris
\cite{Xylouris2009linnik,Xylouris2011least,Xylouris2011nullstellen,Xylouris2018linnik}
subsequently obtained $5.2$, $5.18$, and $5$. Before them, Pan \cite{Pan1957least,Pan1958least} gave
the first numerical values, and Chen, Jutila, Graham and Wang lowered them in a long series of
papers, which includes Graham's introduction of Selberg sieve weights into zero-density estimates
\cite{Graham1977thesis,Graham1981linnik}. Our proof uses the framework of these works. It rests on
the three principles that Heath-Brown isolates at the start of his paper: a zero-free region
with at most one exceptional zero
\cite{delaValleePoussin1896recherches,Gronwall1913series,Landau1918imaginar,Page1935number}, the
repulsion of other zeros by an exceptional zero, known as the Deuring--Heilbronn phenomenon
\cite{Deuring1933imaginare,Heilbronn1934class,Linnik1944deuring}, and a log-free zero-density
estimate \cite{Linnik1944least}. The last means that the density bound has no additional power of
$\log q$, a feature needed when studying zeros at distance $O(1/\log q)$ from $1$. Simpler proofs
of these principles, and of Linnik's theorem, were given by Rodosski\u\i{} \cite{Rodosskii1954least},
Tur\'an \cite{Turan1961density,Turan1984new}, Knapowski \cite{Knapowski1962linnik}, Fogels
\cite{Fogels1965zeros}, Gallagher \cite{Gallagher1970large}, Jutila \cite{Jutila1972density},
Motohashi \cite{Motohashi1975density,Motohashi1978primes} and Bombieri \cite{Bombieri1987grand},
among others; accounts are in the books \cite{Prachar1957primzahlverteilung,Montgomery1971topics,
IwaniecKowalski2004analytic,FriedlanderIwaniec2010opera}.

The conjectured scale is much smaller. Chowla \cite{Chowla1934least} conjectured
$P(a,q)\ll_\eps q^{1+\eps}$ for every $\eps>0$; the maximum over reduced classes is
conjecturally of order $\varphi(q)(\log q)^2$
\cite{Wagstaff1979greatest,Pomerance1980note,GranvillePomerance1990least,LiPrattShakan2017lower}.
The Generalized Riemann Hypothesis gives admissible exponents $2+\eps$; see
\cite{BachSorenson1996explicit,LamzouriLiSoundararajan2015conditional,CarneiroMilinovichQuesadaHerreraRamos2025fourier}
for explicit conditional bounds. Explicit \cite{BennettMartinOBryantRechnitzer2018explicit} and
uniform \cite{ThornerZaman2024refinements} unconditional estimates for primes in progressions are
also known. In the other direction, an exceptional (Siegel) zero helps: Heath-Brown
\cite{HeathBrown1990siegel} showed that a sufficiently strong exceptional zero forces
$P(a,q)\le q^{3+\delta}$, in an effective form that we use below, and $q^{2+\delta}$ ineffectively; see also
\cite{FriedlanderIwaniec2003exceptional,Iwaniec2006conversations,Zaman2016least}. Other proofs of
Linnik's theorem use pretentious methods \cite{Sachpazis2023pretentious}, sieve methods
\cite{FriedlanderIwaniec2023sifting}, or avoid $L$-functions
\cite{MatomakiMerikoskiTeravainen2024primes}, but give larger exponents. For moduli with special
multiplicative structure there are stronger zero-free regions
\cite{BarbanLinnikChudakov1964prime,Gallagher1972primes,Iwaniec1974zeros}, and much smaller
exponents are known \cite{Chang2014short}.

\begin{table}[ht]
\centering
\begin{tabular}{llll}
\toprule
$L$ & Year & Author(s) & Reference \\
\midrule
$10000$ & 1957 & Pan (announced) & \cite{Pan1957least} \\
$5448$ & 1958 & Pan & \cite{Pan1958least} \\
$777$ & 1965 & Chen & \cite{Chen1965least} \\
$630$ & 1971 & Jutila (reported by Tur\'an) & \cite{Turan1971recent} \\
$550$ & 1970 & Jutila & \cite{Jutila1970new} \\
$168$ & 1977 & Chen & \cite{Chen1977least} \\
$80$ & 1977 & Jutila & \cite{Jutila1977linnik} \\
$36$ & 1977 & Graham & \cite{Graham1977thesis} \\
$20$ & 1981 & Graham & \cite{Graham1981linnik} \\
$17$ & 1979 & Chen & \cite{Chen1979least} \\
$16$ & 1986 & Wang & \cite{Wang1986least} \\
$13.5$ & 1989 & Chen and Liu & \cite{ChenLiu1989III,ChenLiu1989IV} \\
$11.5$ & 1991 & Chen and Liu & \cite{ChenLiu1991V} \\
$8$ & 1991 & Wang & \cite{Wang1991least} \\
$5.5$ & 1992 & Heath-Brown & \cite{HeathBrown1992zero} \\
$5.2$ & 2009 & Xylouris & \cite{Xylouris2009linnik} \\
$5.18$ & 2011 & Xylouris & \cite{Xylouris2011least} \\
$5$ & 2011 & Xylouris & \cite{Xylouris2011nullstellen,Xylouris2018linnik} \\
\midrule
$3.99$ & 2026 & this paper & \\
\bottomrule
\end{tabular}
\caption{Admissible values of Linnik's constant, following the tables in
\cite{HeathBrown1992zero} and \cite{Xylouris2011least}, in order of the bound. Graham's value $20$
was submitted before Chen's $17$ appeared.}
\label{tab:history}
\end{table}

\subsection{From zeros to primes}\label{subsec:outline}
Write $\LL=\log q$ and express a zero of a Dirichlet $L$-function as
\[
  \rho=1-\frac{\lambda}{\LL}+i\frac{\mu}{\LL}.
\]
Thus small $\lambda$ means that the zero is close to the line $\sigma=1$; $\mu$ is its
height multiplied by $\log q$. We group each nonprincipal character with its complex
conjugate and call this a \emph{family}. The parameters $\lambda_1,\lambda_2,\lambda_3$
record the first zero of successive families in a specified rectangle near $1$.
The parameter $\lambda'$ records the next zero occurrence in the first family, after
removing the distinguished zero and its conjugate where appropriate. Precise definitions,
including multiplicities and ties, are given in \cref{sec:notation}.

A nonnegative weight $h$, supported in $[A,L]$ with $A=L-2T$ and $T=0.416829$, detects
primes between $q^A$ and $q^L$. Let $H$ be its Laplace transform. Xylouris's positivity
criterion \cite[(3.57)]{Xylouris2011nullstellen} gives
\[
  \sum_{p\equiv a\,(q)}\frac{\log p}{p}\,h\Bigl(\frac{\log p}{\LL}\Bigr)
  \ge \frac{\LL H(0)}{\varphi(q)}\bigl(1-W-\eta\bigr),
\]
where $\eta=10^{-6}$ and
\[
  W=\frac1{H(0)}\sum_{\chi\ne\chi_0}\ \sum_{\rho\in R_P}
       \bigl|H\bigl((1-\rho)\LL\bigr)\bigr|.
\]
Here $R_P$ is a fixed square in the normalized coordinates, and zeros are counted with
multiplicity. It is therefore enough to prove $W\le1-2\eta$. The contribution of an
individual zero decays exponentially with $\lambda$, so zeros close to $1$ require the
most accurate estimates.

We combine three kinds of information about these zeros. First, a \emph{per-character
bound} estimates the total contribution of one $L$-function from a lower bound for the
parameters of its zeros. Second, a \emph{far-density bound} limits a weighted sum over
characters, with weight depending on a selected zero of each character. Third,
\emph{near-density bounds} constrain the characters whose selected zeros lie closest to
$1$. Zero-location estimates for $\lambda_1$, $\lambda'$, $\lambda_2$, and $\lambda_3$ determine
which of these bounds are available in a given case.

\subsection{The new near-density estimate}\label{subsec:newlemma}
Heath-Brown's near-density estimate \cite[Lemma~12.1]{HeathBrown1992zero} bounds the number
$N(\lambda)$ of characters with a zero in $\sigma\ge1-\lambda/\LL$, $|t|\le1$. Its proof applies
the Cauchy--Schwarz inequality to a prime sum at the single point $\beta_1=1-\lambda_1/\LL$, and
compares the resulting Gram form of the characters with the diagonal. Because every character
is tested at the same point, a zero at distance $\lambda$ counts only through the single number
$N(\lambda)$, and the bound is weaker than the density estimate of Heath-Brown's \S11 as soon as
$\lambda\ge1.1$ and void soon after. That density estimate, proved with Selberg sieve weights
after Graham, has the form (11.4) of a sum over characters of a weight that depends on each
character's own distance $\lambda^{(k)}$, and it contains the count (11.5) of the characters with
$\lambda^{(k)}\le\lambda$. In his closing list of possible improvements Heath-Brown writes
\cite[\S16, item~9]{HeathBrown1992zero}:
\begin{quote}
It would be nice to have a weighted version of Lemma~12.1, in the way that (11.4) is a
weighted version of (11.5). Unfortunately, no neat way of achieving this seems available.
\end{quote}

The graded near-density lemma (\cref{thm:graded}) supplies such an estimate. We test each
character at its own \emph{anchor}, a real parameter determining the point at which a
smoothed prime sum is evaluated. Cauchy--Schwarz relates these prime sums to a Gram form
of the characters. We estimate its off-diagonal terms by the local explicit formula and,
where sieve weights are used, Burgess's bound; Graham's sieve asymptotic controls the
diagonal terms. Allowing the anchors to vary is the grading in the name of the lemma.

After normalization, the resulting constraint has the form
\[
  \Bigl(\sum_j a_jv_j\Bigr)_+^2
  \le \sum_j D_ja_j^2+d\Bigl(\sum_j a_j\Bigr)^2
  \qquad(a_j\ge0).
\]
Here $x_+=\max(x,0)$, the index $j$ runs over character--height entries, $v_j$ is a lower
bound for the response of an entry, and $D_j>0$ and $d>0$ bound its diagonal and common
correlation costs. Most entries correspond to one selected zero; some retain a second zero
of the same character, with an explicit correction for their correlation.

\Cref{prop:threshold} turns this inequality into a form suited to computation: there is a
number $\tau\in[0,\sqrt d]$ such that
\[
  \sum_j\frac{(v_j-\tau)_+^2}{D_j}\le1-\frac{\tau^2}{d}.
\]
A single constraint now distinguishes zeros at different distances from $1$.
The sieve weighting has precedents in Graham's work
\cite{Graham1977thesis,Graham1981linnik}, Heath-Brown's far-density argument, and the
sieve-weighted large-sieve inequalities of Motohashi \cite{Motohashi1975density}.

\emph{Where the new lemma enters.} In Heath-Brown's proof the near-density estimate enters at the very end, in the assembly of
\cite[\S15]{HeathBrown1992zero}. There the characters with a zero below a level $\Lambda$ are sorted into
bins $(\Lambda_{r+1},\Lambda_r]$, with $\Lambda_r=\Lambda-0.025r$; each character is charged the cost of the end of its
bin nearer to $s=1$; and the numbers of characters in the bins are controlled only through the
counts $N(\Lambda_r)$, bounded one level at a time by Lemma~12.1 (his Table~13). The far-density lemma
controls the characters beyond $\Lambda$. Xylouris's assembly \cite[\S6]{Xylouris2011nullstellen} has the same
structure, with a sharper form of Lemma~12.1 (his Lemma~5.3). A count $N(\lambda)$ treats all characters
with a zero below $\lambda$ alike, so it cannot tell a few zeros very close to $s=1$ from many zeros
further away. In our proof this place is taken by the \emph{near rows} of the leaf programs
described below. Each near row is one instance of \cref{thm:graded}: a single quadratic constraint
in which every bin of characters enters with its own feature, computed at its own anchor. The
weighted version of Lemma~12.1 that Heath-Brown asked for is thus used exactly where his count was
used, in every leaf program, and it is the main source of the improvement
(\cref{subsec:compare}).

\subsection{The finite case analysis}
We divide the middle range $0.1\le\lambda_1<1.5$ into cases determined by the type of the
first family and intervals for the first few zero parameters. There are four levels in
the computation, serving different purposes.
\begin{enumerate}[(1)]
  \item A \emph{root case} specifies a range of possible zero configurations.
    The $58$ parent rows taken from published zero-location tables lead to $4453$
    root cases, covering the middle range.
  \item A \emph{leaf} is a terminal case after further subdivision and valid
    zero-location exclusions. There are $4788$ leaves that require numerical bounds.
  \item A \emph{leaf program} groups characters into bins according to their selected
    zero parameters. Its nonnegative variables are the numbers of characters in these
    bins, together with weighted masses for the tails. The objective bounds $W$, and
    its constraints express the far-density bound, restrictions on the first families,
    and the near-density inequalities. For fixed thresholds these constraints are linear.
  \item A \emph{threshold box} bounds the auxiliary numbers $\tau$ in the near rows.
    On each box, a linear relaxation is either excluded or bounded by exact dual
    multipliers. Subdividing these boxes covers every possible choice of thresholds.
\end{enumerate}
The case tree partitions information about zeros; the threshold boxes partition the
auxiliary parameters used to bound them. Keeping these two subdivisions distinct is
essential to the coverage argument.

Each leaf, or each subcase of a leaf, has its own certificate, a binary tree of threshold
boxes. The $4788$ leaves give $5591$ such certificate trees, because some leaves are
subdivided further, and these trees contain $4{,}196{,}879$ threshold boxes. The largest certified value is
$0.9999998223\ldots<1$; it bounds the zero sum together with the required error allowance.
All certificate inequalities are checked in integer arithmetic. The tree was constructed
for an earlier computation at exponent $4.30$. Its subdivisions and zero-location
arguments are independent of the target exponent, so we retain them and certify new
leaf programs at $3.99$. \Cref{tab:scale} compares these numbers with the case analyses of
Heath-Brown and Xylouris, and \cref{tab:census} lists every level of the analysis.

The exterior ranges are shorter arguments. An exceptionally small $\lambda_1$ is covered
by Heath-Brown's theorem on Siegel zeros. The rest of $\lambda_1\le0.1$ is handled by his
quantitative Deuring--Heilbronn estimates and zero-location tables. A single further
certificate treats $\lambda_1\ge1.5$. If the relevant rectangle contains no zero, then
$W=0$ and the positivity criterion applies immediately.

\subsection{Comparison with Heath-Brown and Xylouris}\label{subsec:compare}

Our framework is Heath-Brown's \cite{HeathBrown1992zero} with Xylouris's refinements
\cite{Xylouris2011nullstellen,Xylouris2018linnik}. We use the following from their work:
\begin{itemize}
  \item the positivity criterion \cite[(3.57)]{Xylouris2011nullstellen};
  \item the per-character bound \cite[Lemma~3.10]{Xylouris2011nullstellen};
  \item the far-density lemma \cite[Lemma~5.1]{Xylouris2011nullstellen};
  \item the zero-free regions and Deuring--Heilbronn estimates, as printed in their tables;
  \item the local explicit formulas \cite[Lemmas~3.1, 3.2]{Xylouris2011nullstellen}, which are
    \cite[Lemmas~5.3, 5.2]{HeathBrown1992zero}.
\end{itemize}

\emph{The case analyses compared.} Both earlier proofs are computer-assisted, and both are
organized, like ours, as a case analysis on the first zeros; \cref{tab:scale} compares the three.
Heath-Brown's paper has twelve numerical tables of zero-location and density bounds. Its final
assembly \cite[\S15]{HeathBrown1992zero} treats $\lambda_1\ge0.348$ in fourteen intervals, mostly of length
$0.02$, each by a chain of inequalities evaluated numerically; the largest total it reports is
$0.9943$. Heath-Brown writes \cite[\S1]{HeathBrown1992zero} that ``the arguments rely heavily on
numerical calculations. These we do not reproduce in full, nor have we attempted any rigorous
analysis of the rounding and truncation errors in the computer algorithms employed.'' Xylouris's
dissertation adds eleven tables. Its final assembly \cite[\S6.4, pp.~86--87]{Xylouris2011nullstellen}
has $21$ cases: three intervals of $\lambda_1$ for a real zero of a real character, two for a complex
zero of a real character, and sixteen for a nonreal character. Each case is split further
according to the possible values of two zero-counting functions
\cite[(6.46)]{Xylouris2011nullstellen}, and the computations were done in Maple with ten digits
\cite[p.~14]{Xylouris2011nullstellen}; the largest value is $W<0.998$. In both, each case is closed by a
chain of inequalities in which the unknown zeros are estimated separately, one after another.

In our proof the same role is played by $4453$ root cases and $4788$ leaves, and each leaf is
closed by a linear program with an exact certificate instead of a floating-point evaluation. The
$4.2$ million boxes have no counterpart in the earlier work. They do not subdivide the
configurations of zeros. They subdivide the auxiliary thresholds of the near rows, and they arise
because the graded lemma is a quadratic constraint rather than a count. The largest certified
value, $0.99999982$ against $0.9943$ and $0.998$, shows how much finer the analysis must be at
$L=3.99$.

\begin{table}[tb]
\centering
\small
\setlength{\tabcolsep}{6pt}
\renewcommand{\arraystretch}{1.3}
\begin{tabular}{@{}>{\raggedright\arraybackslash}p{3.05cm}
                >{\centering\arraybackslash}p{3.25cm}
                >{\centering\arraybackslash}p{3.25cm}
                >{\columncolor{ourcolumn}\centering\arraybackslash}p{3.65cm}}
\toprule
 & \textbf{Heath-Brown} \cite{HeathBrown1992zero} & \textbf{Xylouris} \cite{Xylouris2011nullstellen} & \textbf{This paper}\\
 & $L=5.5$ \ (1992) & $L=5$ \ (2011) & $L=3.99$\\
\midrule
Near-density input & counts $N(\lambda)$ from Lemma~12.1 & counts $N(\lambda)$ from a sharper Lemma~12.1 &
  the weighted (graded) lemma, \cref{thm:graded}\\
Numerical tables & $12$ & $11$ more & $58$ rows from theirs, plus over $100$ new\\
Final case analysis & $14$ intervals of $\lambda_1$ & $21$ cases, each split by two zero counts &
  $4453$ roots, $4788$ leaves\\
How a case is closed & floating point, without rounding analysis & Maple, $10$-digit floating point &
  exact certificates on $4.2$~million boxes\\
Largest value of $W$ (must be below $1$) & $0.9943$ & $0.998$ & $0.99999982$\\
\bottomrule
\end{tabular}
\caption{Three proofs of an explicit Linnik constant compared, in the middle range of the first
zero ($\lambda_1\ge0.348$ for Heath-Brown and Xylouris, $0.1\le\lambda_1<1.5$ here). The last row is the largest
bound for the normalized zero sum $W$ over all cases of the final assembly: for Heath-Brown the
largest total he reports, for Xylouris his stated bound, and for this paper the largest certified
value, a bound for $W+2\eta$. The entry ``without rounding analysis'' paraphrases Heath-Brown's own
remark \cite[\S1]{HeathBrown1992zero}, quoted above.}
\label{tab:scale}
\end{table}

\emph{Where the gain comes from.} The gain from $5$ to $3.99$ has three sources:
\begin{itemize}
  \item the graded near-density lemma, which replaces the counts $N(\lambda)$ and applies to every
    range of $\lambda$ at once;
  \item the linear-programming combination, which lets all constraints act together on each
    case instead of chaining worst cases;
  \item a fine case analysis of the first zeros, possible because every case is certified by
    machine.
\end{itemize}
The last source alone does not go far. Xylouris estimates that his method, with much larger
computations (smaller intervals and finer grids), would give about $4.96$ \cite[p.~82]{Xylouris2018linnik}.
The earlier form of our method, with the linear programs and a fine case analysis but without the
graded lemma, reached $4.30$, using several ad hoc near-density models. The graded lemma replaces
almost all of them, and it is the main reason for the further gain to $3.99$.

\subsection{The certified exponent and numerical evidence}\label{subsec:why}
We state the result at $3.99$ because this is the exponent for which the complete cover
has been certified. The tightest cases have a nonreal first character and
$\lambda_1\in[0.64,0.80]$. In this range the available lower bounds for $\lambda_2$ are
comparatively small, which restricts the anchors in the near-density estimates.
Several of the corresponding leaf bounds are within $10^{-7}$ of $1$ after the error
allowance, and require fine subdivisions.

Exploratory floating-point computations suggest that a somewhat smaller exponent may
be accessible. One difficult leaf, with the reserved-family columns of
\cref{subsec:leafmodel}, has computed values $0.992$ at $L=3.97$ and $1.020$ at $L=3.95$.
These are values of particular relaxations, not certified global bounds or an obstruction
to other methods. Similarly, the positive numerical margins for the exterior ranges at
smaller exponents do not constitute a complete proof there. The verification reported in
this paper is at $L=3.99$.

\subsection{Verification and a guide to the proof}\label{subsec:checked}
The certificates have been checked by independent programs. Parts of the analytic
argument and the soundness of the checkers have also been formalized in Lean, with the
published analytic inputs stated as hypotheses. This formalization includes a theorem
that a valid certified leaf yields a prime when the modulus satisfies that leaf's
hypotheses about zeros. It does not include the full passage from arbitrary moduli to
leaves, or the exterior arguments. \Cref{sec:verification} gives the verification records
and separates these claims precisely.

\Cref{sec:notation,sec:inputs} give the notation and published inputs.
\Cref{sec:detection,sec:costs,sec:far} reduce prime detection to bounds for the zero sum,
and \cref{sec:location} supplies the zero-location estimates. The main new analytic
argument is in \cref{sec:near}; \cref{sec:rows} derives the constraints used in the
computation. \Cref{sec:lp} proves certificate soundness, while \cref{sec:casetree} shows
that the zero configurations are covered and satisfy the programs. The exterior cases
and the final choice of constants are treated in \cref{sec:exteriors,sec:join}.
A reader interested first in the new density estimate may read \cref{sec:notation,sec:inputs}
and then \cref{sec:near}, returning to the numerical construction afterwards.

\Cref{fig:proofmap} shows how the main results depend on one another, and which of them are
published inputs, new results, computations or formally verified statements. The terms of the
case analysis are introduced in \cref{subsec:vocabulary}, and the fixed constants and the error
budget in \cref{subsec:budget}. \Cref{tab:census} in \cref{sec:casetree} lists all levels of the
finite analysis, and \cref{tab:status} in \cref{sec:verification} records, for each component,
its source and how it was checked.

\begin{figure}[tbp]
\centering
\begin{tikzpicture}[
  x=1cm, y=1cm, font=\footnotesize,
  box/.style={rectangle, rounded corners=1.5pt, align=center, inner sep=3pt,
    minimum height=8.5mm, minimum width=17mm, line width=0.5pt},
  pub/.style={box, fill=black!12, draw=black!55},
  new/.style={box, fill=white, draw=black},
  comp/.style={box, fill=white, draw=black, dash pattern=on 2.2pt off 1.4pt},
  lean/.style={double, double distance=1pt, line width=0.45pt},
  dep/.style={-{Stealth[length=3.4pt,width=2.8pt]}, draw=black!75, line width=0.45pt},
  wire/.style={draw=black!75, line width=0.45pt},
]
% ---------------------------------------------------------------- theorem and regimes
\node[new] (main) at (8,0)
  {\textbf{Theorem~\ref{thm:main}}\\ $P(a,q)<q^{3.99}$ for $q\ge q_0$};
\node[new] (tiny) at (1.2,-1.8)
  {(R2) $\lambda_1<u_{\mathrm{exc}}$:\\ $P(a,q)\le q^{3.5}$\\ Prop.~\ref{prop:tiny}};
\node[new] (small) at (4.2,-1.8)
  {(R2)\\ $u_{\mathrm{exc}}\le\lambda_1<0.1$\\ Thm.~\ref{thm:small}};
\node[new] (r1) at (6.75,-1.8)
  {(R1) no zero\\ in $R(l)$,\\ so $W=0$};
\node[new, lean] (detect) at (9.7,-1.8)
  {$W\le1-2\eta$\\ gives a prime\\ Prop.~\ref{prop:detect}};
\node[pub] (crit) at (13.3,-1.8)
  {Xylouris's\\ criterion\\ Input~\ref{inp:criterion}};
\node[pub] (siegel) at (1.2,-3.6)
  {Siegel zeros\\ (Heath-Brown)\\ Input~\ref{inp:siegel}};
\node[pub] (hb) at (4.2,-3.6)
  {Heath-Brown's\\ lemmas\\ Inputs~\ref{inp:hb132}--\ref{inp:hb111}};
\node[comp] (marg) at (6.75,-3.6)
  {Interval\\ margins};
\node[new] (r3) at (9.7,-3.6)
  {(R3) $0.1\le\lambda_1<1.5$\\ Thm.~\ref{thm:cover}};
\node[new] (r4) at (12.6,-3.6)
  {(R4) $\lambda_1\ge1.5$\\ Thm.~\ref{thm:large}};
% ---------------------------------------------------------------- the middle range
\node[new] (src) at (1.0,-6.3)
  {4453 root\\ cases cover\\ Thm.~\ref{thm:sourcecover}};
\node[comp] (tree) at (3.3,-6.3)
  {Refinement\\ trees, \S\ref{subsec:refinement}\\ 4788 leaves};
\node[comp, lean] (certs) at (5.9,-6.3)
  {Exact\\ certificates:\\ $4{,}196{,}879$ boxes};
\node[new] (real) at (8.6,-6.3)
  {Realization:\\ feasible points\\ Prop.~\ref{prop:realization}};
\node[new, lean] (certthm) at (11.6,-6.3)
  {Certificate\\ soundness\\ Thm.~\ref{thm:certificate}};
\node[comp, lean] (largecert) at (14.4,-6.3)
  {One\\ certificate\\ (36 intervals)};
% ----
\node[new] (parent) at (1.0,-8.5)
  {58 parent\\ rows\\ Prop.~\ref{prop:parentrows}};
\node[new] (newloc) at (3.5,-8.5)
  {New location\\ bounds\\ Props.~\ref{prop:realrows}--\ref{prop:third}};
\node[new] (costs) at (6.3,-8.5)
  {Per-character\\ costs, \S\ref{sec:costs}};
\node[new] (farb) at (8.8,-8.5)
  {Far budget,\\ representatives\\ \S\ref{sec:far}};
\node[new] (rows) at (11.4,-8.5)
  {Near rows\\ hold\\ Prop.~\ref{prop:rowsvalid}};
% ----
\node[pub] (tables) at (1.3,-10.6)
  {Published\\ zero tables\\ Inputs~\ref{inp:rrtables}--\ref{inp:lambda3}};
\node[pub] (env) at (5.9,-10.6)
  {Envelope lemma\\ (Xylouris)\\ Input~\ref{inp:envelope}};
\node[pub] (farin) at (8.9,-10.6)
  {Far density\\ (Xylouris)\\ Input~\ref{inp:far}};
\node[new, lean] (zt) at (11.7,-10.6)
  {Threshold form\\ Prop.~\ref{prop:threshold},\\ Cor.~\ref{cor:zeroform}};
\node[new] (twotest) at (14.4,-10.6)
  {Two-test row\\ Thm.~\ref{thm:twotest}};
% ----
\node[new, lean] (response) at (10.1,-12.6)
  {Response lemma\\ Lem.~\ref{lem:response}};
\node[new, lean] (graded) at (13.5,-12.6)
  {\textbf{Graded}\\ \textbf{near-density lemma}\\ Thm.~\ref{thm:graded}};
\node[pub] (expl) at (10.3,-14.5)
  {Explicit formulas\\ Inputs~\ref{inp:X31}, \ref{inp:X32}};
\node[pub] (burgra) at (13.7,-14.5)
  {Burgess, Graham\\ Inputs~\ref{inp:burgess}, \ref{inp:graham}};
% ---------------------------------------------------------------- arrows
% both parts of (R2) give primes directly
\draw[wire] (tiny.north) |- (main.west);
\draw[dep] (small.north) |- (main.west);
\draw[dep] (detect.north) -- (detect.north |- main.south);
\draw[dep] (r1) -- (detect);
\draw[dep] (crit) -- (detect);
\draw[dep] (r3) -- (detect);
\draw[dep] (r4) -- (detect);
\draw[dep] (siegel) -- (tiny);
\draw[dep] (hb) -- (small);
\draw[dep] (marg) -- (small);
% the ingredients of the cover theorem, joined
\coordinate (bus) at ($(r3.south)+(0,-0.6)$);
\foreach \nd in {src,tree,certs,real}
  \draw[wire] (\nd.north) -- (\nd.north |- bus);
\draw[wire] ([xshift=-3mm]certthm.north) -- ([xshift=-3mm]certthm.north |- bus);
\draw[wire] (src.north |- bus) -- ([xshift=-3mm]certthm.north |- bus);
\draw[dep] (bus) -- (r3.south);
\draw[dep] ([xshift=4mm]certthm.north) -- ([xshift=4mm]certthm.north |- r4.south);
\draw[dep] (largecert) -- (r4);
% below
\draw[dep] (parent) -- (src);
\draw[dep] (newloc) -- (tree);
\draw[dep] (costs) -- (real);
\draw[dep] (farb) -- (real);
\draw[dep] (rows) -- (real);
\draw[dep] (tables) -- (parent);
\draw[dep] (tables) -- (newloc);
\draw[dep] (env) -- (costs);
\draw[dep] (farin) -- (farb);
\draw[dep] (zt) -- (rows);
\draw[dep] (twotest) -- (rows);
\draw[dep] (response) -- (zt);
\draw[dep] (graded) -- (zt);
\draw[dep] (expl) -- (response);
\draw[dep] (expl) -- (graded);
\draw[dep] (burgra) -- (graded);
% ---------------------------------------------------------------- legend
\node[pub, minimum width=7.5mm, minimum height=4.6mm] (k1) at (0.55,-12.3) {};
\node[new, minimum width=7.5mm, minimum height=4.6mm] (k2) at (0.55,-12.95) {};
\node[comp, minimum width=7.5mm, minimum height=4.6mm] (k3) at (0.55,-13.6) {};
\node[new, lean, minimum width=7.5mm, minimum height=4.6mm] (k4) at (0.55,-14.25) {};
\node[anchor=west] at ([xshift=1mm]k1.east) {published input};
\node[anchor=west] at ([xshift=1mm]k2.east) {result proved in this paper};
\node[anchor=west] at ([xshift=1mm]k3.east) {computation};
\node[anchor=west] at ([xshift=1mm]k4.east) {formally verified in Lean};
\end{tikzpicture}
\caption{Proof map for \cref{thm:main}. An arrow points from a result to a result whose
proof uses it; the joined lines into (R3) are the ingredients of \cref{thm:cover}. The
regimes (R1)--(R4) of \cref{sec:join} are defined by the first zero $\lambda_1$, and (R2)
splits at the constant $u_{\mathrm{exc}}$ of \cref{sec:exteriors}. A double border marks a
statement formally verified in Lean, with the published inputs as hypotheses, or, for the
certificates, a computation accepted by a checker proved sound in Lean. The formalization
also covers the far budget and parts of \cref{prop:rowsvalid,prop:realization}; the
zero-location results, the case tree, the costs, the two-test row, the exterior regimes and
the assembly are not formally verified (\cref{subsec:lean}). Some arrows are omitted, for
instance from the explicit formulas to \cref{sec:costs,sec:location} and into
\cref{thm:large}.}
\label{fig:proofmap}
\end{figure}


% !TEX root = ../linnik399.tex
\section{Notation and conventions}\label{sec:notation}

We use standard notation for Dirichlet characters and their $L$-functions. For background on
Dirichlet $L$-functions and their zeros we refer to
\cite{Ingham1932distribution,Titchmarsh1986theory,Karatsuba1993basic,Montgomery1994ten,
Davenport2000multiplicative,IwaniecKowalski2004analytic,MontgomeryVaughan2007multiplicative,
Tenenbaum2015introduction}, and for the history of the subject to \cite{Narkiewicz2000development}. The distinctions between a zero and a zero occurrence, and between
normalized and physical height, will be used throughout the proof.

\subsection{Moduli, characters and zeros}
Throughout, $q\ge2$ is an integer modulus and
\[
  \LL=\log q .
\]
Characters are Dirichlet characters modulo $q$; $\chi_0$ is the principal character. For a
nonprincipal character $\chi$ we consider the zeros $\rho=\beta+i\gamma$ of $L(s,\chi)$ with
$\beta>0$, counted with multiplicity, and we write
\begin{equation}\label{eq:normalized}
  \rho=1-\frac{\lambda}{\LL}+i\frac{\mu}{\LL},\qquad \lambda=(1-\beta)\LL,\quad \mu=\gamma\LL .
\end{equation}
We call $\lambda=\lambda_\rho$ the \emph{parameter} of $\rho$ and $\mu=\mu_\rho$ its
\emph{normalized height}; $\gamma=\gamma_\rho$ is its \emph{physical height}. Thus a height
bound always concerns $|\gamma|$ or $|\mu|$. Smaller $\lambda$ means a zero farther to the
right. The multiplicity of $\rho$ as a zero of the specified $L(s,\chi)$ is $m_\rho$.
A \emph{zero occurrence} is one copy of the pair $(\chi,\rho)$: a zero of multiplicity
$m_\rho$ gives $m_\rho$ occurrences. Unless stated otherwise, sums over zeros count
multiplicity. A zero of an imprimitive character is a zero of the
$L$-function of the primitive character inducing it, with the same multiplicity, as long as
$\beta>0$; all zeros considered below have $\beta>1/2$.

\subsection{Regions}
Following Xylouris \cite[(3.7)]{Xylouris2011nullstellen} we put, for $x>0$,
\[
  R(x)=\Bigl\{\sigma+it:\ 1-\frac{\log\log\LL}{3\LL}\le\sigma\le1,\ |t|\le x\Bigr\}.
\]
There is an integer $l=l(q)$ with $1\le l\le\LL/10$ such that no zero of any $L(s,\chi)$ lies in
$R(10l)\setminus R(l)$ (\cref{inp:rectangle}). We fix one such choice, for example the least
admissible integer. The distinguished zeros (\cref{def:additionalzero}) and the representatives
(\cref{subsec:vocabulary}) are selected in $R(l)$.
The local explicit formula also involves disc zeros, which need not all belong to this
rectangle. Three further regions have separate roles:
\begin{itemize}
  \item the \emph{prime window} $R_P=\{0\le\lambda\le C_P,\ |\mu|\le C_P\}$, whose zeros are the only
    ones that enter the zero sum;
  \item the \emph{buffer} $R_B=\{0\le\lambda\le C_B,\ |\mu|\le C_B\}$, with $C_B>C_P+M$;
  \item the \emph{height-one region} $|\gamma|\le1$, in which the far-density lemma operates.
\end{itemize}
The constants $C_P$, $M$ and $C_B$ are fixed independently of $q$, in the order described in
\cref{sec:join}. For sufficiently large $q$, $R_P\subset R_B\subset R(l)$, and the buffer
has physical height $C_B/\log q<1$. The word \emph{inside} means that the selected first
zero belongs to $R_B$; \emph{outside} means that it does not. It does not refer to $R(l)$.

\subsection{Families and the first zeros}
\begin{definition}[Families and successive minima]\label{def:families}
A \emph{family} is the set $\mathcal F(\chi)=\{\chi,\bchi\}$ for a nonprincipal character
$\chi$. It has one member when $\chi$ is real and two otherwise.
If there is a nonprincipal zero in $R(l)$, select a pair $(\chi_1,\rho_1)$ with least
parameter and write $\lambda_1=\lambda_{\rho_1}$, $\mu_1=\mu_{\rho_1}$, and
$\mathcal F_1=\mathcal F(\chi_1)$. Remove all occurrences belonging to $\mathcal F_1$ and
repeat to define $(\chi_2,\rho_2)$, $\lambda_2$, and $\mathcal F_2$; continue in this way.
Thus $\lambda_1\le\lambda_2\le\lambda_3\le\cdots$. We call $\mathcal F_1$ the \emph{first family},
$\mathcal F_2$ the \emph{second family} and $\mathcal F_3$ the \emph{third family}, and $(\chi_1,\rho_1)$
the \emph{first zero}.
If a subsequent minimum is over an empty set, its value is $+\infty$.
Ties may be resolved by any admissible choice; each result used below holds for all such
choices. If the initial set is empty, we use the empty-region argument of
\cref{sec:exteriors}.
\end{definition}

These are the selections of Xylouris \cite[\S3.1.2]{Xylouris2011nullstellen}.
They order \emph{families}, not the individual zeros of one $L$-function.

\begin{definition}[The additional first-family zero]\label{def:additionalzero}
The \emph{distinguished occurrences} are one copy of $\rho_1$ as a zero of $\chi_1$, together with one
copy of $\overline{\rho_1}$ as a zero of $\chi_1$ if $\chi_1$ is real and $\rho_1$ is not, or as a zero of
$\bchi_1$ if $\chi_1$ is not real. We write $\lambda'$ for the least parameter of a zero occurrence of
$\chi_1$ or $\bchi_1$ in $R(l)$ that is not distinguished; thus $\lambda'=\lambda_1$ when $\rho_1$ is a
multiple zero. If no occurrence remains, put $\lambda'=+\infty$.
\end{definition}

When $\lambda'$ is finite, it can always be realized by a zero of $\chi_1$ itself.
Indeed, zeros of $L(s,\bchi_1)$ are conjugates of zeros of $L(s,\chi_1)$, with the same
multiplicities. Replace any minimizing occurrence of $\bchi_1$ by its conjugate
occurrence of $\chi_1$. We call every resulting non-distinguished occurrence $\rho'$
with parameter $\lambda'$ \emph{admissible} (a second copy of $\rho_1$ is admissible when
$\rho_1$ is a multiple zero), and
write $\gamma'$ for its height; as $\rho'\in R(l)$, $|\gamma'|\le l$. The arguments below hold for every admissible
choice; in the second-zero split of \cref{subsec:subdivision} the choice is the one used there.

The following table fixes the three types and the two counts used in the first-family
bounds. Here $n=|\mathcal F_1|$, while $\alpha$ counts distinguished occurrences
\emph{per character}.
\begin{center}
\begin{tabular}{@{}llll@{}}
\toprule
Type & First character and zero & $n$ & $\alpha$\\
\midrule
$\mathrm{rr}$ & $\chi_1$ real, $\rho_1$ real & $1$ & $1$\\
$\mathrm{rc}$ & $\chi_1$ real, $\rho_1$ nonreal & $1$ & $2$\\
$\mathrm{complex}$ & $\chi_1$ nonreal & $2$ & $1$\\
\bottomrule
\end{tabular}
\end{center}
For type $\mathrm{rc}$ we choose the conjugate with $\mu_1>0$.

\subsection{The vocabulary of the case analysis}\label{subsec:vocabulary}
In the middle range $0.1\le\lambda_1<1.5$ the proof is a finite case analysis, defined
precisely in \cref{sec:lp,sec:casetree}. Its terms are needed in the analytic sections that
precede those, and we introduce them here.
\begin{itemize}
  \item \emph{Anchor.} The local explicit formulas of \cref{sec:inputs} are applied at points
    $s=1-\sigma/\LL+i\gamma$; the real number $\sigma$ is the \emph{anchor}. An inequality anchored
    at $\sigma$ is useful only if the zeros near $s$ with parameter below $\sigma$ are absent or are
    retained explicitly, so anchors are chosen at or below known lower bounds for zero parameters.
  \item \emph{Specification.} A \emph{specification} $\mathfrak s$ (\cref{def:specification})
    records the type of the first family; an interval $[a,b]$ containing $\lambda_1$, the
    \emph{first-zero cell}; a lower bound $p\le\lambda'$ and possibly an interval
    $[\lo',\hi']$ containing $\lambda'$, called a \emph{gap}; a lower bound $l_2\le\lambda_2$; a
    lower bound $r$ for the parameters of the zeros of height at most $1$ of the families other
    than the first; possibly a reservation (next item); and whether $\rho_1$ lies inside or
    outside the buffer. Its \emph{configuration set} $C(\mathfrak s)$ consists of the zero
    configurations compatible with these data.
  \item \emph{Reserved and ordinary characters.} A reservation concerns a family other than the
    first whose least parameter among zeros of height at most $1$ is as small as possible. It
    fixes the number $n_2\in\{1,2\}$ of its characters and an interval $[\lo_2,\hi_2]$ for that
    parameter. This \emph{reserved family} need not be $\mathcal F_2$. The characters outside the
    first and the reserved family are \emph{ordinary}.
  \item \emph{Representatives.} The \emph{height-one representative} of a nonprincipal character
    $\chi$ is a zero of $L(s,\chi)$ in $R(l)$ with $|\gamma|\le1$ and least parameter; its
    \emph{$T^*$-representative} is defined in the same way with $|\gamma|\le T^*$, for a height
    $T^*\in[\frac12,1]$ fixed in \cref{lem:Tstar} (\cref{def:representative}). The leaf programs
    place a character by the parameter of one of its representatives.
  \item \emph{The case tree.} The \emph{roots} are $4453$ specifications whose configuration sets
    cover the middle range. Each root is refined by splitting intervals and by applying
    zero-location results. This gives a finite tree of specifications, the root's
    \emph{refinement tree}; its vertices are \emph{nodes}, and together these trees form the
    \emph{case tree}. A node at which a zero-location result sharpens one of the bounds is a
    \emph{location update}; a node whose configuration set is shown to be empty is
    \emph{excluded}. The terminal nodes that are not excluded are the \emph{leaves}. Some leaves
    are divided further into finitely many \emph{subcases}.
  \item \emph{Leaf programs, bins and columns.} Each leaf or subcase has a \emph{leaf program}
    (\cref{def:leafprogram}) whose value bounds $W$ for every configuration in it. The parameters of
    the ordinary characters (and in some leaves those of the reserved characters) are divided into
    finitely many intervals, the \emph{bins}. The \emph{column} of a bin is the variable of the
    program that counts the characters whose representative zero (\cref{def:representative}) has
    parameter in the bin. The ordinary characters beyond the last bin form the \emph{tail}, whose
    column is a weighted mass rather than a count. In an outside leaf the first family is
    represented by \emph{hidden columns}, which bin the least parameter $t_\chi$ of the zeros in $R_P$
    of each first-family character (\cref{subsec:firstoutside}). The constant $\mathrm{first}$
    bounds the cost of the first family in an inside leaf ($0$ in an outside leaf), plus any fixed
    charge for the reserved family, and $\mathrm{final}$ is an allowance for the error terms.
  \item \emph{Far budget.} The far-density lemma (\cref{inp:far}) bounds a sum over all characters
    of a positive decreasing function $w$ of the parameter of a selected zero of height at most
    $1$. We call $w(\lambda)$ the \emph{far weight} of such a zero, and the bound
    \eqref{eq:farbudget} the \emph{far budget}; in a leaf program it is one linear constraint.
  \item \emph{Near rows.} The other constraints of a leaf program are conditions on counts and the
    \emph{near rows}. A near row is, with one exception (\cref{subsec:twotest}), an instance of the
    graded lemma (\cref{thm:graded}) in the form of \cref{prop:threshold}. Its terms are
    \emph{entries}, one for each character (occasionally each zero) that it includes. Each entry
    has a \emph{feature}, a lower bound for how strongly its zero is detected, and a
    \emph{diagonal}, what it costs on its own, and one \emph{correlation term} bounds the
    interaction between distinct entries (\cref{prop:threshold}). The entries of a family whose
    zeros are known are grouped into \emph{family terms}. A row is built from two test functions,
    a \emph{detector} and a \emph{Gram test} (\cref{def:neardatum}); by the choice of its anchors and
    sieve weight it is a \emph{family}, \emph{shifted}, \emph{graded} or \emph{two-test} row
    (\cref{tab:rowkinds}). It is linear in the columns once an auxiliary number, its
    \emph{threshold} $\tau$, is fixed.
  \item \emph{Certificates.} A \emph{certificate} of a leaf or subcase is a binary tree, the
    \emph{certificate tree}, whose nodes are \emph{boxes} of thresholds, starting from a box that
    contains every admissible threshold; its leaves are the \emph{terminal boxes}. On a terminal
    box each near row is replaced by a linear relaxation, and each choice among the alternatives
    of these relaxations, a \emph{relaxation case}, is either shown to be infeasible (an
    \emph{exclusion}) or bounded below $1$ by integer dual multipliers, in exact integer
    arithmetic (\cref{sec:lp}).
\end{itemize}

\subsection{Elementary notation}
For real $x$ we write $x_+=\max(x,0)$, and $\1_E$ denotes the indicator of a condition or
set $E$. The symbol $\varphi(q)$ is Euler's totient; the separately defined
$\varphi(\chi)$ is a conductor coefficient (\cref{sec:inputs}). A \emph{cost} means an
upper bound for a contribution to the normalized zero sum $W$; a \emph{budget} is the
right-hand side of a constraint on these contributions. All finite decimals used as
parameters are exact rational numbers unless an approximation is explicitly indicated.

\subsection{Asymptotic conventions}
Statements about zeros below are understood for all sufficiently large $q$, unless a
different range is stated. The threshold $q_0$ depends on finitely many fixed choices
(the exponent, test functions, weights, constants and certificates, and the tolerances of the
published results we use); the order in which these are chosen is spelled out in \cref{sec:join}.
We write $o(1)$ for a quantity tending to $0$ as $q\to\infty$, uniformly in everything except
the fixed choices. The number
\[
  \eta=10^{-6}
\]
is a fixed tolerance used throughout, and the exponent is
\[
  L=3.99 .
\]
Certificate coefficients use the scale $S=10^{16}$. Counts of characters are unscaled;
the precise conversion between coefficients and real quantities is given in
\cref{subsec:leafdata}.

\subsection{Constants and the error budget}\label{subsec:budget}
The fixed constants of the proof, with their values, their roles and the stage of
\cref{subsec:order} at which each is chosen, are listed in \cref{tab:constants} in
Appendix~\ref{app:constants}. The tolerance $\eta=10^{-6}$ is spent as follows; the allowance $5\eta$ in each leaf value covers
the analytic errors with room to spare.
\begin{enumerate}[(1)]
  \item \emph{Detection.} With $\eps=\eta H_0$ in \cref{inp:criterion}, the weighted prime sum is at
    least $\frac{\LL H_0}{\varphi(q)}(1-W-\eta)$; so $W\le1-2\eta$ suffices (\cref{prop:detect}).
  \item \emph{Analytic error.} In \eqref{eq:Wsplit} the error $E$ has two parts
    (\cref{subsec:allowance}). The envelope errors, which include the smoothing error by
    \eqref{eq:thetachoice}, total less than $19\eps_1$ for the ordinary characters, since
    $\sum_\chi e^{-A\lambda_\chi}\le K_{\mathrm{far}}<19$ by \eqref{eq:Kfar}, and at most $4\eps_1$ for
    the at most two first-family and two reserved characters; together less than $23\eps_1=\eta$.
    The outside first-family bound (\cref{prop:firstoutside}) adds at most
    $4c_{\mathrm{out}}/(H_0C_B^2)\le\eta$, by the choice of $C_B$ in (S4). Hence $E\le2\eta$.
  \item \emph{Allowance.} Every leaf value includes $\mathrm{final}/S=5\eta$
    (\cref{subsec:leafmodel}); so does the program of \cref{thm:large}.
  \item \emph{Conclusion.} By the proof of \cref{prop:realization},
    $W\le\mathcal V(x)-\mathrm{final}/S+E\le\mathcal V(x)-3\eta$, that is, $W+3\eta\le\mathcal V(x)<1$
    by \cref{thm:certificate}. Hence $W<1-3\eta$, and \cref{prop:detect} applies with a margin $\eta$.
  \item \emph{Other tolerances.} They enter the constraints rather than $E$: $\eps=\eta J^2$ in
    \cref{inp:far} gives the far budget $(1+\eta)V$ of \eqref{eq:farbudget}, the tolerance $\eta'$ is
    built into every feature, diagonal and correlation term of a near row (\cref{cor:zeroform}), and every
    stored number is rounded in the safe direction (\cref{subsec:leafmodel}).
\end{enumerate}

% !TEX root = ../linnik399.tex
\section{Published inputs}\label{sec:inputs}

This section records the analytic estimates used in the proof. Statements labelled
\emph{Input} are taken from the literature, with notation adapted to this paper. Our main
sources are Heath-Brown \cite{HeathBrown1992zero} and Xylouris's dissertation
\cite{Xylouris2011nullstellen}; the dissertation, rather than Xylouris's shorter article,
is the source of the lemma and equation numbers cited below. Any additional hypothesis
needed from a source proof, or extension of a printed statement, is identified explicitly.

\subsection{Test functions}

\begin{definition}[Conditions 1 and 2]\label{def:conditions}
Let $x_0>0$ and let $f:[0,\infty)\to\RR$ be continuous with $f(t)=0$ for $t\ge x_0$.
\begin{itemize}
  \item $f$ satisfies \emph{Condition~1} if it is twice continuously differentiable on
    $(0,x_0)$ with $|f''|\le B$ there for some constant $B$
    \cite[p.~280]{HeathBrown1992zero}, \cite[Bedingung~1, p.~17]{Xylouris2011nullstellen}.
  \item $f$ satisfies \emph{Condition~2} if $f\ge0$ and its Laplace transform
    $F(z)=\int_0^\infty f(t)e^{-zt}\,dt$ satisfies $\Re F(z)\ge0$ for $\Re z\ge0$
    \cite[p.~286]{HeathBrown1992zero}, \cite[Bedingung~2, p.~18]{Xylouris2011nullstellen}.
\end{itemize}
\end{definition}

Condition~1 supplies the regularity needed for the explicit formula. Condition~2 ensures
that zero terms with nonnegative real argument can be discarded in an upper bound.
The two conditions have different roles, so we specify each one when it is needed.

A function satisfying Condition~1 is bounded, its derivative is bounded on $(0,x_0)$, and its
transform $F$ is entire. The \emph{conductor coefficient} of a nonprincipal character $\chi$ modulo
$q$ is \cite[p.~17]{Xylouris2011nullstellen}
\[
  \varphi(\chi)=\begin{cases}\frac14,& q\text{ cube-free, or }\ord\chi\le\LL,\\[2pt]
  \frac13,&\text{otherwise.}\end{cases}
\]
In particular $\varphi(\chi)=\frac14$ for every real character when $q\ge8$ (its order is $2\le\LL$), and
$\varphi(\chi)\le\frac13$ always.

\begin{inp}[Comparison lemma {\cite[Lemma~4.1]{HeathBrown1992zero}, \cite[Lemma~3.5, p.~22]{Xylouris2011nullstellen}}]\label{inp:comparison}
Let $F_1,F_2$ be holomorphic in $\{\Re z\ge0\}$, with $\Re F_1(z)\ge|F_2(z)|$ on $\Re z=0$, and
suppose that $F_1$ and $F_2$ tend to $0$ uniformly as $|z|\to\infty$ in $\Re z\ge0$. Then
$\Re F_1(z)\ge|F_2(z)|$ for all $\Re z\ge0$.
\end{inp}

This is a consequence of the maximum principle
\cite{PhragmenLindelof1908extension,Titchmarsh1939theory}. All transforms below are Laplace
transforms of bounded functions of compact support \cite{Widder1941laplace}.

\subsection{Explicit formulas}
The following two lemmas are Heath-Brown's Lemmas~5.3 and~5.2, as stated by Xylouris
\cite[Lemmas~3.1 and~3.2, p.~18]{Xylouris2011nullstellen}. In them $s=\sigma+it$ with
\begin{equation}\label{eq:lemmaregion}
  |\sigma-1|\le\frac{(\log\LL)^{1/2}}\LL,\qquad|t|\le\LL .
\end{equation}

\begin{inp}[Principal character]\label{inp:X31}
Let $f$ satisfy Condition~1 and let $s$ satisfy \eqref{eq:lemmaregion}. Then
\[
  \sum_{n=1}^\infty\Lambda(n)\frac{\chi_0(n)}{n^s}\,f\Bigl(\frac{\log n}\LL\Bigr)=\LL\,F\bigl((s-1)\LL\bigr)
  +O\Bigl(\frac\LL{\log\LL}\Bigr),
\]
with an implied constant depending only on $f$.
\end{inp}

\begin{inp}[Local explicit formula]\label{inp:X32}
Let $\chi\ne\chi_0$, let $s=\sigma+it$ satisfy \eqref{eq:lemmaregion}, and let $f$ satisfy
Condition~1 with $f(0)\ge0$. For every $\eps>0$ there are
$\delta=\delta(f,\eps)\in(0,1)$ and $q_0=q_0(f,\eps)$, independent of $\chi$ and $s$, such that
for $q\ge q_0$
\[
  \sum_{n=1}^\infty\Lambda(n)\,\Re\frac{\chi(n)}{n^s}\,f\Bigl(\frac{\log n}\LL\Bigr)\le
  -\LL\sum_{|1+it-\rho|\le\delta}\Re F\bigl((s-\rho)\LL\bigr)+\frac{f(0)}2\varphi(\chi)\LL+\eps\LL,
\]
where the sum runs over the nontrivial zeros of $L(s,\chi)$ in the disc, with multiplicity.
\end{inp}

These are smoothed forms of the explicit formula
\cite{Weil1952formules,Davenport2000multiplicative,MontgomeryVaughan2007multiplicative,IwaniecKowalski2004analytic};
the conductor coefficient $\varphi(\chi)$ comes from Burgess's bounds, through Heath-Brown's
growth estimate for $L(s,\chi)$ and a Jensen-type formula \cite[\S\S2--3]{HeathBrown1992zero}. Heath-Brown's printed Lemma~5.2 omits the factor $\Lambda(n)$ on
the left, a misprint that Xylouris's statement corrects. Xylouris notes \cite[p.~36]{Xylouris2011nullstellen}
that the radius $\delta$ and the threshold $q_0$ of \cref{inp:X32}, and the implied constant of
\cref{inp:X31}, depend only on $\eps$ and on upper bounds for $x_0$, $x_0^{-1}$, $\sup|f|$ and
$\sup_{(0,x_0)}(|f'|+|f''|)$; so both inputs hold uniformly for a family of test functions with
such common bounds. In the proof of Heath-Brown's Lemma~3.1, on which Lemma~5.2
rests, the disc radius is $\delta=\min(1/(2k),\eps_0/(3c_0k^2))$ with $k\ge3$, so $\delta\le1/6$.
The same proof works for every smaller radius, and Heath-Brown remarks after Lemma~5.2 that
$\delta=1/\log\LL$ may be taken. We may and do assume $\delta<1/2$.

\begin{inp}[The height $l$ {\cite[Lemma~6.1]{HeathBrown1992zero}, \cite[Lemma~3.3, p.~19]{Xylouris2011nullstellen}}]\label{inp:rectangle}
There are $q_0$ and a number $l=l(q)\in\NN$ with $l\le\LL/10$ such that for $q\ge q_0$ the function
$\prod_{\chi}L(s,\chi)$ has no zeros in $R(10l)\setminus R(l)$.
\end{inp}

The principal $L$-function has no zeros in $R(l)$ for large $q$ \cite[p.~21]{Xylouris2011nullstellen}.
All results we quote from the tables of \cite{HeathBrown1992zero,Xylouris2011nullstellen} concern zeros in
$R(l)$ for this $l$; Heath-Brown's proof of Lemma~6.1 produces $l$ and the later sections use it
only through $1\le l\le\LL/10$ and the zero-free annulus.

\begin{inp}[All-zero envelope {\cite[Lemma~3.10, pp.~35--36]{Xylouris2011nullstellen}}]\label{inp:envelope}
Let $C_0$ be the size of the square in \cref{inp:criterion}, and fix
$\eps,\lambda_{11}>0$ and an integer $m\ge1$. For $\lambda\in(\lambda_{11},C_0]$, suppose that:
\begin{enumerate}[(i)]
  \item each $f^\lambda_{1i}$, $1\le i\le m$, satisfies Condition~1 with support parameter
    $x_{0i}>0$ and $f^\lambda_{1i}(0)\ge0$;
  \item the quantities $x_{0i}$, $x_{0i}^{-1}$, $\sup|f^\lambda_{1i}|$, and
    $\sup_{(0,x_{0i})}(|(f^\lambda_{1i})'|+|(f^\lambda_{1i})''|)$ have common bounds depending
    only on $\lambda_{11}$ and $C_0$, and $f_1^\lambda=\sum_i f^\lambda_{1i}\ge0$;
  \item $H_2$ is holomorphic in $\Re z\ge0$,
    $|H_2(\lambda+it)|\le\Re F_1^\lambda(it)$ for every real $t$, and both transforms tend
    uniformly to $0$ as $|z|\to\infty$ in that half-plane;
  \item $\chi\ne\chi_0$ and $L(s,\chi)$ has no zero with
    $1-\lambda/\LL<\beta\le1$ and $|\gamma|\le1$.
\end{enumerate}
Here $F_1^\lambda$ is the Laplace transform of $f_1^\lambda$.
Then there is an effectively computable $q_0$,
depending on $\eps,m,\lambda_{11},C_0$ but not on $\lambda$, such that for $q\ge q_0$
\[
  {\sum_\rho}'\bigl|H_2\bigl((1-\rho)\LL\bigr)\bigr|\le F_1^\lambda(-\lambda)+\frac{f_1^\lambda(0)}6+\eps,
\]
the sum running over the zeros of $L(s,\chi)$ in the square of \cref{inp:criterion}.
\end{inp}

We have included the hypothesis $f^\lambda_{1i}(0)\ge0$ in (i), although it is omitted
from the printed statement. The proof \cite[pp.~35--36]{Xylouris2011nullstellen} applies \cref{inp:comparison}, then
\cref{inp:X32} at $s=1-\lambda/\LL$ to each $f^\lambda_{1i}$, then \cref{inp:X31}. It therefore uses
$f^\lambda_{1i}(0)\ge0$ for each $i$; every application below has this property.
The zero-free hypothesis is used only to know
that the zeros of $L(s,\chi)$ in the disc of \cref{inp:X32} about $1$ have $\beta\le1-\lambda/\LL$.
We use the lemma in this localized form (\cref{sec:costs}).

\subsection{Character sums and the Selberg sieve}

\begin{inp}[Burgess {\cite[Lemma~2.1]{HeathBrown1992zero}, $k=3$}]\label{inp:burgess}
Let $q\ge1$ and let $\chi$ be a primitive character modulo $q$. Let $N\ge1$ and $1\le H\le q$. Then
for every $\eps>0$
\[
  \sum_{N<n\le N+H}\chi(n)\ll_\eps q^{1/9+\eps}H^{2/3}.
\]
\end{inp}

This is the case $k=3$ of Heath-Brown's statement, which is Burgess's theorem
\cite{Burgess1962primitive,Burgess1962character,Burgess1963character,Burgess1986character}; see
\cite{HeathBrown2013burgess} for an account. It improves on the P\'olya--Vinogradov inequality
\cite{Polya1918verteilung,Vinogradov1918distribution} for short intervals. For special moduli stronger
bounds are known \cite{BarbanLinnikChudakov1964prime,Chang2014short,BanksShparlinski2019bounds,
Kerr2019moments,PetrowYoung2020weyl}, but we do not use them.

The weights $\xi_d$ in the next input are Selberg's sieve weights
\cite{Selberg1947elementary,Selberg1991lectures} in the form used for zero-density estimates by Graham
and Heath-Brown; see \cite{HalberstamRichert1974sieve,Motohashi1983lectures,FriedlanderIwaniec2010opera}
for the Selberg sieve in general.

\begin{inp}[Graham {\cite[p.~84]{Graham1978asymptotic}; \cite[(11.13), $U=1$]{HeathBrown1992zero}}]\label{inp:graham}
Let $z\ge2$ and $\xi_d=\mu(d)\log(z/d)/\log z$ for $d\le z$, $\xi_d=0$ otherwise. Then for $N\ge z$
\[
  \sum_{n\le N}\Bigl(\sum_{d\mid n}\xi_d\Bigr)^2=\frac N{\log z}+O\Bigl(\frac N{\log^2z}\Bigr).
\]
\end{inp}

\subsection{Log-free zero density}

\begin{inp}[Log-free density {\cite[Theorem~1]{Jutila1977linnik}, as quoted in the proof of \cite[Lemma~6.1]{HeathBrown1992zero}}]\label{inp:logfree}
For $T\ge1$, $\frac45\le\sigma\le1$ and every $\eps>0$,
\[
  \sum_{\chi\bmod q}N(\sigma,T,\chi)\ll_\eps(qT)^{(2+\eps)(1-\sigma)},
\]
where $N(\sigma,T,\chi)$ counts the nontrivial zeros $\beta+i\gamma$ of $L(s,\chi)$, with multiplicity, in
$\beta\ge\sigma$, $|\gamma|\le T$.
\end{inp}

Estimates of this kind go back to Linnik \cite{Linnik1944least,Linnik1944deuring}; see
\cite{Turan1961density,Fogels1965zeros,Gallagher1970large,Montgomery1971topics,Jutila1977linnik,
Motohashi1975density,Bombieri1987grand} and \cite[(1.4)]{HeathBrown1992zero},
\cite[Prinzip~3, p.~11]{Xylouris2011nullstellen}; explicit versions are in
\cite{ThornerZaman2024explicit}, and related estimates in
\cite{Ingham1940estimation,Huxley1972difference,Jutila1972density,HeathBrown1979density,
SoundararajanThorner2019weak,Pintz2019some}. We use it only to see that the number of zeros of all
$L(s,\chi)$ modulo $q$, counted with multiplicity, with $\lambda\le3$ and $|\gamma|\le T$ is bounded
independently of $q$, for $T=2$ and for $T=\LL$: take $1-\sigma=3/\LL$; then
$(qT)^{(2+\eps)(1-\sigma)}\ll(q\LL)^{7/\LL}\ll1$.

\subsection{Exceptional zeros}

\begin{inp}[Siegel zeros {\cite[Corollary~1, p.~406]{HeathBrown1990siegel}}]\label{inp:siegel}
Suppose $L(\beta_0,\chi)=0$ for a real character $\chi$ modulo $q$, not necessarily primitive,
and a real $\beta_0\ge1-1/(3\log q)$, and put $\eta_0=((1-\beta_0)\log q)^{-1}$. For every $\delta>0$
there is an effectively computable constant $\eta_{\mathrm{HB}}(\delta)$ such that $\eta_0\ge\eta_{\mathrm{HB}}(\delta)$
implies $P(a,q)\le q^{3+\delta}$ for every $a$ coprime to $q$.
\end{inp}

The zero-location results of Heath-Brown and Xylouris that we use are stated in
\cref{sec:location}, where they are applied, and the inputs for the small first-zero regime in
\cref{sec:exteriors}.

% !TEX root = ../linnik399.tex
\section{Prime detection}\label{sec:detection}

Our goal is to make a nonnegative weighted sum over primes in a reduced residue class
strictly positive. The explicit formula expresses this as a comparison between a main
term and a sum over zeros. We first fix the weight, then state the precise zero-sum
bound that will suffice.

\subsection{The kernel}\label{subsec:kernel}
Fix
\[
  T=0.416829,\qquad \kappa=T/16,
\]
and the sixteen step heights $b_0,\dots,b_{15}$ of \cref{tab:beta}. Put
\[
  \psi(t)=\sum_{i=0}^{15}b_i\,\1_{[i\kappa,(i+1)\kappa)}(t),\qquad
  \Psi(z)=\int_0^T\psi(t)e^{-zt}\,dt .
\]
For the target exponent $L$ put $A=L-2T$; at $L=3.99$ we have $A=3.156342$. The \emph{prime
weight} is
\[
  h(t)=(\psi*\psi)(t-A),
\]
which is supported in $[A,L]$. Its Laplace transform is
\begin{equation}\label{eq:Hkernel}
  H(z)=\int h(t)e^{-zt}\,dt=e^{-Az}\,\Psi(z)^2,\qquad H_0=H(0)=\Psi(0)^2 .
\end{equation}
Numerically $\Psi(0)=\kappa\sum_ib_i=0.193642\ldots$ and $H_0=0.037497\ldots$.

\begin{table}[ht]
\centering
\small
\begin{tabular}{rl@{\qquad}rl@{\qquad}rl@{\qquad}rl}
\toprule
$i$ & $b_i$ & $i$ & $b_i$ & $i$ & $b_i$ & $i$ & $b_i$ \\
\midrule
0 & 0.02318741 & 4 & 0.23237155 & 8 & 0.47877833 & 12 & 0.75172165 \\
1 & 0.07159337 & 5 & 0.29081583 & 9 & 0.54505502 & 13 & 0.82236428 \\
2 & 0.12265085 & 6 & 0.35147264 & 10 & 0.61279917 & 14 & 0.89340300 \\
3 & 0.17627704 & 7 & 0.41418551 & 11 & 0.68177389 & 15 & 0.96451862 \\
\bottomrule
\end{tabular}
\caption{The step heights of the kernel $\psi$. They are exact decimal numbers.}
\label{tab:beta}
\end{table}

Xylouris \cite[(3.51)]{Xylouris2011nullstellen}, following Heath-Brown \cite[\S13]{HeathBrown1992zero},
uses the triangular weights
\[
  h_{L',K}(t)=\begin{cases}
    t-(L'-2K), & L'-2K\le t\le L'-K,\\
    L'-t, & L'-K\le t\le L',\\
    0, & \text{otherwise},
  \end{cases}
\]
for $L'>2K>0$; the transform of $h_{L',K}$ is $e^{-(L'-2K)z}\bigl((1-e^{-Kz})/z\bigr)^2$.

\begin{lemma}[Decomposition into positive triangular weights]\label{lem:triangles}
For the step function $\psi$ and shifted convolution $h$ defined above,
\[
  h=\sum_{k=0}^{30}c_k\,h_{A+(k+2)\kappa,\,\kappa},\qquad
  c_k=\sum_{\substack{0\le i,j\le15\\ i+j=k}}b_ib_j>0 .
\]
Every triangle in this sum has left endpoint $A+k\kappa\ge A$.
\end{lemma}

\begin{proof}
For $i,j\ge0$ the convolution $\1_{[i\kappa,(i+1)\kappa)}*\1_{[j\kappa,(j+1)\kappa)}$ is the tent
function of height $\kappa$ supported on $[(i+j)\kappa,(i+j+2)\kappa]$, which is
$h_{(i+j+2)\kappa,\kappa}$. Expanding $\psi*\psi$ and shifting by $A$ gives the formula. Each
$c_k$ is a sum of products of positive numbers, and the sum over $i+j=k$ is nonempty for
$0\le k\le30$.
\end{proof}

\subsection{The positivity criterion}
The following is Xylouris's form of the explicit formula for a finite combination of triangles
\cite[(3.56)--(3.58) and the lines between (3.56) and (3.57), pp.~34--35]{Xylouris2011nullstellen}; it
refines Heath-Brown's Lemmas~13.1 and~13.2 \cite{HeathBrown1992zero}.

\begin{inp}[Xylouris's criterion]\label{inp:criterion}
Let $h=\sum_i\alpha_ih_{L_i,K_i}$ be a finite real combination of triangles with
$L_i>2K_i+3$ for every $i$, and let $H$ be its Laplace transform. For every $\eps>0$ there are
$C_0=C_0(\eps)$ and $q_0=q_0(\eps)$ such that for every $q\ge q_0$ and every $a$ coprime to $q$,
\[
  \sum_{p\equiv a\,(q)}\frac{\log p}{p}\,h\Bigl(\frac{\log p}{\LL}\Bigr)\ \ge\
  \frac{\LL}{\varphi(q)}\Bigl[H(0)-\sum_{\chi\ne\chi_0}\ {\sum_\rho}'\ \bigl|H\bigl((1-\rho)\LL\bigr)\bigr|
  -\eps\Bigr],
\]
where $\sum'_\rho$ runs over the zeros $\rho=\beta+i\gamma$ of $L(s,\chi)$, with multiplicity, in the
square $1-C_0/\LL\le\beta\le1$, $|\gamma|\le C_0/\LL$.
\end{inp}

Xylouris states the extension to combinations without a separate proof, and it follows from
Heath-Brown's proof of his Lemma~13.2 \cite[\S13]{HeathBrown1992zero}. That proof starts from the explicit
formula \cite[Lemma~13.1]{HeathBrown1992zero}, which is linear in the weight, and bounds the zeros
outside the square by summing over rectangles $1-\frac{m+1}\LL\le\beta\le1-\frac m\LL$,
$\frac n{2\LL}\le|\gamma|\le\frac n\LL$ with $\max(m,n)\ge R$, using a zero-density bound $\ll e^{3m}(1+n/\LL)^{3/2}$
for each rectangle and $|F((1-\rho)\LL)|\ll e^{-(L'-2K)m}n^{-2}$ for each zero. For a combination with positive
coefficients $c_k$ one has $|H|\le\sum_kc_k|H_k|$, so the same estimate holds with $L'-2K$ replaced by the
smallest left endpoint, which exceeds $3$; and the zeros in the square contribute at most
$\sum|H((1-\rho)\LL)|$.

In the normalized coordinates \eqref{eq:normalized} the square is
$\{0\le\lambda\le C_0,\ |\mu|\le C_0\}$. Enlarging it only adds nonnegative terms to the zero
sum, so the conclusion holds for every square $\{\lambda\le C_P,\,|\mu|\le C_P\}$ with $C_P\ge
C_0$. We fix
\begin{equation}\label{eq:CP}
  C_P\ge\max\bigl(C_0(\eta H_0),3\bigr),
\end{equation}
which fixes the prime window $R_P=\{0\le\lambda\le C_P,\ |\mu|\le C_P\}$ of \cref{sec:notation}. Since $C_P$ is fixed and
$\log\log\LL\to\infty$, for large $q$ every zero in $R_P$ lies in the rectangle $R(l)$ of
\cref{sec:notation}.

\begin{definition}[Normalized zero sum]\label{def:zerosum}
For the kernel $H$ in \eqref{eq:Hkernel} and the prime window fixed by \eqref{eq:CP}, put
\[
  W=W(q)=\frac1{H_0}\sum_{\chi\ne\chi_0}\ \sum_{\rho\in R_P}\bigl|H\bigl((1-\rho)\LL\bigr)\bigr| ,
\]
where each zero is counted with its multiplicity.
\end{definition}

\begin{proposition}[A zero-sum bound that guarantees a prime]\label{prop:detect}
Let $L>3+2T$. There is $q_0$ such that for every $q\ge q_0$ with $W(q)\le1-2\eta$ and every $a$ coprime
to $q$ there is a prime $p$ with
\[
  p\equiv a\pmod q,\qquad q^{A}<p<q^{L},\qquad A=L-2T.
\]
Here $W(q)$ is formed from the kernel and prime window for this fixed $L$; the threshold
$q_0$ is independent of $a$ and of the zero configuration of $q$.
\end{proposition}

\begin{proof}
By \cref{lem:triangles}, $h$ is a finite positive combination of triangles $h_{L_i,\kappa}$ with
$L_i-2\kappa=A+k\kappa\ge A>3$. So \cref{inp:criterion} applies with $\eps=\eta H_0$; it
gives a $q_0$ such that for every $q\ge q_0$ with $W(q)\le1-2\eta$,
\[
  \sum_{p\equiv a\,(q)}\frac{\log p}{p}\,h\Bigl(\frac{\log p}{\LL}\Bigr)\ge
  \frac{\LL H_0}{\varphi(q)}\bigl(1-W-\eta\bigr)\ge\frac{\LL H_0\,\eta}{\varphi(q)}>0 .
\]
The sum is therefore nonempty. Since $h$ vanishes outside $(A,L)$, some prime $p\equiv a\pmod q$
satisfies $A<\log p/\LL<L$.
\end{proof}

At $L=3.99$ the condition $L>3+2T=3.833658$ holds. Two features of \cref{inp:criterion} matter
below.
\begin{enumerate}[(i)]
  \item \emph{Occurrence form.} The zero sum charges each zero occurrence separately, through the
    absolute value of its own term. Every bound in \cref{sec:costs} is a bound for a sum of such
    individual terms. In particular we never need the cancellation between zeros that a grouped
    form of the explicit formula would record.
  \item \emph{Size of a term.} Since $|\Psi(\lambda-i\mu)|\le\Psi(\lambda)$ for $\lambda\ge0$,
    \begin{equation}\label{eq:termbound}
      \bigl|H\bigl((1-\rho)\LL\bigr)\bigr|=e^{-A\lambda}\,|\Psi(\lambda-i\mu)|^2\le
      e^{-A\lambda}\,\Psi(\lambda)^2\le H_0e^{-A\lambda}.
    \end{equation}
    A zero at distance $\lambda$ from $\sigma=1$ costs at most $e^{-A\lambda}$ times $H_0$. The
    difficulty is entirely with the number of zeros close to $\sigma=1$.
\end{enumerate}

If no nonprincipal $L(s,\chi)$ has a zero in $R(l)$, then for large $q$ it has none in $R_P$, and
$W=0$. This is the first exterior regime (\cref{sec:exteriors}).

% !TEX root = ../linnik399.tex
\section{Zero costs}\label{sec:costs}

The zero sum $W$ of \cref{sec:detection} is a sum over characters. This section bounds the
contribution of a single character, in terms of an \emph{anchor} (\cref{subsec:vocabulary}): a
number $\lambda_*$ such that no zero of $L(s,\chi)$ near $s=1$ has parameter below $\lambda_*$, at which
the explicit formula is evaluated. For most characters the anchor is
the parameter of a representative zero (\cref{def:representative}); the first family needs a
finer treatment. The argument is Xylouris's proof of his Lemma~3.10 \cite[pp.~35--36]{Xylouris2011nullstellen},
which we reproduce in the localized form we need.

\subsection{The envelope functions}\label{subsec:envelope}
The function $\mathcal G_\phi(\lambda)$ below bounds the total cost of one character.
The auxiliary functions $B_\phi$, $\tilde h$, and $C$ separate the common exponential
decay from the correction for a distinguished zero. The parameter $\phi$ will be
$\varphi(\chi)$, or an upper bound for it.

For real $\lambda$ let $\psi_\lambda(u)=\psi(u)e^{-\lambda u}$ and
\[
  f_\lambda(t)=2\int_0^\infty\psi_\lambda(u)\,\psi_\lambda(u+t)\,du
  =2\int_0^{T-t}\psi(u)\psi(u+t)e^{-\lambda(2u+t)}\,du\qquad(t\ge0),
\]
so that $f_\lambda(t)=0$ for $t\ge T$. Let $F_\lambda$ be its Laplace transform. For
$0\le\phi\le\frac13$, $\lambda\ge0$ and $p\ge a\ge0$ put
\begin{align*}
  B_\phi(\lambda)&=\frac1{H_0}\Bigl[2\int_0^T\psi(u)e^{-2\lambda u}\int_u^T\psi(v)\,dv\,du
    +\phi\int_0^T\psi(u)^2e^{-2\lambda u}\,du\Bigr],\\
  \mathcal G_\phi(\lambda)&=e^{-A\lambda}B_\phi(\lambda),\qquad \mathcal G=\mathcal G_{1/3},\qquad
  \tilde h(\lambda)=\frac{\Psi(\lambda)^2}{H_0},\\
  C(p,a)&=\frac2{H_0}\int_0^T\psi(u)e^{-2pu}\int_u^T\psi(v)e^{-a(v-u)}\,dv\,du .
\end{align*}

\begin{lemma}[Properties of the envelope functions]\label{lem:envelope}
With the functions just defined, the following hold. Parts (i)--(iii) hold for every
real $\lambda$; in (iv)--(v), take $0\le\phi\le1/3$, $\lambda\ge0$, and $p\ge a\ge0$.
\begin{enumerate}[(i)]
  \item $f_\lambda\ge0$, $f_\lambda$ is Lipschitz on $[0,T]$, and $f_\lambda(T)=0$.
  \item $\Re F_\lambda(iy)=|\Psi(\lambda+iy)|^2$ for real $y$.
  \item $\Re F_\lambda(z)\ge|\Psi(\lambda+z)|^2$ for $\Re z\ge0$. In particular $\Re F_\lambda(z)\ge0$
    for $\Re z\ge0$.
  \item $H_0B_\phi(\lambda)=F_\lambda(-\lambda)+\frac\phi2f_\lambda(0)$ and
    $H_0C(p,a)=F_p(a-p)$. Moreover $C(a,a)=\tilde h(a)$.
  \item $B_\phi$, $\mathcal G_\phi$ and $\tilde h$ are positive and decreasing on $[0,\infty)$, and so is
    $\mathcal G_\phi/w$, where $w$ is either far weight of \cref{sec:far}.
\end{enumerate}
\end{lemma}

\begin{proof}
(i) Nonnegativity is clear. For $0\le t<t'\le T$,
$|f_\lambda(t)-f_\lambda(t')|\le2\|\psi_\lambda\|_\infty\int|\psi_\lambda(u+t)-\psi_\lambda(u+t')|\,du$,
which is $O(t'-t)$ because $\psi_\lambda$ has bounded variation.

(ii) Let $a_\lambda(t)=\int\psi_\lambda(u)\psi_\lambda(u+|t|)\,du$ for $t\in\RR$; it is even and
integrable, and its Fourier transform is $|\widehat{\psi_\lambda}(y)|^2=|\Psi(\lambda+iy)|^2$.
Hence $\Re F_\lambda(iy)=\int_0^\infty f_\lambda(t)\cos(yt)\,dt=\int_\RR a_\lambda(t)e^{-iyt}\,dt
=|\Psi(\lambda+iy)|^2$.

(iii) Apply \cref{inp:comparison} to $F_1=F_\lambda$ and $F_2(z)=\Psi(\lambda+z)^2$. Both are
entire. By (i), $F_\lambda(z)=z^{-1}\bigl(f_\lambda(0)+\int_0^Tf_\lambda'(t)e^{-zt}\,dt\bigr)$, so
$|F_\lambda(z)|\le(f_\lambda(0)+\|f'_\lambda\|_1)/|z|$ for $\Re z\ge0$; similarly
$|\Psi(\lambda+z)|\le(\psi_\lambda(0^+)+V(\psi_\lambda))/|z|$ with $V$ the total variation. So both tend to
$0$ uniformly, and on $\Re z=0$ we have equality by (ii).

(iv) Substitute $v=u+t$:
$F_\lambda(-\lambda)=2\int\!\!\int\psi(u)\psi(v)e^{-\lambda(u+v)}e^{\lambda(v-u)}\,dv\,du$ over $0\le u\le v\le T$,
which is the first term of $H_0B_\phi$; and $f_\lambda(0)=2\int\psi^2e^{-2\lambda u}$. The same
substitution gives $F_p(a-p)=2\int_0^T\psi(u)e^{-2pu}\int_u^T\psi(v)e^{-a(v-u)}\,dv\,du$. Finally
$H_0C(a,a)=2\int\!\!\int_{u\le v}\psi(u)\psi(v)e^{-a(u+v)}=\Psi(a)^2$.

(v) All integrands are positive and decrease in $\lambda$. For the last claim,
\[
  \frac{\mathcal G_\phi(\lambda)}{w(\lambda)}=B_\phi(\lambda)\int_{u_0}^xw_0(t)^2e^{(2t-A)\lambda}\,dt,
\]
and $2t-A<0$ on $[u_0,x]$ (\cref{lem:decay}).
\end{proof}

The number $\mathcal G(\lambda)$ is the cost, in units of $H_0$, that the argument below assigns to a
character whose zeros near $s=1$ all have parameter at least $\lambda$. For a real character
the factor $\phi=\frac13$ may be replaced by $\frac14$ (\cref{prop:envelope}).

\subsection{Smoothing}\label{subsec:smoothing}
Since $\psi$ is a step function, $f_\lambda$ has corners at the multiples of $\kappa$ and does
not satisfy Condition~1. We therefore work with smooth approximations of $\psi$ in the
auxiliary functions only; the prime weight $h$ and its transform $H$ are never changed. For
$\theta\in(0,1)$ fix a function $\psi^\theta\in C^\infty(\RR)$ with support in $(0,T)$ such that
\[
  0\le\psi^\theta\le1,\qquad\|\psi-\psi^\theta\|_{L^1}\le\theta,
\]
for instance a mollification of $\psi\,\1_{[\theta/4,\,T-\theta/4]}$. Let $\Psi^\theta$,
$f^\theta_\lambda$, $F^\theta_\lambda$, $B^\theta_\phi$, $\tilde h^\theta$ and $C^\theta$ be defined as above
with $\psi^\theta$ in place of $\psi$, and $H_0$ kept.

\begin{lemma}[Uniform smooth approximation]\label{lem:smoothing}
Fix $\theta\in(0,1)$ and $\Lambda\ge1$, and choose $\psi^\theta$ as above.
\begin{enumerate}[(i)]
  \item For every $\lambda\in[0,\Lambda]$, $f^\theta_\lambda$ satisfies Condition~1 and
    Condition~2 with $x_0=T$. The quantities $x_0$, $x_0^{-1}$, $\sup|f^\theta_\lambda|$ and
    $\sup_{(0,T)}(|(f^\theta_\lambda)'|+|(f^\theta_\lambda)''|)$ are bounded independently of
    $\lambda\in[0,\Lambda]$; the bounds may depend on $\theta$ and $\Lambda$.
    \Cref{lem:envelope}(i)--(iv) hold for $\psi^\theta$.
  \item For $\Re z\ge0$, $\bigl|\Psi(z)^2-\Psi^\theta(z)^2\bigr|\le3\theta$.
  \item As $\theta\to0$, $B^\theta_\phi\to B_\phi$, $\tilde h^\theta\to\tilde h$ and $C^\theta\to C$ uniformly on
    $\{0\le a\le p\le\Lambda\}$ and $0\le\phi\le\frac13$.
\end{enumerate}
\end{lemma}

\begin{proof}
(i) The function $(t,\lambda)\mapsto f^\theta_\lambda(t)$ is smooth on $\RR^2$, and it vanishes for
$t\ge T$ because $\psi^\theta$ is supported in a compact subset of $(0,T)$; this gives
Condition~1 and the uniform bounds. The proof of \cref{lem:envelope} applies verbatim to
$\psi^\theta$, and Condition~2 is \cref{lem:envelope}(i) and (iii) for $\psi^\theta$.
(ii) For $\Re z\ge0$ we have $|\Psi(z)|\le\|\psi\|_1\le T<\frac12$, $|\Psi^\theta(z)|\le\|\psi\|_1+\theta$
and $|\Psi(z)-\Psi^\theta(z)|\le\theta$; hence $|\Psi^2-(\Psi^\theta)^2|\le\theta(2T+\theta)\le3\theta$.
(iii) Each quantity is a bilinear expression in $(\psi,\psi)$ with a kernel bounded by $1$ on
$[0,T]^2$ for the parameters in question, and $\|\psi\|_\infty,\|\psi^\theta\|_\infty\le1$.
\end{proof}

\subsection{The disc and a uniform zero count}\label{subsec:zerocount}
Throughout this section $\lambda_{11}=0.05$, the smoothing parameter $\theta$ is the one fixed in stage
(S3) of \cref{subsec:order}, and $\eps_1=\eta/23$. Let $\mathcal A\subset[\lambda_{11},C_P]$ be the finite set
consisting of the points of the grid $\lambda_{11}+h_{\mathcal A}\ZZ$ in $[\lambda_{11},C_P]$ together with the
anchors of the first-family bounds of all leaves (\cref{subsec:firstinside,subsec:firstoutside}). For a leaf
with first-zero cell $[a,b]$ and lower bound $p^*$ for $\lambda'$ (the lower end of its gap if it has
one, and $p$ otherwise; \cref{subsec:vocabulary}), these are $a$ and $p^*$ if the leaf is inside and
$p^\circ=\max(a,p^*)$ if it is outside (\cref{subsec:leafmodel});
the mesh $h_{\mathcal A}>0$ is chosen so that $|B_\phi(x)-B_\phi(y)|\le\eps_1/4$ for $\phi\in\{\frac14,\frac13\}$
and $x,y\in[0,C_P]$ with $|x-y|\le h_{\mathcal A}$. In stage (S5) we fix a radius $\delta_0\in(0,\frac13]$ such
that \cref{inp:X32} holds with radius $\delta_0$ for each of the finitely many pairs of a test
function and a tolerance used in this paper: the function $f^{(1)}$ below with tolerance $1$, the
functions $f^\theta_\sigma$ ($\sigma\in\mathcal A$) with tolerance $H_0\eps_1/8$, and the detectors (\cref{def:neardatum}) of the near
rows with the tolerances of \cref{lem:response}. This is possible because \cref{inp:X32} remains
true for every smaller radius (see the remark after it). The \emph{disc} of a character $\chi$ is the
set of zeros $\rho$ of $L(s,\chi)$ with $|1-\rho|\le\delta_0$. For fixed $C$ and large $q$ the square
$\{\lambda\le C,\,|\mu|\le C\}$ lies in the disc, since there $|1-\rho|\le\sqrt2C/\LL$.

\begin{lemma}[A uniform bound for zeros in a fixed window]\label{lem:zerocount}
Let $C\ge1$. There are $N=N(C)$ and $q_0$ such that for $q\ge q_0$ the following holds. If
$\chi\ne\chi_0$ and every zero in the disc of $\chi$ has parameter at least $\lambda_{11}$, then $L(s,\chi)$
has at most $N$ zeros, counted with multiplicity, with $\lambda\le C$ and $|\mu|\le C$.
\end{lemma}

\begin{proof}
Let $K_1=1/(4C)$ and $\vartheta(z)=(1-e^{-K_1z})/z$, the transform of $\1_{[0,K_1]}$. Let
$f^{(1)}(t)=2\int_0^{K_1-t}e^{-\lambda_{11}(2u+t)}\,du$ for $0\le t\le K_1$ and $f^{(1)}(t)=0$ for
$t\ge K_1$. This is the function $f_{\lambda_{11}}$ of \cref{subsec:envelope} for the step
$\1_{[0,K_1]}$; it equals $(e^{-\lambda_{11}t}-e^{-\lambda_{11}(2K_1-t)})/\lambda_{11}$ on $[0,K_1]$, so
it satisfies Condition~1, and $f^{(1)}(0)\ge0$. As in \cref{lem:envelope}(iii),
$\Re F^{(1)}(z)\ge|\vartheta(\lambda_{11}+z)|^2\ge0$ for $\Re z\ge0$.

Apply \cref{inp:X32} to $f^{(1)}$ at $s=1-\lambda_{11}/\LL$, with tolerance $1$ and a radius
at most $\delta_0$. For large $q$, the fixed normalized square lies in this smaller disc.
All its zeros still have parameter at least $\lambda_{11}$. In the remainder of this proof,
``the disc'' refers to this smaller disc. The
point $s$ lies in the range \eqref{eq:lemmaregion} for large $q$, and $(s-\rho)\LL=(\lambda_\rho-\lambda_{11})-i\mu_\rho$.
Since $-\Re\chi(n)\le\chi_0(n)$ for every $n$ and $f^{(1)}\ge0$, \cref{inp:X31} gives
\[
  -\sum_n\Lambda(n)\Re\frac{\chi(n)}{n^s}f^{(1)}\Bigl(\frac{\log n}\LL\Bigr)\le
  \sum_n\Lambda(n)\frac{\chi_0(n)}{n^s}f^{(1)}\Bigl(\frac{\log n}\LL\Bigr)=\LL F^{(1)}(-\lambda_{11})+o(\LL).
\]
Hence
$\sum_{\rho\text{ in the disc}}\Re F^{(1)}\bigl((\lambda_\rho-\lambda_{11})-i\mu_\rho\bigr)\le
F^{(1)}(-\lambda_{11})+\frac16f^{(1)}(0)+2=:c_1$ for large $q$. Every term is nonnegative, because
$\lambda_\rho\ge\lambda_{11}$. For a zero with $\lambda\le C$, $|\mu|\le C$ put $z=\lambda-i\mu$; then
$|K_1z|\le\sqrt2/4<\frac12$ and $|1-e^{-w}|\ge|w|-|w|^2e^{|w|}/2\ge|w|/2$ for $|w|\le\frac12$, so
the term is at least $|\vartheta(z)|^2\ge K_1^2/4$. The number of such zeros is therefore at most
$N=\lfloor4c_1/K_1^2\rfloor$.
\end{proof}

\subsection{The anchored envelope}
The following inequality is the common source of all the costs below. For $\sigma\in\mathcal A$, a character
$\chi\ne\chi_0$ and a zero $\rho$ put
\[
  \Phi^\theta_\sigma(\rho)=\Re F^\theta_\sigma\bigl((\lambda_\rho-\sigma)-i\mu_\rho\bigr).
\]
By \cref{lem:envelope}(iii), $\Phi^\theta_\sigma(\rho)\ge|\Psi^\theta(\lambda_\rho-i\mu_\rho)|^2\ge0$ whenever
$\lambda_\rho\ge\sigma$.

\begin{lemma}[Anchored inequality]\label{lem:anchored}
Fix the smoothing, finite anchor set $\mathcal A$, and common disc radius as in
\cref{subsec:zerocount}. There is $q_0$, uniform over $\chi\ne\chi_0$ and
$\sigma\in\mathcal A$, such that for $q\ge q_0$,
\[
  \sum_{\rho\text{ in the disc of }\chi}\Phi^\theta_\sigma(\rho)\le H_0B^\theta_{\varphi(\chi)}(\sigma)+\tfrac14H_0\eps_1 .
\]
\end{lemma}

\begin{proof}
Apply \cref{inp:X32} to $f=f^\theta_\sigma$ at $s=1-\sigma/\LL$ (so $t=0$), with tolerance $H_0\eps_1/8$ and
radius $\delta_0$, and bound the prime sum as in the proof of \cref{lem:zerocount}, using
\cref{inp:X31}:
\[
  \LL\sum_{\rho}\Phi^\theta_\sigma(\rho)\le\frac{f^\theta_\sigma(0)}2\varphi(\chi)\LL+\LL F^\theta_\sigma(-\sigma)
  +\tfrac18H_0\eps_1\LL+O\Bigl(\frac\LL{\log\LL}\Bigr).
\]
By \cref{lem:envelope}(iv) for $\psi^\theta$ the right side is at most
$\LL(H_0B^\theta_{\varphi(\chi)}(\sigma)+\frac14H_0\eps_1)$ for large $q$. Only the finitely many functions
$f^\theta_\sigma$, $\sigma\in\mathcal A$, occur, so one $q_0$ serves them all.
\end{proof}

\begin{proposition}[Localized envelope]\label{prop:envelope}
For $q\ge q_0$ the following holds. Let $\chi\ne\chi_0$ and $\lambda_*\in[\lambda_{11},C_P]$, and suppose that
every zero in the disc of $\chi$ has parameter at least $\lambda_*$. Then
\[
  \frac1{H_0}\sum_{\rho\in R_P}\bigl|H\bigl((1-\rho)\LL\bigr)\bigr|\le
  \mathcal G_{\varphi(\chi)}(\lambda_*)+\eps_1e^{-A\lambda_*},
\]
the sum running over the zeros of $L(s,\chi)$ in $R_P$ with multiplicity. In particular the
left side is at most $\mathcal G(\lambda_*)+\eps_1e^{-A\lambda_*}$, and at most
$\mathcal G_{1/4}(\lambda_*)+\eps_1e^{-A\lambda_*}$ if $\chi$ is real.
\end{proposition}

\begin{proof}
Let $N=N(C_P)$ be as in \cref{lem:zerocount}. The parameter $\theta$ of stage (S3) satisfies
\begin{equation}\label{eq:thetachoice}
  3(N+6)\theta+H_0\sup\bigl|B^\theta_\phi-B_\phi\bigr|\le\tfrac14H_0\eps_1,
\end{equation}
the supremum over $[0,C_P]$ and $\phi\in\{\frac14,\frac13\}$ (\cref{lem:smoothing}(iii)). Let
$\sigma\in\mathcal A$ be the largest grid point with $\sigma\le\lambda_*$, so that $\lambda_*-\sigma\le h_{\mathcal A}$. Let
$\rho\in R_P$. Then $\rho$ is in the disc, so $\lambda_\rho\ge\lambda_*\ge\sigma$, and by \eqref{eq:termbound},
\cref{lem:smoothing}(ii) and \cref{lem:envelope}(iii) for $\psi^\theta$,
\[
  \bigl|H\bigl((1-\rho)\LL\bigr)\bigr|=e^{-A\lambda_\rho}|\Psi(\lambda_\rho-i\mu_\rho)|^2
  \le e^{-A\lambda_*}\bigl(\Phi^\theta_\sigma(\rho)+3\theta\bigr).
\]
The zeros of the disc that are not in $R_P$ have $\Phi^\theta_\sigma\ge0$. Summing over the at most $N$
zeros in $R_P$ and using \cref{lem:anchored} and \eqref{eq:thetachoice},
\[
  \sum_{\rho\in R_P}\bigl|H\bigl((1-\rho)\LL\bigr)\bigr|\le
  e^{-A\lambda_*}\Bigl(H_0B^\theta_{\varphi(\chi)}(\sigma)+\tfrac14H_0\eps_1+3N\theta\Bigr)
  \le e^{-A\lambda_*}H_0\bigl(B_{\varphi(\chi)}(\sigma)+\tfrac12\eps_1\bigr),
\]
and $B_{\varphi(\chi)}(\sigma)\le B_{\varphi(\chi)}(\lambda_*)+\frac14\eps_1$ by the choice of $h_{\mathcal A}$. Finally
$\varphi(\chi)\le\frac13$ always, $\varphi(\chi)=\frac14$ for real $\chi$ (whose order is $2\le\LL$), and $B_\phi$
increases with $\phi$.
\end{proof}

\begin{remark}\label{rem:lemma310}
\Cref{prop:envelope} is a localized form of Xylouris's Lemma~3.10 \cite[pp.~35--36]{Xylouris2011nullstellen},
applied with $m=1$, $f_1=f^\theta_\sigma$ and $H_2=(\Psi^\theta)^2$, for which his condition (3.59) holds
with equality by \cref{lem:envelope}(ii); the prime weight $H$ is then compared with $(\Psi^\theta)^2$
through \cref{lem:smoothing}(ii). We require only that the zeros in the disc have parameter at least
$\lambda_*$, where Xylouris requires that $L(s,\chi)$ have no zero with $1-\lambda_*/\LL<\beta\le1$ and
$|\gamma|\le1$; his proof uses this hypothesis only for the zeros in the disc, which have
$|\gamma|\le\delta_0<1$. His printed statement allows $f_1$ to be a sum of $m$ pieces; for $m>1$ his proof
needs in addition that $f_{1i}(0)\ge0$ for each piece, and that the sum over the square be enlarged to
a sum over a disc common to all pieces before \cref{inp:X32} is applied to them. With $m=1$ neither
point arises. Finally, we use finitely many anchors, so that no uniformity in the anchor is needed.
The coefficient $\varphi(\chi)/2$ comes from \cref{inp:X32}; the printed Lemma~3.10 uses
$\varphi(\chi)\le\frac13$.
\end{remark}

\subsection{Ordinary and reserved characters}\label{subsec:ordinarycost}
Let $\chi$ be a nonprincipal character (in the case tree, one outside the first family; in
\cref{thm:large}, any), and let $\lambda_\chi$ be the parameter of its $T^*$-representative or of its
height-one representative (\cref{def:representative}). If $L(s,\chi)$ has a zero in $R_P$, then that zero lies in $R(l)$ and
has $|\gamma|\le C_P/\LL<T^*$, so the representative exists and $\lambda_\chi\le C_P$. In the regimes
in which we use these costs $\lambda_1\ge0.1$, so $\lambda_\chi\ge\lambda_{11}$.

\begin{corollary}[Cost of a character with a representative]\label{cor:ordinarycost}
Suppose $\lambda_1\ge0.1$. Let $\chi\ne\chi_0$ have a zero in $R_P$, and let
$\lambda_\chi$ be the parameter of either its $T^*$-representative or its height-one
representative. For sufficiently large $q$, its contribution to $W$ is at most
\[
  \mathcal G(\lambda_\chi)+\eps_1e^{-A\lambda_\chi},\qquad \eps_1=\eta/23.
\]
If $\chi$ is real, $\mathcal G$ may be replaced by $\mathcal G_{1/4}$.
The threshold is uniform over these choices of character and representative.
\end{corollary}

\begin{proof}
We check the hypothesis of \cref{prop:envelope} with $\lambda_*=\lambda_\chi$. A zero in the disc
with parameter at most $C_P$ lies in $R(l)$ for large $q$, since $C_P<\frac13\log\log\LL$ and
$|\gamma|\le\delta_0<1\le l$; it has $|\gamma|\le\delta_0<\frac12\le T^*$, so its parameter is at least
$\lambda_\chi$ by the minimality of the representative. A zero in the disc with parameter
above $C_P$ has parameter above $\lambda_\chi$.
\end{proof}

\subsection{The first family inside the buffer}\label{subsec:firstinside}
Suppose that $\rho_1\in R_B$, that $\lambda_1\in[a,b]$ with $a\ge\lambda_{11}$, and that every zero
occurrence of the first family in $R(l)$ other than the \emph{distinguished} ones has parameter at
least $p$, where $p\le C_P$; equivalently, $p\le\lambda'$. We assume that $a,p\in\mathcal A$ (in the leaves they are
the numbers $a$ and $p^*$ of \cref{subsec:zerocount}). The distinguished occurrences and their
number $\alpha$ per character are those of \cref{def:additionalzero} and the table after it. Let $\phi_1=\frac14$ if $\chi_1$ is real and $\phi_1=\frac13$ otherwise, and put
\begin{align*}
  J_{\mathrm{old}}&=n\Bigl[e^{-Ap}\bigl(B_{\phi_1}(a)-\alpha\tilde h(a)\bigr)_++\alpha e^{-Aa}\tilde h(a)\Bigr],\\
  J_{\mathrm{new}}&=n\Bigl[\alpha e^{-Aa}\tilde h(a)+e^{-Ap}\bigl(B_{\phi_1}(p)-\alpha C(p,a)\bigr)\Bigr]
  \qquad(p\ge b).
\end{align*}

\begin{proposition}[Cost of the first family inside the buffer]\label{prop:firstinside}
Suppose $\rho_1\in R_B$, $\lambda_1\in[a,b]$, $a,p\in\mathcal A$,
$a\ge\lambda_{11}$, and $p\le\min(\lambda',C_P)$.
For sufficiently large $q$, the contribution of all $n$ characters of the first family
to $W$ is at most $J_{\mathrm{old}}+n\eps_1$. If also $p\ge b$, it is at most
$J_{\mathrm{new}}+n\eps_1$. Here $J_{\mathrm{old}}$ and $J_{\mathrm{new}}$ are the displayed
functions of $a,p,n,\alpha$, and $\phi_1$.
\end{proposition}

The two bounds use different anchors. The bound $J_{\mathrm{old}}$ uses the lower endpoint
$a$ of the first-zero cell. The bound $J_{\mathrm{new}}$ uses $p$, the lower bound for the
remaining zeros, and estimates the distinguished zero together with its correction in
the explicit formula. Their combined expression is smaller than a separate estimate
of the two terms.

\begin{proof}
It suffices to treat one first-family character $\chi$; the bound for its conjugate is the same,
because $R_P$ and the disc are symmetric under conjugation. A zero of the disc of $\chi$ lies in $R(l)$
or has parameter above $\frac13\log\log\LL>C_P$. So every zero of the disc has parameter at least
$\lambda_1\ge a$, and every non-distinguished one has parameter at least $p$. Write $\cS$ for the
distinguished occurrences of $\chi$; they lie in $R_B$, hence in the disc.

\emph{The bound $J_{\mathrm{old}}$.} Every non-distinguished occurrence has parameter at least
$\max(p,\lambda_1)\ge\max(p,a)$, and $J_{\mathrm{old}}$ decreases as $p$ increases; so we may and do
assume $p\ge a$. For a non-distinguished $\rho\in R_P$ we then have $\lambda_\rho\ge p\ge a$, and
$|H((1-\rho)\LL)|\le e^{-Ap}|\Psi(\lambda_\rho-i\mu_\rho)|^2\le e^{-Ap}(\Phi^\theta_a(\rho)+3\theta)$. By
\cref{lem:anchored} at $\sigma=a$, since every term $\Phi^\theta_a$ in the disc is nonnegative and there are
at most $N$ zeros in $R_P$,
\[
  \sum_{\rho\in R_P\setminus\cS}\bigl|H\bigl((1-\rho)\LL\bigr)\bigr|\le
  e^{-Ap}\Bigl(H_0B^\theta_{\phi_1}(a)+\tfrac14H_0\eps_1+3N\theta-\sum_{\rho\in\cS}\Phi^\theta_a(\rho)\Bigr).
\]
For $\rho\in\cS$ put $X_\rho=|\Psi(\lambda_1-i\mu_\rho)|^2$; then $\Phi^\theta_a(\rho)\ge X_\rho-3\theta$ and
$|H((1-\rho)\LL)|=e^{-A\lambda_1}X_\rho$. We add $e^{-A\lambda_1}X_\rho$ for every $\rho\in\cS$; for those in $R_P$
this is their contribution, and for the others it is a nonnegative addition. By \eqref{eq:thetachoice}
the total is at most
\[
  e^{-Ap}H_0B_{\phi_1}(a)+\sum_{\rho\in\cS}X_\rho\bigl(e^{-A\lambda_1}-e^{-Ap}\bigr)+H_0\eps_1 .
\]
If $\lambda_1\le p$, each summand increases with $X_\rho$, and
$X_\rho\le\Psi(\lambda_1)^2\le\Psi(a)^2=H_0\tilde h(a)$; so the total is at most
$H_0[e^{-Ap}(B_{\phi_1}(a)-\alpha\tilde h(a))+\alpha e^{-A\lambda_1}\tilde h(a)+\eps_1]$, which is at most
$H_0(J_{\mathrm{old}}/n+\eps_1)$. If $\lambda_1>p$, each summand is $\le0$, and the total is at
most $H_0[e^{-Ap}B_{\phi_1}(a)+\eps_1]\le H_0[e^{-Ap}(B_{\phi_1}(a)-\alpha\tilde h(a))_++\alpha e^{-Ap}\tilde h(a)+\eps_1]$,
which is again at most $H_0(J_{\mathrm{old}}/n+\eps_1)$ because $p\ge a$.

\emph{The bound $J_{\mathrm{new}}$.} Now $p\ge b\ge\lambda_1$. Apply \cref{lem:anchored} at
$\sigma=p$. The non-distinguished zeros of the disc have $\Phi^\theta_p\ge0$, and for those in $R_P$,
$|H((1-\rho)\LL)|\le e^{-Ap}(\Phi^\theta_p(\rho)+3\theta)$. Therefore the contribution of $\chi$ is at
most
\[
  e^{-Ap}\Bigl(H_0B^\theta_{\phi_1}(p)+\tfrac14H_0\eps_1+3N\theta\Bigr)
  +\sum_{\rho\in\cS}\Bigl(e^{-A\lambda_1}|\Psi(\lambda_1-i\mu_\rho)|^2\,\1_{\rho\in R_P}-e^{-Ap}\Phi^\theta_p(\rho)\Bigr).
\]
The indicator may be replaced by $1$, since the term it multiplies is nonnegative, and
$|\Psi(\lambda_1-i\mu_\rho)|^2\le|\Psi^\theta(\lambda_1-i\mu_\rho)|^2+3\theta$. For a real number $y$ consider
\[
  D(y)=e^{-A\lambda_1}|\Psi^\theta(\lambda_1+iy)|^2-e^{-Ap}\,\Re F^\theta_p(\lambda_1-p+iy)
  =\int_0^\infty k(t)\cos(yt)\,dt,
\]
where, by \cref{lem:envelope}(ii) for $\psi^\theta$ and the definition of $f^\theta_p$,
\[
  k(t)=e^{-A\lambda_1}f^\theta_{\lambda_1}(t)-e^{-Ap}e^{(p-\lambda_1)t}f^\theta_p(t)
  =2\int\psi^\theta(u)\psi^\theta(u+t)e^{-\lambda_1t}\bigl(e^{-\lambda_1(A+2u)}-e^{-p(A+2u)}\bigr)\,du .
\]
Since $\lambda_1\le p$, the integrand is nonnegative, so $k\ge0$ and $D(y)\le D(0)=\int_0^\infty
k(t)\,dt$. Moreover $k(t)$ decreases as $\lambda_1$ increases, for each $t$, because both
$e^{-\lambda_1t}$ and the bracket decrease. Hence
\[
  D(y)\le\int_0^\infty k(t)\big|_{\lambda_1=a}\,dt=e^{-Aa}\Psi^\theta(a)^2-e^{-Ap}F^\theta_p(a-p)
  =H_0\bigl(e^{-Aa}\tilde h^\theta(a)-e^{-Ap}C^\theta(p,a)\bigr),
\]
by \cref{lem:envelope}(iv). By \cref{lem:smoothing}(ii),(iii), $H_0|\tilde h^\theta(a)-\tilde h(a)|\le3\theta$ and
$H_0|C^\theta(p,a)-C(p,a)|\le3\theta$. So each of the $\alpha\le2$ distinguished occurrences contributes at
most $H_0(e^{-Aa}\tilde h(a)-e^{-Ap}C(p,a))+9\theta$, and with \eqref{eq:thetachoice} the contribution of $\chi$ is
at most $H_0(J_{\mathrm{new}}/n+\eps_1)$.
\end{proof}

We write $J_1=\min(J_{\mathrm{old}},J_{\mathrm{new}})$, where $J_{\mathrm{new}}$ is omitted when
$p<b$. In every leaf of the case tree the anchor $p$ is at most $2.293<3\le C_P$, so the hypothesis
$p\le C_P$ holds. In $J_{\mathrm{new}}$ the quantity $B_{\phi_1}(p)-\alpha C(p,a)$ may be negative; the
argument does not require it to be positive.

\subsection{The first family outside the buffer}\label{subsec:firstoutside}
Suppose now that $\rho_1\notin R_B$. Since $\lambda_1<C_B$ in the regimes where this occurs,
this means $|\mu_1|>C_B$. The global zero $\rho_1$ is not in $R_P$ and contributes nothing to $W$,
but the first-family characters may have other zeros in $R_P$. As before, the distinguished
occurrences are $\rho_1$ (and $\overline{\rho_1}$) for $\chi_1$ and $\overline{\rho_1}$ for $\bchi_1$, and
every other zero occurrence of the first family in $R(l)$ has parameter at least $p\ge a$, where $p\in\mathcal A$
(in the leaves, $p=p^\circ$).

For a first-family character $\chi$ with a zero in $R_P$, let $t_\chi$ be the least parameter of
its zeros in $R_P$. Then $t_\chi\ge p$, and the corresponding zero has $|\gamma|\le C_P/\LL\le1$.

\begin{proposition}[Cost of the first family outside the buffer]\label{prop:firstoutside}
Suppose $\rho_1\notin R_B$, $\lambda_1\ge a\ge\lambda_{11}$, and
$p\in\mathcal A$ satisfies $a\le p\le\lambda'$. For a first-family character with a
zero in $R_P$, let $t_\chi$ be the least parameter of its zeros in $R_P$.
Let $\alpha$ be the number of distinguished occurrences per character, and put
$\phi_1=1/4$ for a real $\chi_1$ and $\phi_1=1/3$ otherwise.
There is a constant $c_{\mathrm{out}}$, depending only on the fixed choices of stages (S1)--(S3) of
\cref{subsec:order} (in particular on $\theta$), such that for $q\ge q_0$ the following holds. Each
first-family character $\chi$ with a zero in $R_P$ contributes at most
\[
  e^{-At_\chi}B_{\phi_1}(p)+\eps_1+\frac{\alpha\,c_{\mathrm{out}}}{H_0C_B^2}
\]
to $W$, and a first-family character without a zero in $R_P$ contributes nothing.
\end{proposition}

\begin{proof}
Apply \cref{lem:anchored} at $\sigma=p\in\mathcal A$. For $\rho\in R_P$ we have $\lambda_\rho\ge t_\chi\ge p$ and
$|H((1-\rho)\LL)|\le e^{-At_\chi}(\Phi^\theta_p(\rho)+3\theta)$. The non-distinguished zeros of the disc
have $\Phi^\theta_p\ge0$, and so does a distinguished occurrence with $\lambda_\rho\ge p$. A distinguished
occurrence $\rho$ in the disc with $\lambda_\rho<p$ has $\lambda_\rho-p\in[a-p,0)\subseteq[-C_P,0]$ and
$|\mu_\rho|>C_B$. Since $f^\theta_p$ is smooth on $[0,T]$ and vanishes near $T$, two integrations
by parts give, for $x\in[-C_P,0]$ and real $Y\ne0$,
\[
  \bigl|\Re F^\theta_p(x+iY)\bigr|=\Bigl|\Re\Bigl(\frac{f^\theta_p(0)}{x+iY}+\frac{(f^\theta_p)'(0)}{(x+iY)^2}
  +\frac1{(x+iY)^2}\int_0^T(f^\theta_p)''(t)e^{-(x+iY)t}\,dt\Bigr)\Bigr|\le\frac{c_{\mathrm{out}}}{Y^2},
\]
where $c_{\mathrm{out}}=C_Pf^\theta_p(0)+|(f^\theta_p)'(0)|+e^{C_PT}\|(f^\theta_p)''\|_1$, maximized over the
finitely many anchors $p$ in use; the first term is bounded by $f^\theta_p(0)|x|/Y^2$. So the
distinguished terms subtracted in \cref{lem:anchored} total at least $-\alpha c_{\mathrm{out}}/C_B^2$,
and summing as in \cref{prop:envelope}, with \eqref{eq:thetachoice}, gives the claim.
\end{proof}

In the leaf programs these contributions are the \emph{hidden columns}: at most one per
first-family character, carrying the value $t_\chi$, the objective $e^{-At_\chi}B_{\phi_1}(p)$ and,
because the zero realizing $t_\chi$ has $|\gamma|\le1$, the far weight $w(t_\chi)$
(\cref{subsec:leafmodel}).

\subsection{The objective and the error allowance}\label{subsec:allowance}
In a leaf of the case tree (\cref{subsec:vocabulary}) the unknown characters are described as follows (\cref{subsec:leafmodel}): the
first family, possibly a \emph{reserved} second family of $n_2\in\{1,2\}$ characters with
height-one representatives in a given interval $[\lo_2,\hi_2]$, and the \emph{ordinary}
characters, described by their $T^*$-representatives. By
\cref{cor:ordinarycost,prop:firstinside,prop:firstoutside},
\begin{equation}\label{eq:Wsplit}
  W\le\mathcal J_1+\sum_{\chi\ \mathrm{reserved}}\mathcal G_{\phi_2}(\nu_\chi)
  +\sum_{\chi\ \mathrm{ordinary}}\mathcal G(\lambda_\chi)+E,
\end{equation}
where $\phi_2=\frac14$ if the reserved family is a single real character and $\phi_2=\frac13$
otherwise, $\nu_\chi$ is the parameter of the height-one representative of a reserved
character. The first-family term $\mathcal J_1$ is $J_1$ for an inside leaf and
$\sum_\chi e^{-At_\chi}B_{\phi_1}(p)$ for an outside leaf. In the latter case it is
represented by hidden columns, rather than by the constant part of the leaf objective
(\cref{subsec:leafmodel}).

There are two sources of error in \eqref{eq:Wsplit}. The envelope errors total less
than $19\eps_1$ for the ordinary characters, by \eqref{eq:Kfar}, and at most
$4\eps_1$ for the first and reserved families. The outside first-family bound has the
additional error $4c_{\mathrm{out}}/(H_0C_B^2)$. We choose
\[
  \eps_1=\frac\eta{23},\qquad
  \frac{4c_{\mathrm{out}}}{H_0C_B^2}\le\eta,
\]
so $E\le2\eta$. The error $\eta$ in the prime-detection criterion is accounted for
separately in \cref{prop:detect}.

Every leaf objective includes an allowance of $5\eta$. Consequently its value bounds
$W+3\eta$, and therefore also $W+2\eta$, as required by the positivity criterion.
These choices are made uniformly over the finite collection of leaves in
\cref{subsec:order}.

% !TEX root = ../linnik399.tex
\section{Far density and representatives}\label{sec:far}

The per-character costs must be summed over many characters. A weighted density
estimate supplies a common budget for that sum. Its zeros have physical height at most
$1$, so we also specify how to choose representatives within this height range.

\subsection{Xylouris's weighted density lemma}
The following is Xylouris's improvement \cite[Lemma~5.1, (5.19), p.~66]{Xylouris2011nullstellen} of
Heath-Brown's Lemma~11.1 \cite{HeathBrown1992zero}; the latter is the case $J=1$, $w_0\equiv1$.
(Xylouris writes $M$ for our $J$; we reserve $M$ for a separation constant.)

\begin{inp}[Far density]\label{inp:far}
Let $\eps,c_1,c_2>0$, $J\in\NN$, and $\alpha_1,\dots,\alpha_J\ge0$ with $\sum_i\alpha_i=1$. Put
\[
  x=\tfrac23+3c_1+c_2,\qquad u_i=\tfrac13+2c_1+\frac{ic_2}{J}\quad(0\le i\le J),
\]
and let $w_0:[u_0,x]\to\RR$ be continuous and continuously differentiable except at finitely many
points, with $1\ll w_0\ll1$ and $w_0'\ll1$. Choose for each nonprincipal character $\chi$ modulo
$q$ at most one zero $\rho_\chi$ of $L(s,\chi)$ with $|\Im\rho_\chi|\le1$ and
$\lambda_\chi\le\lambda_0=\frac13\log\log\LL$. Then for $q\ge q_0$, with $q_0$ depending on all
the parameters,
\[
  \sum_\chi\Bigl(\int_{u_0}^xw_0(t)^2e^{2\lambda_\chi t}\,dt\Bigr)^{-1}\le
  \frac{J^2+\eps}{c_1c_2^2}\sum_{i=1}^J\alpha_i^2\int_{u_{i-1}}^xw_0(t)^{-2}
  \min\{t-u_{i-1},\,u_i-u_{i-1}\}\,dt .
\]
\end{inp}

In \cite{Xylouris2011nullstellen} the statement is made for the zeros $\rho^{(k)}$ that are chosen
for the characters counted by $N(\lambda_0)$, one zero per character \cite[\S3.2.2,
pp.~25--26]{Xylouris2011nullstellen}, \cite[\S11]{HeathBrown1992zero}; any such choice is allowed, and
the estimates in the proof \cite[pp.~67--72]{Xylouris2011nullstellen} are uniform over
these choices. This uniformity is needed because the selected zeros, including the
$T^*(q)$-representatives below, depend on $q$. A nonreal character and
its conjugate are different characters, and each contributes its own term; the principal
character is excluded, since $L(s,\chi_0)$ has no zero in the relevant region for large $q$
\cite[p.~21]{Xylouris2011nullstellen}. The hypotheses on $w_0$ are satisfied by our profiles, which
are those of Xylouris's own application \cite[(6.20)--(6.21), p.~80]{Xylouris2011nullstellen}.

\subsection{The two weight profiles}
We use two profiles of the shape \cite[(6.20)]{Xylouris2011nullstellen}, both with $J=10$ and
$\eps_0=10^{-7}$:
\[
  w_0(t)^2=e^{-\theta_{\mathrm{far}}t}\sqrt{\min(t-u_0+\eps_0,\ c_2+\eps_0)}\qquad(u_0\le t\le x).
\]
\emph{Profile~I} has $c_1=0.09035$, $c_2=0.235968$ and $\theta_{\mathrm{far}}=1.28683$, and
\emph{profile~II} has $c_1=0.0821922$, $c_2=0.2170903$ and $\theta_{\mathrm{far}}=1.4964274$. The weights
$\alpha_i$ are the exact decimals of \cref{tab:alpha}; in each profile they sum to $1$. Both
profiles satisfy the hypotheses of \cref{inp:far}: $w_0$ is continuous and positive on
$[u_0,x]$, and it is smooth except at $t=u_0+c_2$.

\begin{table}[ht]
\centering
\small
\begin{tabular}{rll}
\toprule
$i$ & profile~I $\alpha_i$ & profile~II $\alpha_i$ \\
\midrule
1 & 0.0788827218 & 0.0806195583 \\
2 & 0.0849386148 & 0.0862705300 \\
3 & 0.0895629779 & 0.0905489307 \\
4 & 0.0938231516 & 0.0944644867 \\
5 & 0.0979710491 & 0.0982543287 \\
6 & 0.1021284916 & 0.1020315591 \\
7 & 0.1063732939 & 0.1058668825 \\
8 & 0.1107651869 & 0.1098131140 \\
9 & 0.1153565492 & 0.1139152197 \\
10 & 0.1201979632 & 0.1182153903 \\
\bottomrule
\end{tabular}
\caption{The weights $\alpha_i$ of the two far profiles.}
\label{tab:alpha}
\end{table}

Put
\[
  w(\lambda)=\Bigl(\int_{u_0}^xw_0(t)^2e^{2\lambda t}\,dt\Bigr)^{-1},
\]
a positive decreasing function of $\lambda$. Let $V$ be the right-hand side of \cref{inp:far}
with $\eps=0$. Rigorous enclosures give
\[
  V=175.26640331\ldots\ \text{(profile I)},\qquad V=243.32098105\ldots\ \text{(profile II)}.
\]
Taking $\eps=\eta J^2$ in \cref{inp:far}, we obtain for $q\ge q_0$ and every choice of zeros as in
\cref{inp:far}:
\begin{equation}\label{eq:farbudget}
  \sum_\chi w(\lambda_\chi)\le(1+\eta)V .
\end{equation}
Each leaf of the case tree uses one of the two profiles, recorded in its data. Both
are fixed finite choices within \cref{inp:far}.

\begin{lemma}[Comparison of the prime weight and the far weight]\label{lem:decay}
For both profiles, $x\le1.1737$. Hence $A>2x$ for $L=3.99$, and for every $\lambda\ge0$ the function
$\lambda\mapsto e^{-A\lambda}/w(\lambda)$ is decreasing. In particular
$e^{-A\lambda}\le w(\lambda)/w(0)$.
\end{lemma}

\begin{proof}
The values are $x=1.1736846\ldots$ and $x=1.1303335\ldots$, so $2x<2.348<A=3.156342$. Now
$e^{-A\lambda}/w(\lambda)=\int_{u_0}^xw_0(t)^2e^{(2t-A)\lambda}\,dt$, and $2t-A<0$ on
$[u_0,x]$.
\end{proof}

With $w(0)=10.7054\ldots$ (profile~I) and $w(0)=13.1730\ldots$ (profile~II), \cref{lem:decay} and
\eqref{eq:farbudget} give, for any selection satisfying \cref{inp:far},
\begin{equation}\label{eq:Kfar}
  \sum_\chi e^{-A\lambda_\chi}\le\frac{(1+\eta)V}{w(0)}=K_{\mathrm{far}},\qquad
  K_{\mathrm{far}}<16.38\ \text{(profile I)},\quad K_{\mathrm{far}}<18.48\ \text{(profile II)}.
\end{equation}
This bound controls errors that are summed over unboundedly many characters.

\subsection{Representatives}\label{subsec:representatives}
The leaf programs describe an ordinary character through one parameter, the parameter of a
\emph{representative} zero. For the near-density argument of \cref{sec:near} the representative
must be separated in height from any zero with a smaller parameter, and we arrange this by
choosing the height of the strip in which representatives are taken.

Fix $s_{\max}=3$; every anchor used in the near rows is at most $2.3$ (\cref{sec:rows}), and every
ordinary entry uses an anchor at most $s_{\max}$. Let $K_0$ be an upper bound for
the number of zeros, of all nonprincipal characters modulo $q$ and counted with multiplicity,
with $\lambda\le s_{\max}$ and $|\gamma|\le2$. By a log-free zero-density estimate
(\cref{inp:logfree}), $K_0$ may be taken independent of $q$. Next fix $M$ as in \cref{sec:join};
it depends on $K_0$.

\begin{lemma}[A strip boundary separated from low-parameter zeros]\label{lem:Tstar}
Let $K_0$ bound the number of zeros with $\lambda\le s_{\max}=3$ and $|\gamma|\le2$,
as above, and fix $M>0$. If $\LL>4M(K_0+1)$, there is $T^*\in[1/2,1]$ such that no nonprincipal $L(s,\chi)$ has a zero with
$\lambda\le s_{\max}$ and $\bigl||\mu|-T^*\LL\bigr|<M$. We take $T^*=T^*(q)$ to be the least such number.
\end{lemma}

\begin{proof}
The at most $K_0$ zeros in question exclude at most $K_0$ open intervals of $T$ of length $2M/\LL$
each. The interval $[1/2,1]$ has length $1/2>2MK_0/\LL$, so the remaining set is a nonempty
finite union of closed intervals, and it has a least element.
\end{proof}

\begin{definition}[Representatives]\label{def:representative}
For a nonprincipal character $\chi$, its \emph{$T^*$-representative} is a zero of $L(s,\chi)$ in
$R(l)$ with $|\gamma|\le T^*$ and least parameter, if there is one; its parameter is $\lambda_\chi$.
Representatives are used for the characters outside the first family; in \cref{thm:large} all
characters use their height-one representatives. The \emph{height-one representative} is defined in the same way with the bound
$|\gamma|\le1$. The reserved family of a leaf (\cref{subsec:vocabulary,def:specification}) and an inside first family use
height-one representatives.
\end{definition}

Every representative has $|\gamma|\le1$, so the far bound \eqref{eq:farbudget} applies to any
selection of representatives, one per character. Representatives for conjugate characters
may be chosen as conjugate zeros, with the same parameter, because every selection
region is symmetric in height.
The $T^*$-representative of $\chi$ satisfies $\lambda_\chi\ge$ the parameter of its height-one
representative. Shrinking the strip can only increase the minimum parameter; in terms
of real parts, the selected zero can only move to the left.

\subsection{The far constraint of a leaf}
In a leaf program (\cref{subsec:vocabulary}), the characters other than the first family are placed
in bins by the parameters of their representatives, and each bin has a column
(\cref{sec:lp,sec:casetree}). The column of a bin $[\lo,\hi]$ has far weight
$W_c=\lfloor Sw(\hi)\rfloor$; since $w$ is decreasing, each character in the bin contributes at least
$w(\hi)$ to the left side of \eqref{eq:farbudget}. The
\emph{tail column} collects the ordinary characters $\chi$ with a zero in $R_P$ and $\lambda_\chi\ge R$, where
$R=\max(3,r)$ and $r$ is the leaf's lower bound for ordinary representatives. Its value is the mass
$x_{\mathrm{tail}}=\sum w(\lambda_\chi)$ over these characters, with far weight $S$. The far budget of
the leaf is
\[
  F=\bigl\lceil S(1+\eta)V\bigr\rceil-\1_{\mathrm{inside}}\,n\bigl\lfloor Sw(b)\bigr\rfloor .
\]
Here the last term is present for an inside leaf, that is when $\rho_1\in R_B$. Then
$|\gamma_1|\le C_B/\LL\le1$, so $\rho_1$ and $\overline{\rho_1}$ are the height-one representatives of $\chi_1$ and
$\bchi_1$, and they contribute $n\,w(\lambda_1)\ge n\,w(b)$ to \eqref{eq:farbudget}. A reserved second family is charged either through its own columns, which carry its far
weight, or by subtracting the fixed amount $n_2\lfloor Sw(\hi_2)\rfloor$ from $F$ (\cref{subsec:leafmodel}).

% !TEX root = ../linnik399.tex
\section{Zero location}\label{sec:location}

This section collects the lower bounds for $\lambda_1$, $\lambda'$, $\lambda_2$ and $\lambda_3$ that the
case tree uses. Most are printed results of Heath-Brown \cite{HeathBrown1992zero} and Xylouris
\cite{Xylouris2011nullstellen}; we state each in the form used, with its scope. The others are new
consequences of the same explicit formula, proved here by the positivity method; one of them uses
a sharpened conductor coefficient, which we prove in \cref{subsec:conductor}. None of the results of
this section depends on the exponent $L$.

\subsection{Conventions}\label{subsec:locconventions}
We fix once and for all an admissible height function $l(q)$ as in \cref{inp:rectangle}, for
instance the least one. All the zero-location results we quote are proved for an arbitrary
integer $l$ with $1\le l\le\LL/10$ for which $R(10l)\setminus R(l)$ contains no zero; their proofs use
$l$ only through these two properties, and their thresholds depend only on fixed data. So they
hold simultaneously for our $l(q)$.

The zeros $\rho_1,\rho_2,\rho_3$ are selected as in \cite[(3.11)--(3.12), p.~21]{Xylouris2011nullstellen}
(which follows \cite[\S6]{HeathBrown1992zero}). At step $k$, remove the characters in the
families $\mathcal F_1,\ldots,\mathcal F_{k-1}$ and choose a zero of maximal real part
in $R(l)$ among the remaining nonprincipal $L$-functions. Each real character is
removed once. This is the successive-minimum convention of \cref{def:families}.
The additional zero
$\rho'$ of the first family is selected as in \cite[pp.~21--22]{Xylouris2011nullstellen}: $\rho'=\rho_1$ if
$\rho_1$ is multiple, and otherwise $\rho'$ has maximal real part among the zeros of $L(s,\chi_1)$ in
$R(l)$ other than $\rho_1$ (and other than $\overline{\rho_1}$ if $\chi_1$ is real). Its parameter is the
$\lambda'$ of \cref{sec:notation} (Heath-Brown also writes $\lambda'$; Xylouris indexes this zero by $0$). When several zeros qualify
(ties in the real part, or several characters attaining a minimum), any of them may be selected:
the proofs of the quoted results use only the extremal property of the selection, so they hold for
every such choice. This is why the configuration sets of \cref{def:specification} quantify over the
admissible choices. All the results of this section concern zeros in $R(l)$ at any height up to
$l$, so they apply whether $\rho_1$ is inside or outside the buffer.

\emph{Printed tables.} Every printed bound is used as an exact implication for $q\ge q_0$. Each
of them was obtained by showing that a certain inequality fails with a positive margin, which
persists when the error term $\eps$ is small enough \cite[p.~43, p.~44 and footnote~2 on
p.~45]{Xylouris2011nullstellen}; for Heath-Brown's Tables~2 and~3 the printed values lie a little below
the computed roots \cite[\S8]{HeathBrown1992zero}. Heath-Brown states no convention for his
Tables~4 and~7, and we have recomputed their roots from his printed parameters; every recomputed
root exceeds the printed value, and each printed value is the root rounded down. The closest case
in Table~4 is row $1.05$, with root $1.43904$ against $1.439$, and the closest in Table~7 are rows
$0.20$ and $0.55$, with $2.01015$ against $2.01$ and $1.00036$ against $1.00$. The
rows of Table~4 are derived under the side conditions \cite[(8.6), (8.8)]{HeathBrown1992zero}; we have
checked from the printed parameters that these hold for every row from $0.35$ on (the tightest is
(8.6), with slacks $3.8\cdot10^{-6}$, $3.8\cdot10^{-6}$ and $2.3\cdot10^{-6}$ for the rows $1.05$, $1.15$ and
$1.294$). For the row $(0.3,2.293)$ they need not hold, but that row also follows from Heath-Brown's
Table~3, which gives $\lambda'\ge2.43$ for $\lambda_1\le0.30$ \cite[\S8]{HeathBrown1992zero}. These checks, and the
recomputation of the roots, were made in floating point, separately from the replay. A separate
program compares every printed number that the proof uses, in the code and in this paper, with
the text of the source pages and with the printed scope of its row; each agrees with the printed
value or is weaker than it in the safe direction.
Results stated in the form $c-\eps$ are used with an explicit reduction: $0.857<\frac67$ in
\cref{inp:lambda3}, and $1.09<\frac{12}{11}$ and $\eps=0.001$ in \cref{sec:exteriors}.

\subsection{Types and the first zero}

\begin{proposition}[Lower bounds for a nonreal first zero]\label{prop:firstlower}
For all sufficiently large $q$, the first family satisfies
\[
  \lambda_1>0.440\quad\text{if }\chi_1\text{ is nonreal},\qquad
  \lambda_1>0.628\quad\text{if }\chi_1\text{ is real and }\rho_1\text{ is nonreal}.
\]
Thus $\lambda_1\le0.44$ forces both $\chi_1$ and $\rho_1$ to be real.
These bounds are \cite[Lemma~4.5 and Table~11, pp.~60--61]{Xylouris2011nullstellen}.
\end{proposition}

Xylouris's lemma gives $\lambda_1>0.440$, $0.493$, $0.478$, $0.498$ and $0.628$ for characters of order
$\ge6$, $=5$, $=4$, $=3$ and $=2$ respectively; the minimum over orders at least $3$ is $0.440$.

\subsection{The test function}
The location inequalities use one family of test functions. We define it here so that
all parameters in the statements below are specified before they occur.
For $\gamma>0$, put
\begin{equation}\label{eq:ftest}
  f_\gamma(t)=\begin{cases}\frac{16\gamma^5}{15}(1-u)^3(1+3u+u^2),&0\le t\le2\gamma,\ u=t/(2\gamma),\\
  0,&t\ge2\gamma,\end{cases}
\end{equation}
The function $f_\gamma$ is the autocorrelation of $x\mapsto(\gamma^2-x^2)_+$, the test function of
\cite[(3.14), (3.17)--(3.18)]{Xylouris2011nullstellen}; it satisfies Conditions~1 and~2
(\cref{lem:parabolic}), and its transform $F$ is positive and decreasing on $\RR$.

\subsection{Printed zero-location inputs}

\begin{inp}[Real first zero]\label{inp:rrtables}
Let $\chi_1$ and $\rho_1$ be real. For $q\ge q_0$:
\begin{enumerate}[(a)]
  \item \cite[Tables~3 and~4, Lemmas~8.2 and~8.3]{HeathBrown1992zero}. If $\lambda_1\le t$ then
    $\lambda'\ge c$, for the $24$ pairs $(t,c)$ from $(0.3,\,2.293)$ to $(1.294,\,1.294)$ of the table.
    For a real $\rho'$ this follows from $\lambda'\ge2.427$ (Lemma~8.2).
  \item \cite[Lemma~8.4]{HeathBrown1992zero}. For every $\eps>0$: $\lambda'\ge(2-\eps)\log\lambda_1^{-1}$ if
    $\lambda_1\le0.2$ and $q\ge q(\eps)$; and $\lambda'\ge\max\bigl(\frac32\log\lambda_1^{-1},1.294\bigr)$ for every $\lambda_1$.
  \item \cite[Table~7 and Lemma~8.7]{HeathBrown1992zero}. If $\lambda_1\le t$ then $\lambda_2\ge c$, for
    the $16$ pairs $(t,c)$ from $(0.12,\,2.56)$ to $(0.745,\,0.745)$ of the table; these hold whether or
    not $\chi_2^4=\chi_0$ (the row $t=0.10$ is not used).
  \item \cite[Lemma~8.8]{HeathBrown1992zero}. $\lambda_2\ge\max\bigl((\frac{12}{11}-\eps)\log\lambda_1^{-1},0.745\bigr)$
    for $q\ge q(\eps)$.
\end{enumerate}
\end{inp}

Heath-Brown remarks that the estimates of his \S8 have not been proved for extremely small values
of $\lambda_1$ \cite[\S8]{HeathBrown1992zero}. We use (a)--(d) only for $\lambda_1$ bounded below by a
fixed positive constant: $\lambda_1\ge0.1$ in the case tree, and, in \cref{thm:small}, $\lambda_1$ at
least the fixed constant $u_{\mathrm{exc}}>0$ of \cref{sec:exteriors}.

\begin{inp}[Nonreal first character or zero]\label{inp:complextables}
Let $\chi_1$ or $\rho_1$ be nonreal. For $q\ge q_0$:
\begin{enumerate}[(a)]
  \item \cite[Table~2$'$, p.~62]{Xylouris2011nullstellen}. If $\lambda_1\le t$ then $\lambda'>c$, for the pairs
    $(0.46,1.85),\dots,(0.827,0.827)$ of the table.
  \item \cite[Table~3, p.~45]{Xylouris2011nullstellen}. If $\ord\chi_1\in\{2,3,4\}$ and $\lambda_1\le t$, then
    $\lambda'>c$, for the pairs $(0.38,2.53),\dots,(1.099,1.099)$; in particular $\lambda_1\le0.86$ implies
    $\lambda'>1.35$.
  \item \cite[Tables~6 and~7, p.~53]{Xylouris2011nullstellen}. If $\lambda_1\le t$ then $\lambda_2>c$; Table~6
    concerns a real $\chi_1$ (Fall~7), Table~7 all cases.
  \item \cite[Lemma~9.4 and Table~10]{HeathBrown1992zero}. $\lambda_2\ge0.702$ always, and
    $\lambda_1\le0.70$ implies $\lambda_2\ge0.704$.
\end{enumerate}
\end{inp}

\begin{inp}[The third family]\label{inp:lambda3}
For $q\ge q_0$:
\begin{enumerate}[(a)]
  \item \cite[Lemma~10.3]{HeathBrown1992zero}. $\lambda_3\ge\frac67-\eps$; we use $\lambda_3>0.857$.
  \item \cite[Table~8, column ``alle F\"alle'', p.~55]{Xylouris2011nullstellen}. If $\chi_1$ or $\rho_1$ is
    nonreal and $\lambda_1\le t$, then $\lambda_3>c$, for $(t,c)=(0.52,1.320)$, $(0.54,1.243)$,
    $(0.56,1.160)$, $(0.58,1.079)$, $(0.60,1.001)$, $(0.62,0.933)$.
  \item \cite[Lemma~4.4 and Table~10, p.~55]{Xylouris2011nullstellen}. If $\chi_1$ and $\rho_1$ are real, then
    $\lambda_1\in[0.44,0.60]$, $[0.60,0.70]$, $[0.70,0.80]$ imply $\lambda_3>1.176$, $1.055$, $0.952$.
  \item \cite[(4.28)--(4.29), pp.~56--57]{Xylouris2011nullstellen}. Let $\chi_1$ be nonreal and
    $\lambda_1\in[0.44,0.85]$, and let $f$ be the test function \eqref{eq:ftest} with $\gamma=\frac54$. Then
    \[
      0\le F(-\lambda_1)-F(\lambda_3-\lambda_1)-F(\lambda_2-\lambda_1)-F(0)+\tfrac76f(0)+\eps,
    \]
    provided that $\sup_{y\in\RR}\Re\{F(-\lambda_1+iy)-2F(iy)\}\le\frac16f(0)$.
  \item \cite[(4.31) and (4.34), pp.~58--59]{Xylouris2011nullstellen}. Let $\chi_1$ and $\rho_1$ be real,
    $\lambda_1\in[0.44,0.80]$ and $\lambda_2\in[0.44,1.176]$, and let $f$ be \eqref{eq:ftest} with
    $\gamma=1.04$. Then
    \[
      0\le F(-\lambda_2)-F(\lambda_3-\lambda_2)-F(0)-F(\lambda_1-\lambda_2)+\tfrac98f(0)+\eps,
    \]
    provided that $\sup_{t\in\RR}\Re\{F(-\lambda_2+it)-F(\lambda_1-\lambda_2+it)-F(it)\}\le\frac5{48}f(0)$ on this box.
\end{enumerate}
\end{inp}

In (d), the side condition is \cite[(4.29)]{Xylouris2011nullstellen}, needed when a product character in his argument
is principal; he verified it numerically, and we have verified it with interval arithmetic for
all $\lambda_1\in[0.44,0.85]$: the supremum is at most $0.00335$, while $\frac16f(0)=0.5425\ldots$. In (e),
Xylouris's argument \cite[(4.32)--(4.34), p.~59]{Xylouris2011nullstellen} needs
\[
  \sup_{t\in\RR}\Re\bigl\{F(-\lambda_2+it)-F(\lambda_1-\lambda_2+it)-F(it)\bigr\}\le\tfrac5{48}f(0)
\]
on the whole box, and he states that the supremum is below $0.10$, while $\frac5{48}f(0)=0.1351\ldots$. We have
verified this with interval arithmetic as well: over $\lambda_1\in[0.44,0.80]$, $\lambda_2\in[0.44,1.18]$ and all real
$t$ the supremum is at most $0.0031$. (Both checks enclose the transforms on a grid of heights with
step $\frac1{100}$ up to $30$, with second-derivative interpolation errors in all variables, and use a
closed-form bound beyond height $30$.) (Xylouris derives (4.28) from Heath-Brown's (10.2)--(10.6), whose derivation does not use
the assumption $\lambda_3\le\frac67$ made in \cite[\S10]{HeathBrown1992zero}; he himself applies it to
prove bounds above $\frac67$.)

\begin{proposition}[Parent rows]\label{prop:parentrows}
The $58$ \emph{parent rows} $(\mathfrak t,A_j,B_j,p_j,r_j)$ of \cref{tab:parents}, from which the roots of
the case tree are built (\cref{subsec:sourcecover}), are valid implications: for $q\ge q_0$, every
configuration of type $\mathfrak t$ with $\lambda_1\in[A_j,B_j]$ has $\lambda'\ge p_j$ and $\lambda_2\ge r_j$.
\end{proposition}

\begin{proof}
Each entry is at most the largest of the following bounds, each valid on the whole interval $[A_j,B_j]$ for
the type $\mathfrak t$: (i) the applicable rows $(t,c)$ of \cref{inp:rrtables,inp:complextables} with $t\ge B_j$; (ii) the
bound $c$ of an applicable row $(t,c)$ with $c=t$ (for $\lambda_1\le t$ the row gives it, and for
$\lambda_1>t$ the trivial bound gives $\lambda',\lambda_2\ge\lambda_1>t=c$), as for the rows
$(0.827,0.827)$ of Table~$2'$ and $(1.099,1.099)$ of Table~3 of \cite{Xylouris2011nullstellen} and $(0.745,0.745)$
of Table~7 of \cite{HeathBrown1992zero}; (iii) the bounds of \cref{inp:rrtables}(b),(d) and
\cref{inp:complextables}(d), which hold for every $\lambda_1$; and (iv) the trivial bounds
$\lambda',\lambda_2\ge\lambda_1\ge A_j$.
\end{proof}

\begin{table}[tp]
\centering
\small
\begin{tabular}{@{}lccc@{\qquad\qquad}lccc@{}}
\toprule
Type & $[A_j,B_j]$ & $p_j$ & $r_j$ & Type & $[A_j,B_j]$ & $p_j$ & $r_j$\\
\midrule
$\mathrm{rr}$ & $[0.1,0.12]$ & $2.293$ & $2.56$ & $\mathrm{rc}$ & $[0.628,0.66]$ & $1.67$ & $0.93$ \\
$\mathrm{rr}$ & $[0.12,0.14]$ & $2.293$ & $2.39$ & $\mathrm{rc}$ & $[0.66,0.7]$ & $1.59$ & $0.93$ \\
$\mathrm{rr}$ & $[0.14,0.16]$ & $2.293$ & $2.25$ & $\mathrm{rc}$ & $[0.7,0.74]$ & $1.52$ & $0.93$ \\
$\mathrm{rr}$ & $[0.16,0.18]$ & $2.293$ & $2.12$ & $\mathrm{rc}$ & $[0.74,0.78]$ & $1.46$ & $0.93$ \\
$\mathrm{rr}$ & $[0.18,0.2]$ & $2.293$ & $2.01$ & $\mathrm{rc}$ & $[0.78,0.82]$ & $1.40$ & $0.82$ \\
$\mathrm{rr}$ & $[0.2,0.25]$ & $2.293$ & $1.77$ & $\mathrm{rc}$ & $[0.82,0.86]$ & $1.35$ & $0.82$ \\
$\mathrm{rr}$ & $[0.25,0.3]$ & $2.293$ & $1.58$ & $\mathrm{rc}$ & $[0.86,0.9]$ & $1.30$ & $0.82$ \\
$\mathrm{rr}$ & $[0.3,0.348]$ & $2.195$ & $1.42$ & $\mathrm{rc}$ & $[0.9,0.94]$ & $1.25$ & $0.82$ \\
$\mathrm{rr}$ & $[0.348,0.35]$ & $2.195$ & $1.42$ & $\mathrm{rc}$ & $[0.94,0.98]$ & $1.21$ & $0.82$ \\
$\mathrm{rr}$ & $[0.35,0.4]$ & $2.108$ & $1.29$ & $\mathrm{rc}$ & $[0.98,1.02]$ & $1.17$ & $0.82$ \\
$\mathrm{rr}$ & $[0.4,0.45]$ & $2.03$ & $1.18$ & $\mathrm{rc}$ & $[1.02,1.06]$ & $1.13$ & $0.82$ \\
$\mathrm{rr}$ & $[0.45,0.5]$ & $1.958$ & $1.08$ & $\mathrm{rc}$ & $[1.06,1.099]$ & $1.099$ & $0.82$ \\
$\mathrm{rr}$ & $[0.5,0.55]$ & $1.893$ & $1.0$ & $\mathrm{rc}$ & $[1.099,1.5]$ & $1.099$ & $0.82$ \\
$\mathrm{rr}$ & $[0.55,0.6]$ & $1.832$ & $0.92$ & $\mathrm{complex}$ & $[0.44,0.58]$ & $1.36$ & $1.04$ \\
$\mathrm{rr}$ & $[0.6,0.65]$ & $1.776$ & $0.85$ & $\mathrm{complex}$ & $[0.58,0.64]$ & $1.15$ & $0.85$ \\
$\mathrm{rr}$ & $[0.65,0.7]$ & $1.724$ & $0.79$ & $\mathrm{complex}$ & $[0.64,0.66]$ & $1.08$ & $0.79$ \\
$\mathrm{rr}$ & $[0.7,0.745]$ & $1.676$ & $0.745$ & $\mathrm{complex}$ & $[0.66,0.68]$ & $1.02$ & $0.74$ \\
$\mathrm{rr}$ & $[0.745,0.75]$ & $1.676$ & $0.745$ & $\mathrm{complex}$ & $[0.68,0.70]$ & $0.96$ & $0.704$ \\
$\mathrm{rr}$ & $[0.75,0.8]$ & $1.63$ & $0.745$ & $\mathrm{complex}$ & $[0.70,0.72]$ & $0.93$ & $0.702$ \\
$\mathrm{rr}$ & $[0.8,0.85]$ & $1.587$ & $0.8$ & $\mathrm{complex}$ & $[0.72,0.74]$ & $0.91$ & $0.72$ \\
$\mathrm{rr}$ & $[0.85,0.9]$ & $1.547$ & $0.8$ & $\mathrm{complex}$ & $[0.74,0.76]$ & $0.89$ & $0.74$ \\
$\mathrm{rr}$ & $[0.9,0.95]$ & $1.509$ & $0.8$ & $\mathrm{complex}$ & $[0.76,0.78]$ & $0.86$ & $0.74$ \\
$\mathrm{rr}$ & $[0.95,1]$ & $1.473$ & $0.8$ & $\mathrm{complex}$ & $[0.78,0.80]$ & $0.84$ & $0.78$ \\
$\mathrm{rr}$ & $[1,1.05]$ & $1.439$ & $0.8$ & $\mathrm{complex}$ & $[0.80,0.82]$ & $0.83$ & $0.78$ \\
$\mathrm{rr}$ & $[1.05,1.1]$ & $1.406$ & $0.8$ & $\mathrm{complex}$ & $[0.82,1.5]$ & $0.827$ & $0.82$ \\
$\mathrm{rr}$ & $[1.1,1.15]$ & $1.375$ & $0.8$ &  & & &  \\
$\mathrm{rr}$ & $[1.15,1.175]$ & $1.36$ & $0.8$ &  & & &  \\
$\mathrm{rr}$ & $[1.175,1.2]$ & $1.346$ & $0.8$ &  & & &  \\
$\mathrm{rr}$ & $[1.2,1.225]$ & $1.331$ & $0.8$ &  & & &  \\
$\mathrm{rr}$ & $[1.225,1.25]$ & $1.318$ & $0.8$ &  & & &  \\
$\mathrm{rr}$ & $[1.25,1.275]$ & $1.304$ & $0.8$ &  & & &  \\
$\mathrm{rr}$ & $[1.275,1.294]$ & $1.294$ & $0.8$ &  & & &  \\
$\mathrm{rr}$ & $[1.294,1.5]$ & $1.294$ & $0.8$ &  & & &  \\
\bottomrule
\end{tabular}
\caption{Parent implications: for the stated type and $\lambda_1\in[A_j,B_j]$,
$\lambda'\ge p_j$ and $\lambda_2\ge r_j$. Each panel is read downwards, and each row is
valid on the whole closed interval. The left panel gives the $33$ real-first-zero rows;
the right gives the $13$ real-character/nonreal-zero rows and $12$ nonreal-character rows.}
\label{tab:parents}
\end{table}

\subsection{The positivity method}
The use of nonnegative trigonometric polynomials to locate zeros goes back to de la Vall\'ee
Poussin \cite{delaValleePoussin1899fonction}; it underlies all zero-free regions for Dirichlet
$L$-functions \cite{Gronwall1913series,Landau1918imaginar,Page1935number,McCurley1984explicit,
Kadiri2018explicit} and for $\zeta$ \cite{Kadiri2005region,MossinghoffTrudgian2015nonnegative}, and,
combined with smoothed explicit formulas, the tables of Heath-Brown and Xylouris and their
refinements \cite{LiuWang1998numerical}. All the new zero-location results come from the following
consequence of
\cite[Lemmas~3.1 and~3.4]{Xylouris2011nullstellen}. Lemma~3.4 there is a ``working version'' of
\cref{inp:X32}: if $f$ satisfies Conditions~1 and~2, $\chi\ne\chi_0$ and $s\in R(9l)$, and $A_1$ is the
set of zeros $\rho\in R(l)$ of $L(s,\chi)$ with $\Re\rho>\Re s$ (with multiplicity, at most $N$ of them)
and $A_2$ is any set of at most $N$ zeros $\rho\in R(l)$ with $\Re\rho\le\Re s$, then for every $\eps>0$
and $q\ge q_0(f,\eps,N)$
\begin{equation}\label{eq:X34}
  \sum_n\Lambda(n)\Re\frac{\chi(n)}{n^s}f\Bigl(\frac{\log n}\LL\Bigr)\le
  -\LL\sum_{\rho\in A_1\cup A_2}\Re F\bigl((s-\rho)\LL\bigr)+\frac{\varphi(\chi)}2f(0)\LL+\eps\LL .
\end{equation}

\begin{lemma}[Positivity method]\label{lem:positivity}
Let $f$ satisfy Conditions~1 and~2, with Laplace transform $F$. Fix a nonnegative
anchor $a$, and suppose $a\le\lambda_1$. Let finitely many triples
$(c_j,\chi^{(j)},t_j)$ be given, with $c_j>0$, $\chi^{(j)}$ a character modulo $q$ and $|t_j|\le9l$, such
that
\[
  \sum_jc_j\Re\bigl(\chi^{(j)}(n)n^{-it_j}\bigr)\ge0\qquad\text{for all $n$ coprime to $q$.}
\]
For each nonprincipal $\chi^{(j)}$ let $A_2(j)$ be a set of at most $N$ zeros of $L(s,\chi^{(j)})$ in $R(l)$,
with multiplicity. Then for every $\eps>0$ and $q\ge q_0$,
\[
  \sum_{\chi^{(j)}=\chi_0}c_j\Re F(-a+it_j\LL)-\sum_{\chi^{(j)}\ne\chi_0}c_j\sum_{\rho\in A_2(j)}\Re
  F\bigl((1+it_j-\rho)\LL-a\bigr)+\frac{f(0)}2\sum_{\chi^{(j)}\ne\chi_0}c_j\varphi(\chi^{(j)})+\eps\ge0 .
\]
\end{lemma}

\begin{proof}
Multiply the hypothesis by $\Lambda(n)n^{-(1-a/\LL)}f(\log n/\LL)\ge0$ and sum over $n$ (for $(n,q)>1$
every term vanishes, the principal ones included). For principal $\chi^{(j)}$ use \cref{inp:X31} at
$1-a/\LL+it_j$. For the others use \eqref{eq:X34} at $s_j=1-a/\LL+it_j$, which lies in $R(9l)$;
there $A_1=\emptyset$, because every zero in $R(l)$ of a nonprincipal $L$-function has parameter at
least $\lambda_1\ge a$. Divide by $\LL$.
\end{proof}

In every application the retained zeros satisfy $\lambda_\rho\ge a$, the transform $F$ is decreasing on
$\RR$, and the terms that we cannot evaluate are bounded over all real heights with interval
arithmetic. A necessary inequality between the parameters of the retained zeros results, and a
case is excluded when that inequality fails with a positive margin. Since there are finitely many
such inequalities, each with a fixed positive margin, one $q_0$ serves all of them.

\subsection{A sharper conductor coefficient}\label{subsec:conductor}
Heath-Brown remarks \cite[\S16, item~4]{HeathBrown1992zero} that his estimate for $\lambda_2$ in the
case of a real first character can be improved by replacing the conductor coefficient $\frac{11}{24}$
by $\frac{97}{216}$; he describes the ingredients but gives no proof. We need a weighted form of this
refinement, and we prove it here. For a modulus $q$ put
\[
  v=\prod_{p^e\parallel q,\ e\ge3}p^e,\qquad \theta_q=\frac{\log v}{\log q}\in[0,1] .
\]
We also write $\operatorname{rad}q=\prod_{p\mid q}p$.

\begin{lemma}[Growth bounds with conductor dependence]\label{lem:Lbounds}
Let $k\ge3$ and $\eps>0$. For every nonprincipal character $\chi$ modulo $q$,
\[
  |L(\sigma+it,\chi)|\ll_{\eps,k}q^{\phi_q(\chi)(1-\sigma)(1+k^{-1})+\eps}(1+|t|)\qquad
  \Bigl(1-\frac1k\le\sigma\le1+\frac{\log\LL}\LL\Bigr),
\]
where $\phi_q(\chi)=\frac14\frac{\log(4\operatorname{rad}q)}{\log q}$ if $\chi$ is real and
$\phi_q(\chi)=\min\bigl(\frac13,\frac14+\frac34\theta_q\bigr)$ otherwise.
\end{lemma}

\begin{proof}
Let $\chi$ be induced by the primitive character $\chi^*$ modulo $q^*$; then $q^*\mid q$ and
$|L(s,\chi)|\ll q^\eps|L(s,\chi^*)|$ for $\sigma\ge\frac12$ \cite[p.~10]{HeathBrown1992zero}. Heath-Brown
proves (\cite[pp.~9--10]{HeathBrown1992zero}, (2.2)--(2.4)) that a bound
$\sum_{N<n\le N+H}\chi^*(n)\ll_{\eps,k}Q^{(k+1)/(4k^2)}q^{*\eps}H^{1-1/k}$ for $1\le H\le q^*$ gives, by partial
summation with the cut $q_1=Q^{(k+1)/(4k)}$ and the P\'olya--Vinogradov inequality beyond $q^*$,
$L(s,\chi^*)\ll Q^{(1-\sigma)(1+k^{-1})/4}q^{*\eps}(1+|t|)$ in the stated range (if $q_1>q^*$ the trivial
bound up to $q^*$ is already stronger). For $\sigma\le1$ a bound of this form with a smaller $Q$ is
stronger; for $1<\sigma\le1+\log\LL/\LL$ every factor $Q^{(1-\sigma)c}$ with $1\le Q\le q^4$ and
$0\le c\le1$ lies between $\LL^{-4}$ and $1$, so all such bounds agree up to a factor $q^{\eps}$.

If $\chi$ is real, then $\chi^*$ is a real primitive character, so $q^*$ is the absolute value of a
fundamental discriminant and $q^*\le4\operatorname{rad}(q^*)\le4\operatorname{rad}(q)$. Its order is $2$, so by \cite[Lemma~2.4]{HeathBrown1992zero}, which bounds the
sums with the factor $l^{3/(2k)}$ for a character of order $l$, the character-sum bound holds with
$Q=q^*$ for every $k$. This gives the real case.

In general, \cite[Lemma~2.3]{HeathBrown1992zero} gives the bound with
$Q^{(k+1)/(4k^2)}=v^{*(3k-1)/(4k^2)}q^{*(k+1)/(4k^2)}\le(v^{*3}q^*)^{(k+1)/(4k^2)}$, where $v^*$ is the
cube-full part of $q^*$; since $q^*\mid q$ we have $v^*\mid v$ and $v^{*3}q^*\le v^3q$. This gives the
exponent $\frac14(1+\frac{3\log v}{\log q})=\frac14+\frac34\theta_q$. The exponent $\frac13$ is
\cite[Lemma~2.5]{HeathBrown1992zero}. Taking the smaller bound gives the general case. (Bounds that are uniform in $q$ and $t$
together are studied in \cite{HeathBrown1978hybrid}; here the factor $1+|t|$ is harmless, since
$|t|\le\LL+1$.)
\end{proof}

\begin{lemma}[The refined coefficient in the explicit formula]\label{lem:explicitphi}
For every $\eps>0$, \cref{inp:X32} and \eqref{eq:X34} remain true, with a suitable $q_0$, when $\varphi(\chi)$ is replaced by
$\phi_q(\chi)+\eps$, with the radius $\delta$ that the proof of \cref{inp:X32} provides for $\varphi=\frac13$ (which
serves all characters).
\end{lemma}

\begin{proof}
Both are deduced from Heath-Brown's Lemma~3.1 \cite[pp.~12--14, 23--24]{HeathBrown1992zero}
\cite[pp.~19--21]{Xylouris2011nullstellen}, whose proof uses the bound (2.5) of his Lemma~2.5 only through
the inequality
\[
  \log|L(s_0+Re^{i\vartheta},\chi)|\le\bigl\{\phi(1+k^{-1})(1-\sigma_0-R\cos\vartheta)+2\eps\bigr\}\LL
\]
for $\frac\pi2\le\vartheta\le\frac{3\pi}2$, $R\le1/k$ and large $q$. By \cref{lem:Lbounds} this holds with
$\phi=\phi_q(\chi)$, with an implied constant independent of $q$ and $\chi$. The remaining choices in
that proof ($k=3+[3\phi/(2\eps_0)]$, $R=1/k$, and the radius) may be made with the upper bound
$\phi\le\frac13$ in place of $\phi$, so they do not depend on $q$. The deduction of Lemma~5.2 of
\cite{HeathBrown1992zero} (\cref{inp:X32}) and of \eqref{eq:X34} from Lemma~3.1 is unchanged.
\end{proof}

\begin{proposition}[A coupled saving in the conductor term]\label{prop:conductor}
Let $\chi_1$ be a real nonprincipal character and let $\chi_2$ be a character with $\chi_2^k\ne\chi_0$ and $\chi_1\chi_2^k\ne\chi_0$
for $k\in K$, a finite set of positive integers. Let $c_k>0$ ($k\in K$) with $\Sigma=\sum_kc_k\ge\frac19$.
Suppose the polynomial $1+\sum_{k\in K}c_k\cos(k\vartheta)$ is nonnegative for every real
$\vartheta$. In an application of \cref{lem:positivity} to
$(1+\chi_1(n))\bigl(1+\sum_{k\in K}c_k\Re\bigl((\chi_2(n)n^{-i\gamma_2})^k\bigr)\bigr)$,
the conductor term may, for every fixed $\eps>0$ and sufficiently large $q$, be taken to be
\[
  \Bigl(\frac18+\frac\Sigma3-\frac1{108}\Bigr)f(0)+\eps .
\]
\end{proposition}

\begin{proof}
The nonprincipal characters are $\chi_1$ with coefficient $1$, and $\chi_2^k$ and $\chi_1\chi_2^k$, each with
coefficient $c_k$. By \cref{lem:explicitphi} the conductor term is at most
$\frac{f(0)}2\bigl(\phi_q(\chi_1)+2\Sigma\max_\chi\phi_q(\chi)\bigr)+\eps$, the maximum being over the
characters $\chi_2^k$ and $\chi_1\chi_2^k$. Since $p\le(p^e)^{1/3}$ when $e\ge3$,
$\operatorname{rad}q\le qv^{-2/3}$, so $\phi_q(\chi_1)\le\frac14-\frac{\theta_q}6+o(1)$, and
$\phi_q(\chi)\le\min(\frac13,\frac14+\frac34\theta_q)+o(1)$ for every nonprincipal $\chi$ (for a real $\chi$
because $\frac14-\frac{\theta_q}6\le\frac14$).
The conductor term is therefore at most $\mathcal K(\theta_q)f(0)+2\eps$ for large $q$, with
\[
  \mathcal K(x)=\frac18-\frac x{12}+\Sigma\min\Bigl(\frac13,\frac14+\frac{3x}4\Bigr)\qquad(0\le x\le1).
\]
On $[0,\frac19]$ the slope of $\mathcal K$ is $\frac34\Sigma-\frac1{12}\ge0$, and on $[\frac19,1]$ it is
$-\frac1{12}$. Hence $\mathcal K(\theta_q)\le\mathcal K(\frac19)=\frac18+\frac\Sigma3-\frac1{108}$.
\end{proof}

For $\Sigma=1$ this is Heath-Brown's $\frac{97}{216}=\frac{11}{24}-\frac1{108}$.

\subsection{Lower bounds for \texorpdfstring{$\lambda_2$}{lambda2} when the first zero is real}

\begin{proposition}[Second-family bounds for a real first zero]\label{prop:realrows}
Let $\chi_1,\rho_1$ be real with $\lambda_1\in[a,b]$, where $0\le a\le b$, and let $h\ge b$. Let
$P(x)=1+c_1x+c_2(2x^2-1)+c_5(16x^5-20x^3+5x)$, that is
$P(\cos\vartheta)=1+c_1\cos\vartheta+c_2\cos2\vartheta+c_5\cos5\vartheta$, with
\[
  (c_1,c_2,c_5)=\frac{(1.41866466,\ 0.43287503,\ 0.01421037)}{1.000001},
\]
so that $P\ge0$ on $[-1,1]$ and $\Sigma=c_1+c_2+c_5\ge\frac19$.
Choose $f=f_\gamma$ from \eqref{eq:ftest}, with transform $F$ and fixed $\gamma>0$.
If $\lambda_2\le h$, then for every $\eps>0$ and all sufficiently large $q$, the inequality
corresponding to the character relation below must hold. The generic case means that
$\chi_2$ is nonreal and none of the relations in (ii)--(iii) holds.
\begin{enumerate}[(i)]
  \item (generic) $F(b-a)+c_1F(h-a)\le F(-a)+(\frac18+\frac\Sigma3-\frac1{108})f(0)+\eps$;
  \item ($\chi_2^2=\chi_1$) $F(b-a)+c_1F(h-a)\le F(-a)+(\frac{1+c_2}8+\frac{c_1+c_5}3)f(0)+\sup_\Delta R(\Delta)+\eps$,
    where $R(\Delta)=c_2\{\Re F(-a+i\Delta)-\Re F(\lambda_1-a+i\Delta)\}-c_1\Re F(\lambda_2-a+i\Delta)$, the supremum
    over $\Delta\in\RR$, $\lambda_1\in[a,b]$ and $\lambda_2\in[a,h]$;
  \item ($\chi_2^5\in\{\chi_0,\chi_1\}$) $F(b-a)+c_1F(h-a)\le(1+c_5)F(-a)+(\frac{1+c_5}8+\frac{c_1+c_2}4)f(0)+\eps$;
  \item ($\chi_2$ real) $F(b-a)+F(h-a)\le F(-a)+\frac38f(0)+\eps$, with a second test function
    \eqref{eq:ftest}.
\end{enumerate}
Here $\chi_2$ is a character of the second family. The test parameter in (iv) may be
chosen independently of the one in (i)--(iii). If all four possible cases are excluded
with a positive margin at $\eps=0$, then $\lambda_2>h$ for all sufficiently large $q$.
\end{proposition}

\begin{proof}
Put $z_n=\chi_2(n)n^{-i\gamma_2}$ and apply \cref{lem:positivity} with anchor $a$ to
$(1+\chi_1(n))P(\Re z_n)\ge0$ (with $\Re z_n^k=T_k(\Re z_n)$ for $|z_n|=1$), retaining $\rho_1$ in the term of $\chi_1$
and $\rho_2$ in the term $c_1\Re z_n$. The heights are at most $5l\le9l$. The nonzero frequencies are
$1$, $2$ and $5$, and $\chi_2\ne\chi_1$. A product character is principal only if $\chi_2$ is real, or
$\chi_2^2=\chi_1$ (order $4$), or $\chi_2^5\in\{\chi_0,\chi_1\}$ (order $5$ or $10$); the last two cannot
occur together, since together they would give $\chi_2^2=\chi_0$. Generic case: there is no principal
product, and \cref{prop:conductor} gives (i), using that $F$ is decreasing, $\lambda_1\le b$ and $\lambda_2\le h$.
Case $\chi_2^2=\chi_1$: then $\chi_1\chi_2^2=\chi_0$ at height $2\gamma_2$, $\chi_2^2=\chi_1$ at height $2\gamma_2$
retains $\rho_1$, and $\chi_1\chi_2=\bchi_2$ retains $\overline{\rho_2}$; with $\Delta=2\mu_2$ these give $R(\Delta)$,
while the characters of frequencies $1$ and $5$ are nonprincipal and are charged $\frac13$ (conservatively,
since some of them have order $4$). Case
$\chi_2^5\in\{\chi_0,\chi_1\}$: exactly one term is principal, bounded by $c_5F(-a)$ since $f\ge0$, and
every other character has order at most $10\le\LL$, so $\varphi=\frac14$. Real $\chi_2$: apply the
method to $(1+\chi_1(n))(1+\Re z_n)\ge0$, whose three nonprincipal characters are real.
\end{proof}

The rows used are given in \cref{tab:realrows}: $22$ cells of width $0.0025$ covering $[0.700,0.755]$.
For each cell $[a,b]$ we give $h$ and the test parameter $\gamma$ of the generic case; each of the four
inequalities fails by the margin shown (normalized by $f(0)$), so that $\lambda_2>h$ throughout the
cell. The margins were computed with interval arithmetic: the correlation term $R(\Delta)$ is enclosed
on a grid in $\Delta\in[0,30]$ with second-derivative interpolation errors in all variables, and by a
closed-form bound for $|\Delta|\ge30$. The nonnegativity of $P$ was verified at the points
$x=j/2000$ with a second-derivative bound, and independently by exact Bernstein expansion. In
addition, on cells inside $[0.700,0.7025]$ we use the degree-$2$ row
$P(\cos\vartheta)=(\frac78+\cos\vartheta)^2/\frac{81}{64}$ without the conductor refinement, which gives
$\lambda_2>0.762$ (\cref{prop:realfirstsecond}).

\begin{proposition}[A second-family bound near $0.70$]\label{prop:realfirstsecond}
If $\chi_1,\rho_1$ are real and $0.700\le\lambda_1\le0.7025$, then $\lambda_2>0.762$ for $q\ge q_0$.
\end{proposition}

\begin{proof}
As in \cref{prop:realrows} with $c_1=\frac{112}{81}$, $c_2=\frac{32}{81}$, $c_5=0$, the printed conductor
coefficients and $\gamma=1.09$ (and $\gamma=0.82$ for a real $\chi_2$); the normalized margins of the
generic, order-$4$ and real cases are at least $3.6\cdot10^{-4}$, $0.082$ and $0.033$.
\end{proof}

\begin{table}[ht]
\centering
\scriptsize
\begin{tabular}{@{}cccccccc@{}}
\toprule
cell $[a,b]$ & $h$ & $\gamma$ & generic & order $4$ & order $5,10$ & real & \\
\midrule
$[.7000,.7025]$ & $.7793$ & $1.09624$ & $1.02\cdot10^{-4}$ & $.0804$ & $.1280$ & $.0289$\\
$[.7025,.7050]$ & $.7783$ & $1.09522$ & $5.86\cdot10^{-5}$ & $.0803$ & $.1279$ & $.0287$\\
$[.7050,.7075]$ & $.7772$ & $1.09419$ & $8.66\cdot10^{-5}$ & $.0804$ & $.1280$ & $.0285$\\
$[.7075,.7100]$ & $.7762$ & $1.09317$ & $4.92\cdot10^{-5}$ & $.0803$ & $.1279$ & $.0284$\\
$[.7100,.7125]$ & $.7751$ & $1.09216$ & $8.35\cdot10^{-5}$ & $.0804$ & $.1280$ & $.0282$\\
$[.7125,.7150]$ & $.7741$ & $1.09115$ & $5.21\cdot10^{-5}$ & $.0803$ & $.1279$ & $.0280$\\
$[.7150,.7175]$ & $.7730$ & $1.09014$ & $9.25\cdot10^{-5}$ & $.0804$ & $.1279$ & $.0278$\\
$[.7175,.7200]$ & $.7720$ & $1.08914$ & $6.71\cdot10^{-5}$ & $.0804$ & $.1279$ & $.0276$\\
$[.7200,.7225]$ & $.7710$ & $1.08814$ & $4.46\cdot10^{-5}$ & $.0803$ & $.1279$ & $.0275$\\
$[.7225,.7250]$ & $.7699$ & $1.08715$ & $9.41\cdot10^{-5}$ & $.0804$ & $.1279$ & $.0273$\\
$[.7250,.7275]$ & $.7689$ & $1.08615$ & $7.75\cdot10^{-5}$ & $.0804$ & $.1279$ & $.0271$\\
$[.7275,.7300]$ & $.7679$ & $1.08517$ & $6.39\cdot10^{-5}$ & $.0804$ & $.1279$ & $.0269$\\
$[.7300,.7325]$ & $.7669$ & $1.08418$ & $5.31\cdot10^{-5}$ & $.0804$ & $.1279$ & $.0267$\\
$[.7325,.7350]$ & $.7659$ & $1.08321$ & $4.53\cdot10^{-5}$ & $.0804$ & $.1279$ & $.0265$\\
$[.7350,.7375]$ & $.7649$ & $1.08223$ & $4.03\cdot10^{-5}$ & $.0804$ & $.1278$ & $.0263$\\
$[.7375,.7400]$ & $.7639$ & $1.08126$ & $3.83\cdot10^{-5}$ & $.0804$ & $.1278$ & $.0261$\\
$[.7400,.7425]$ & $.7629$ & $1.08029$ & $3.91\cdot10^{-5}$ & $.0804$ & $.1278$ & $.0260$\\
$[.7425,.7450]$ & $.7619$ & $1.07933$ & $4.28\cdot10^{-5}$ & $.0804$ & $.1278$ & $.0258$\\
$[.7450,.7475]$ & $.7609$ & $1.07837$ & $4.93\cdot10^{-5}$ & $.0804$ & $.1278$ & $.0256$\\
$[.7475,.7500]$ & $.7599$ & $1.07741$ & $5.87\cdot10^{-5}$ & $.0804$ & $.1278$ & $.0254$\\
$[.7500,.7525]$ & $.7589$ & $1.07645$ & $7.09\cdot10^{-5}$ & $.0804$ & $.1278$ & $.0252$\\
$[.7525,.7550]$ & $.7579$ & $1.07550$ & $8.60\cdot10^{-5}$ & $.0804$ & $.1278$ & $.0250$\\
\bottomrule
\end{tabular}
\caption{The degree-$5$ rows of \cref{prop:realrows}: on each cell, $\lambda_2>h$. The values of $\gamma$
are rounded to five decimals (the exact rational values are in the data). The margins are
normalized by $f(0)$. The real case uses $\gamma=0.82$.}
\label{tab:realrows}
\end{table}

\subsection{Lower bounds for \texorpdfstring{$\lambda_2$}{lambda2} and \texorpdfstring{$\lambda'$}{lambda'} when the first character is nonreal}
For $t>0$ put
\[
  P_t(e^{i\vartheta})=\frac{(t+\cos\vartheta)^2}{t^2+\frac12}=1+c_1\cos\vartheta+c_2\cos2\vartheta,\qquad
  c_1=\frac{2t}{t^2+\frac12},\quad c_2=\frac{1/2}{t^2+\frac12},
\]
so that $c_1/c_2=4t$.

\begin{proposition}[Second-family bounds for a nonreal first character]\label{prop:complexrows}
Let $\chi_1$ be nonreal, let $\lambda_1\in[a,b]$ with $0\le a\le b$, and let $h\ge b$.
Fix $t>0$, let $c_1,c_2$ be the coefficients of $P_t$ above, and let
$f=f_\gamma$ with $\gamma>0$ and transform $F$.
Suppose $\lambda_2\le h$. For every $\eps>0$ and sufficiently large $q$, the following
necessary inequalities hold:
\begin{enumerate}[(i)]
  \item if $\chi_2$ is nonreal and
    \begin{equation}\label{eq:alias}
      \sup_{\lambda\in[a,h],\,y\in\RR}\Re\{F(-a+iy)-4tF(\lambda-a+iy)\}\le\tfrac16f(0),
    \end{equation}
    then $c_1\{F(b-a)+F(h-a)\}\le F(-a)+\frac{(1+c_1+c_2)^2-1}6f(0)+\eps$;
  \item if $\chi_2$ is real, choose $t'>0$ and let $c_1',c_2'$ be the coefficients of $P_{t'}$.
    If \eqref{eq:alias} holds with $t'$ in place of $t$ and $[a,b]$ in place of
    $[a,h]$, then $c_1'F(b-a)+F(h-a)\le F(-a)+(\frac18+\frac{c_1'+c_2'}3)f(0)+\eps$.
\end{enumerate}
\end{proposition}

\begin{proof}
(i) Apply \cref{lem:positivity} to $P_t(\chi_1(n)n^{-i\gamma_1})P_t(\chi_2(n)n^{-i\gamma_2})\ge0$, expanded as
\[
  \sum_{w\in\{-2,\dots,2\}^2}a_{|w_1|}a_{|w_2|}\,\chi^w(n)\,n^{-iw\cdot\gamma},\qquad
  a_0=1,\ a_1=\tfrac{c_1}2,\ a_2=\tfrac{c_2}2,\ \chi^w=\chi_1^{w_1}\chi_2^{w_2}.
\] The indices $(\pm1,0)$ retain $\rho_1$ (or $\overline{\rho_1}$) and $(0,\pm1)$ retain $\rho_2$
(or $\overline{\rho_2}$). If no nonzero index is principal, the conductor term is at most
$\frac16((1+c_1+c_2)^2-1)f(0)$. A nonzero principal index $w$ is not on an axis, since $\chi_1$ and $\chi_2$
are nonreal, and does not have $|w_1|,|w_2|\le1$, since $\chi_2\notin\{\chi_1,\bchi_1\}$. So some $|w_i|=2$;
choose such $i$, with $i=1$ in case of a tie, and put $T(w)=w-\operatorname{sgn}(w_i)e_i$. Then
$\chi^{T(w)}=\chi_i^{-\operatorname{sgn}w_i}$ is nonprincipal and not an axis index of length one, its
coefficient is $4t$ times that of $w$, and the zero of $\chi_i^{-\operatorname{sgn}w_i}$ selected at the axis
index sits at relative height $w\cdot\mu$. The map $T$ is injective: a collision could only come from
$(2s,\sigma)$ and $(s,2\sigma)$ with $s,\sigma=\pm1$, and if both were principal then $\chi_1^s=\chi_2^\sigma$,
which is excluded. Retaining at $T(w)$ the zero of $\chi_i^{-\operatorname{sgn}w_i}$, the principal term at $w$
and the extra retained term together change the inequality by at most
$a_w\{\Re F(-a+iy)-\frac16f(0)-4t\Re F(\lambda_i-a+iy)\}\le0$, by \eqref{eq:alias}. (ii) Use
$P_{t'}(\chi_1(n)n^{-i\gamma_1})\bigl(1+\Re\chi_2(n)n^{-i\gamma_2}\bigr)\ge0$. Now $\chi_2$ is real, with coefficient
$1$ and $\varphi=\frac14$; the only possible principal indices are $(\pm2,\pm1)$, which are matched to
$(\pm1,\pm1)$ with target $\overline{\rho_1}$ or $\rho_1$, as before.
\end{proof}

\begin{proposition}[Additional-zero bounds for a nonreal first character]\label{prop:polyrows}
Let $\chi_1$ be nonreal of order at least $5$, let $\lambda_1\in[a,b]$ with $0\le a\le b$, and suppose
$\lambda'\le h$. Fix $t,\gamma>0$, let $c_1,c_2$ be the coefficients of $P_t$, and let
$F$ be the transform of $f=f_\gamma$. Choose an admissible $\rho'$ of $\chi_1$ with
parameter $\lambda'$ and height $\gamma'$. Put
$x=\lambda_1-a$, $y=\lambda'-a$ and $\Delta=\LL(\gamma_1-\gamma')$, and
\[
  M_t=(1+c_1+c_2)^2-1-\tfrac12(c_1^2+c_2^2),\quad p_1=\tfrac12c_1^2,\quad p_2=\tfrac12c_2^2,\quad
  q_1=c_1(1+\tfrac12c_2),\quad q_2=\tfrac12c_1c_2 .
\]
Then for every $\eps>0$ and all sufficiently large $q$,
\[
  c_1[F(x)+F(y)]\le F(-a)+\tfrac{M_t}6f(0)+R(\Delta)+\eps,
\]
where
\begin{align*}
 R(\Delta)={}&p_1\Re F(-a+i\Delta)+p_2\Re F(-a+2i\Delta)\\
 &-q_1\Re\{F(x+i\Delta)+F(y+i\Delta)\}\\
 &-q_2\Re\{F(x+2i\Delta)+F(y+2i\Delta)\}.
\end{align*}
\end{proposition}

\begin{proof}
Apply \cref{lem:positivity} to $P_t(\chi_1(n)n^{-i\gamma_1})P_t(\chi_1(n)n^{-i\gamma'})\ge0$, whose terms are
$\chi_1^{k+l}$ at heights $k\gamma_1+l\gamma'$, $|k|,|l|\le2$. Since the order of $\chi_1$ is at least $5$ and
$|k+l|\le4$, the principal terms are exactly those with $k+l=0$; they have total mass $p_1$ at height
$\pm\Delta$ and $p_2$ at $\pm2\Delta$. In every term with $k+l=\pm1$ retain both $\rho_1$ and $\rho'$ (or their
conjugates); a multiple zero is retained with its multiplicity. The terms $(1,0),(0,1)$ and their
conjugates give $-c_1[F(x)+F(y)]-c_1\Re[F(x+i\Delta)+F(y+i\Delta)]$, and the terms $(2,-1),(-1,2)$ and their
conjugates give $-\frac12c_1c_2\Re[F(x+i\Delta)+F(y+i\Delta)+F(x+2i\Delta)+F(y+2i\Delta)]$.
\end{proof}

For $\chi_1$ of order $3$ or $4$, \cref{inp:complextables}(b) gives $\lambda'>1.35$ whenever
$\lambda_1\le0.86$, which is stronger than every bound in \cref{prop:polyrows} that we use.

The rows used are the following, each on a cell of width $0.0025$ and each with a positive
normalized margin. There are $25$ rows of \cref{prop:complexrows}, on cells covering
$[0.66,0.7225]$, with conclusions from $\lambda_2>0.763$ (on $[0.66,0.6625]$) to $\lambda_2>0.725$ (on
$[0.72,0.7225]$), parameters $\gamma\in\{1.25,1.3\}$, $t\in\{0.9,0.95\}$ and, for a real $\chi_2$,
test parameter $\gamma_{\mathrm r}\in\{1.1,1.15\}$ and $t'=0.85$; their generic margins are at least $8.2\cdot10^{-4}$, and the margins
in \eqref{eq:alias} at least $0.163$. There are $69$ rows of \cref{prop:polyrows}, on $69$ of the $70$ cells of width $0.0025$ in
$[0.68,0.855]$ (all except $[0.6975,0.70]$), with conclusions from $\lambda'>0.975$ (on $[0.68,0.6825]$) to $\lambda'>0.855$ (on
$[0.8525,0.855]$), parameters $\gamma\in\{1.05,1.1,1.15\}$ and $t\in\{0.9,0.95\}$, and margins at least
$1.5\cdot10^{-3}$. For these rows, if $\lambda'\le h$, the left side of \cref{prop:polyrows} is at least
$c_1[F(b-a)+F(h-a)]$, since $F$ is decreasing, and $R(\Delta)$ is replaced by its supremum over all real
$\Delta$, $x\in[0,b-a]$ and $y\in[p-a,h-a]$, where $p=\max(p_j,a)\le\lambda'$ comes from the bound $p_j$ of the
parent row of the cell (\cref{prop:parentrows}; it comes from \cite[Table~$2'$]{Xylouris2011nullstellen}) and the
trivial bound $\lambda'\ge\lambda_1\ge a$: $p=0.96$, $0.93$, $0.91$, $0.89$, $0.86$, $0.84$, $0.83$ on the cells in $[0.68,0.70]$,
$[0.70,0.72]$, \dots, $[0.80,0.82]$, $p=0.827$ on $[0.82,0.8275]$, and $p=a$ on the cells beyond. So each row proves
that $\lambda_1\in[a,b]$ and $\lambda'\ge p$ imply $\lambda'>h$, and since $\lambda'\ge p$ holds for every configuration
with $\lambda_1\in[a,b]$, $\lambda_1\in[a,b]$ implies $\lambda'>h$. The complete list is in Appendix~\ref{app:rows}. On $11$ further cells the inequalities of
\cref{prop:complexrows} are evaluated on the actual cell of the case tree.

\subsection{A positivity exclusion}

\begin{proposition}[Excluding a nearby real second character]\label{prop:positivity}
Let $\chi_1,\rho_1$ be real with $\lambda_1\in[a,b]$, and let $\chi_2\ne\chi_1$ be a real nonprincipal
character with a zero of height at most $1$ and parameter $\nu\le h$. If
\[
  F(b-a)+F(h-a)-F(-a)-\tfrac38f(0)>0
\]
for a test function \eqref{eq:ftest}, then this situation does not occur for $q\ge q_0$.
\end{proposition}

\begin{proof}
Apply \cref{lem:positivity} to $(1+\chi_1(n))(1+\Re\chi_2(n)n^{-i\gamma_2})\ge0$. Its three nonprincipal
characters $\chi_1$, $\chi_2$, $\chi_1\chi_2$ are real, so $\varphi=\frac14$; retain $\rho_1$ and the given zero of
$\chi_2$, which lies in $R(l)$. This gives $F(\lambda_1-a)+F(\nu-a)\le F(-a)+\frac38f(0)+\eps$.
\end{proof}

It excludes $197$ nodes of the case tree (\cref{subsec:vocabulary}), all of type $\mathrm{rr}$ with
a reserved real family; the smallest normalized margin is $7.8\cdot10^{-5}$.

\subsection{Applying the rows}

\begin{lemma}[Preservation of configurations under location updates]\label{lem:updates}
In the notation of \cref{def:specification}, suppose that a result of this section proves $\lambda_2>h$ for every configuration of type $\mathfrak t$ with
$\lambda_1\in[a',b']$, and let $\mathfrak s$ be a specification of type $\mathfrak t$ with cell $[a,b]\subseteq[a',b']$.
Then replacing $l_2$ and $r$ by $\max(l_2,h)$ and $\max(r,h)$ (and $\lo_2$ by the new $r$) does not
change $C(\mathfrak s)$, and $C(\mathfrak s)=\emptyset$ if $\mathfrak s$ is reserved with $\hi_2\le h$. Similarly a proof of
$\lambda'>h$ allows $p$ and $\lo'$ to be replaced by $\max(\cdot,h)$, and makes $C(\mathfrak s)$ empty if
$\hi'\le h$.
\end{lemma}

\begin{proof}
Every family other than the first has all its zeros in $R(l)$ at parameter at least $\lambda_2>h$; in
particular $\nu(\mathcal F)>h$.
\end{proof}

\subsection{The third family}

\begin{proposition}[Lower bounds for the third family]\label{prop:third}
Consider a configuration satisfying a specification of type $\mathfrak t$ with
$\lambda_1\in[a,b]$. Put $h=\hi_2$ if a family is reserved and $h=\infty$ otherwise;
then $\lambda_2\le h$ by \cref{lem:familyfacts}.
Each applicable rule below gives a strict lower bound for $\lambda_3$, for all
sufficiently large $q$. Let $\lambda_3^{\lo}$ be the maximum of the finitely many bounds
selected for the leaf. Then $\lambda_3>\lambda_3^{\lo}$.
\begin{enumerate}[(1)]
  \item $0.857$ (\cref{inp:lambda3}(a));
  \item for $\mathfrak t=\mathrm{rr}$ and $[a,b]$ inside $[0.44,0.60]$, $[0.60,0.70]$ or $[0.70,0.80]$: $1.176$,
    $1.055$, $0.952$ (\cref{inp:lambda3}(c));
  \item for $\mathfrak t\ne\mathrm{rr}$ and $b\le t$: the value $c$ of the smallest applicable row of
    \cref{inp:lambda3}(b);
  \item for $\mathfrak t=\mathrm{complex}$ and $0.44\le a\le b\le0.85$: any $t$ with
    \[
      F(t-b)+F(\min(t,h)-a)-F(-b)+F(0)-\tfrac76f(0)>10^{-5}f(0)\qquad(\gamma=\tfrac54);
    \]
  \item for $\mathfrak t=\mathrm{rr}$, $0.44\le a\le b\le0.80$ and $t\in[0.952,1.176]$: any $t$ such that
    $[a,\min(t,h)]$ is covered by intervals $[c,c']$ with
    \[
      F(-c')-F(t-c')-F(0)-F(b-c)+\tfrac98f(0)<-10^{-5}f(0)\qquad(\gamma=\tfrac{26}{25}).
    \]
\end{enumerate}
\end{proposition}

\begin{proof}
(1)--(3) are printed. For (4), suppose $\lambda_3\le t$; then $\lambda_2\le\min(t,h)$. The function
$\lambda\mapsto F(-\lambda)-F(t-\lambda)=\int f(u)e^{\lambda u}(1-e^{-tu})\,du$ is increasing, so
$F(-\lambda_1)-F(\lambda_3-\lambda_1)\le F(-\lambda_1)-F(t-\lambda_1)\le F(-b)-F(t-b)$, and
$F(\lambda_2-\lambda_1)\ge F(\min(t,h)-a)$. So the right side of \cref{inp:lambda3}(d) is below
$-10^{-5}f(0)+\eps<0$, a contradiction; the side condition \cite[(4.29)]{Xylouris2011nullstellen} holds by our verification. For (5),
suppose $\lambda_3\le t\le1.176$. Then $\lambda_2\in[a,\min(t,h)]\subseteq[0.44,1.176]$ and
$\lambda_1\in[0.44,0.80]$, so \cref{inp:lambda3}(e) applies. If $\lambda_2\in[c,c']$, the same monotonicity gives
$F(-\lambda_2)-F(\lambda_3-\lambda_2)\le F(-c')-F(t-c')$ and $F(\lambda_1-\lambda_2)\ge F(b-c)$, and the right side of
(e) is negative for small $\eps$, a contradiction; the side condition \cite[(4.34)]{Xylouris2011nullstellen} holds
by our verification.
\end{proof}

For example, on the $\mathrm{rr}$ cell $[0.7025,0.705]$ rule (5) gives $\lambda_3>1.0502$ without a
reservation, and $\lambda_3>1.0888$ with a reserved family below $0.85$; Table~10 gives $0.952$.
Rules (4) and (5) are evaluated with interval arithmetic, and a value $t$ is accepted only with a
normalized margin exceeding $10^{-5}$.

% !TEX root = ../linnik399.tex
\section{The graded near-density lemma}\label{sec:near}

The purpose of this section is to constrain many characters at once while retaining
the location of each selected zero. We first bound a quadratic form in smoothed prime
sums, called \emph{responses} (\cref{thm:graded}). We then bound each response from below
using selected zeros (\cref{lem:response}). Finally, an elementary optimization argument
introduces the single threshold used in a near row (\cref{prop:threshold}).
The analytic inputs are stated in \cref{sec:inputs}; the scope of the formalization is described
in \cref{subsec:lean}. \Cref{subsec:nearidea} explains the idea of the proof, and
\cref{ex:twolevels} shows what the resulting constraint can do that separate counts cannot.

\subsection{Data}\label{subsec:neardata}
The detector determines the response to a zero. The Gram test and sieve weight control
correlations between characters. The anchors may vary between entries, but all test
functions and allowed offsets are fixed before $q$ tends to infinity.

\begin{definition}[Near-density data]\label{def:neardatum}
A \emph{near datum} consists of the following fixed objects; Conditions~1 and~2 are
those of \cref{def:conditions}.
\begin{itemize}
  \item A \emph{detector} $f$ satisfying Condition~1, with transform $F$.
  \item A \emph{Gram test} $g$ satisfying Conditions~1 and~2, with transform $G$.
  \item Real numbers $s_1$ and $s$, called the \emph{safe anchor} and the
    \emph{reference anchor}.
  \item A \emph{sieve weight}: cells $t_0<t_1<\dots<t_m$ with $t_0>\frac13+2\eps'$ for a fixed
    $\eps'>0$, and heights $h_0,\dots,h_{m-1}\ge0$. We put
    \[
      \Xi(t)=\sum_{k<m}h_k\1_{[t_k,t_{k+1})}(t),\qquad
      \omega(t)=g(t)e^{s_1t}+\Xi(t),
    \]
    and require that $\omega$ be bounded below by a positive constant on the support of $f$.
  \item A finite set $\Delta\subset[0,\infty)$ of \emph{offsets}.
\end{itemize}
We put $\vartheta_k=\frac12(t_k-\frac13)-\eps'$ (the \emph{sieve levels}), so $0<\vartheta_k<\frac12t_k$, and
\[
  I_B=\int_0^\infty\frac{e^{2st}f(t)^2}{\omega(t)}\,dt,\qquad
  D(\delta)=\int_0^\infty g(t)e^{(s_1-2\delta)t}\,dt+\sum_{k<m}h_ke^{-2\delta t_k}
    \frac{t_{k+1}^2-t_k^2}{2\vartheta_k},\qquad d_1=\frac{g(0)}6 .
\]
Here the integrand of $I_B$ is taken to be $0$ where $f=0$. We require $I_B>0$.
\end{definition}

The number $I_B$ is the first Cauchy--Schwarz factor, $D(\delta)$ bounds a diagonal
entry at offset $\delta$, and $d_1$ bounds the correlation between distinct characters.
Heath-Brown's
Lemma~12.1 is the formal analogue of the case $g=f$, $s=s_1=\lambda_1$, $\Xi=0$ and all offsets $0$ (in that
case $\omega=fe^{s_1t}$ is not bounded below on the support of $f$, and Heath-Brown writes $f=\sqrt
f\cdot\sqrt f$ instead).

\begin{definition}[Entries and responses]\label{def:entries}
An \emph{entry} is a triple $(\chi_j,\gamma_j,\delta_j)$ of a nonprincipal character modulo $q$, a
real height $\gamma_j$ and an offset $\delta_j\in\Delta$. Its \emph{response anchor} is
$s_j=s-\delta_j$, and its \emph{response} is
\[
  R_j=-\frac1\LL\sum_{n\ge1}\Lambda(n)\,\Re\Bigl(\chi_j(n)\,n^{-(1-s_j/\LL+i\gamma_j)}\Bigr)
  f\Bigl(\frac{\log n}\LL\Bigr).
\]
For a family of entries, the \emph{pair excess} of entry $j$ is
\[
  e_j=\sum_{\substack{k\ne j\\ \chi_k=\chi_j}}\Bigl(\Re G\bigl(-\sigma_{jk}+i(\gamma_j-\gamma_k)\LL\bigr)
  -d_1\Bigr)_+,\qquad \sigma_{jk}=s_1-\delta_j-\delta_k .
\]
\end{definition}

\begin{definition}[Safe anchor]\label{def:safeanchor}
For $\Gamma'>0$, the \emph{safe-anchor hypothesis} $\mathrm{SA}(s_1,\Gamma')$ says that no nonprincipal
$L(s,\chi)$ modulo $q$ has a zero $\rho$ with $|\Im\rho|\le\Gamma'$ and $\Re\rho>1-s_1/\LL$.
\end{definition}

The pair excess $e_j$ is needed only when two entries use the same character. It records
the part of their correlation that exceeds the bound $d_1$ for distinct characters.
The safe-anchor hypothesis ensures that discarded zero terms in the Gram estimate have
the required sign.

In the application $s_1\le\lambda_1$. Then $\mathrm{SA}(s_1,2l+1)$ holds for large $q$, by the
definition of $\lambda_1$ and \cref{inp:rectangle}: a zero with $|\gamma|\le2l+1\le10l$ and
$\Re\rho>1-s_1/\LL$ lies in $R(10l)$, hence in $R(l)$, so its parameter is at least
$\lambda_1\ge s_1$.

\subsection{Statement}

\begin{theorem}[Graded near-density lemma]\label{thm:graded}
Assume \cref{inp:X31,inp:X32,inp:burgess,inp:graham}, Mertens' estimate, and the divisor
bound (the number of divisors of $q$ is $O_\eps(q^\eps)$). Fix a near datum as in \cref{def:neardatum} and a
tolerance $\eta'>0$. There is a threshold $q_0$, depending only on this fixed data, with
the following property.

For $q\ge q_0$, let $J$ be a finite index set and let
$(\chi_j,\gamma_j,\delta_j)_{j\in J}$ be entries satisfying
\begin{enumerate}[(i)]
  \item $|\gamma_j|\le\Gamma$ for some $0\le\Gamma\le\LL/3$;
  \item $\mathrm{SA}(s_1,2\Gamma+1)$;
  \item the characters $\chi_j$ are pairwise distinct whenever $\Xi\not\equiv0$.
\end{enumerate}
Then, for every choice of real coefficients $a_j\ge0$,
\[
  \Bigl(\sum_ja_jR_j\Bigr)_+^2\le(1+\eta')\,I_B\Bigl[\sum_j\bigl(D(\delta_j)-d_1+e_j\bigr)a_j^2
  +(d_1+\eta')\Bigl(\sum_ja_j\Bigr)^2\Bigr].
\]
The responses $R_j$ and pair excesses $e_j$ are those of \cref{def:entries}.
The threshold $q_0$ is uniform in the number of entries, their characters and heights,
and the coefficients $a_j$.
\end{theorem}

Two special cases of the pair excess are used.
\begin{enumerate}[(i)]
  \item \emph{Separated heights.} If $\chi_k=\chi_j$ implies $|\gamma_j-\gamma_k|\LL\ge M$ for $k\ne j$,
    and $M$ is so large that $\Re G(-\sigma+iY)\le d_1$ for $|Y|\ge M$ and $\sigma$ in the finite set
    of values $\sigma_{jk}$, then $e_j=0$ (\cref{lem:decaytransform}; this holds once $g(0)>0$).
  \item \emph{A controlled pair.} If entry $j$ shares its character with exactly one other entry
    $k$, with $\delta_j=\delta_k=0$ and $|(\gamma_j-\gamma_k)\LL|$ in a range on which
    $\Re G(-s_1+iy)\le c^{\hi}$, then $e_j\le(c^{\hi}-d_1)_+$.
\end{enumerate}

\subsection{The idea of the proof}\label{subsec:nearidea}
Each response $R_j$ is a smoothed sum over primes, twisted by $\chi_j(n)n^{-i\gamma_j}$ and
evaluated at the point $1-s_j/\LL+i\gamma_j$. By the explicit formula (\cref{lem:response} below),
every zero $\rho$ of $L(s,\chi_j)$ near $1+i\gamma_j$ contributes about
$\Re F\bigl((\lambda_\rho-s_j)+i(\gamma_j\LL-\mu_\rho)\bigr)$ to $R_j$, and this is large when $\lambda_\rho$ is small.
So characters with zeros close to $1$ have large responses.

On the other hand, all the responses are correlations of one vector, the weighted primes,
with the twisted vectors $(\chi_j(n)n^{-i\gamma_j})_n$. By the Cauchy--Schwarz inequality a
combination $\sum_ja_jR_j$ is at most the length of the prime vector, which gives the factor
$I_B$, times the length of $\sum_ja_j\chi_j(n)n^{-i\gamma_j}$ measured with the weight $\omega$. The
latter is a Gram form in the characters. Its diagonal entries are the numbers $D(\delta_j)$, and
distinct characters are nearly orthogonal on the primes: their correlation is at most
$d_1=g(0)/6$. Indeed, by the explicit formula the correlation is governed by the zeros of the
quotient character $\chi_j\bchi_k$ to the right of the safe anchor, and the safe-anchor
hypothesis says that there are none. Hence only a bounded number of characters can have large
responses.

Heath-Brown's Lemma~12.1 is the case in which all entries are tested at the same anchor. The
grading gives each entry its own anchor $s_j=s-\delta_j$. An entry tested at a larger offset
$\delta_j$ has a smaller diagonal $D(\delta_j)$, because of the factor $e^{-2\delta_jt}$, so characters
whose zeros are further from $1$ are charged less. The sieve weight $\Xi$ adds Selberg-sieve
majorants of the primes to $\omega$ on the ranges $[q^{t_k},q^{t_{k+1}})$. This makes $\omega$ larger and
the first factor $I_B$ smaller. The price is a sieve contribution to the diagonal, computed with
Graham's asymptotic, and off-diagonal sums of the quotient characters over long intervals. These
are small by Burgess's estimate provided the quotient characters are nonprincipal, that is,
provided the characters are distinct.

\subsection{A Burgess bound for all nonprincipal characters}

\begin{lemma}[Burgess bound for imprimitive characters and long intervals]\label{lem:burgessall}
Assume \cref{inp:burgess}. For every $\eps>0$ there is $C$ such that for every nonprincipal
character $\chi$ modulo $q$ and all real $N\ge0$ and $H\ge1$,
\[
  \Bigl|\sum_{N<n\le N+H}\chi(n)\Bigr|\le Cq^{1/9+\eps}\min(H,q)^{2/3}.
\]
\end{lemma}

\begin{proof}
Let $\chi$ be induced by the primitive character $\chi^*$ modulo $q^*>1$, and let $r$ be the product
of the primes dividing $q$ but not $q^*$. Then $\chi(n)=\chi^*(n)\sum_{e\mid(n,r)}\mu(e)$, so the sum
is $\sum_{e\mid r}\mu(e)\chi^*(e)\sum_{N/e<m\le(N+H)/e}\chi^*(m)$. The sum only depends on the integers
in the range, so we may take the endpoints real. In each inner sum remove complete periods of
$\chi^*$, over which $\chi^*$ sums to $0$ because $\chi^*$ is nonprincipal; what remains is a sum over at
most $\min(H/e,q^*)$ consecutive integers, which by periodicity may be shifted to start at a point
$\ge1$. \Cref{inp:burgess} bounds it by $Cq^{*1/9+\eps}\min(H/e,q^*)^{2/3}\le Cq^{1/9+\eps}\min(H,q)^{2/3}$
(for an interval of length less than $1$ the bound is trivial). By the divisor bound, the number
of $e$ is $O_\eps(q^\eps)$.
\end{proof}

\subsection{Proof of \texorpdfstring{\cref{thm:graded}}{the theorem}}
Write $t_n=\log n/\LL$. For $(n,q)=1$ put
\[
  Y_n=\sum_ja_je^{-\delta_jt_n}\chi_j(n)n^{-i\gamma_j}.
\]
Since $n^{-(1-s_j/\LL+i\gamma_j)}=n^{-1}e^{st_n}e^{-\delta_jt_n}n^{-i\gamma_j}$, and $\chi_j(n)=0$
for $(n,q)>1$,
\[
  \sum_ja_jR_j=-\Re\,\frac1\LL\sum_{(n,q)=1}\frac{\Lambda(n)}n\,e^{st_n}f(t_n)\,Y_n .
\]
\emph{Cauchy--Schwarz.} Since $\omega>0$ on the support of $f$ and $\omega\ge0$ everywhere,
\begin{equation}\label{eq:CS}
  \Bigl(\sum_ja_jR_j\Bigr)_+^2\le
  \Bigl(\frac1\LL\sum_n\frac{\Lambda(n)}n\,\frac{e^{2st_n}f(t_n)^2}{\omega(t_n)}\Bigr)
  \Bigl(\frac1\LL\sum_{(n,q)=1}\frac{\Lambda(n)}n\,\omega(t_n)\,|Y_n|^2\Bigr).
\end{equation}

\emph{The first factor.} The function $\varphi_B(t)=e^{2st}f(t)^2/\omega(t)$ is bounded, because
$\omega$ is bounded below on the support of $f$; it vanishes outside the support of $f$ and is
continuous except at finitely many points, so it is Riemann integrable. Given $\eps_B>0$ choose a
partition $0=u_0<\dots<u_K$ of an interval containing the support of $f$ whose upper Riemann
sum is at most $I_B+\eps_B/2$. By Mertens' estimate $\sum_{n\le x}\Lambda(n)/n=\log x+O(1)$ (see \cite{RosserSchoenfeld1962approximate}
for explicit forms) we have
$\LL^{-1}\sum_{q^{u_i}<n\le q^{u_{i+1}}}\Lambda(n)/n=u_{i+1}-u_i+O(1/\LL)$, so the first factor is at most
$I_B+\eps_B$ for large $q$.

\emph{The second factor.} Split it as $Q_1+Q_2$ according to $\omega=ge^{s_1t}+\Xi$. Expanding
$|Y_n|^2$ and writing $\chi_{jk}=\chi_j\bchi_k$,
\[
  Q_1=\sum_{j,k}a_ja_k\,X_{jk},\qquad
  X_{jk}=\Re\frac1\LL\sum_{(n,q)=1}\Lambda(n)\chi_{jk}(n)\,n^{-(1-\sigma_{jk}/\LL)-i(\gamma_j-\gamma_k)}\,
  g(t_n).
\]
All the points $w_{jk}=1-\sigma_{jk}/\LL+i(\gamma_j-\gamma_k)$ involved satisfy $|\Re w_{jk}-1|\le C/\LL$ and
$|\Im w_{jk}|\le2\Gamma\le\LL$, so \cref{inp:X31,inp:X32} apply uniformly. The error terms
$o(1)$ below are uniform in the characters, heights and entries.
\begin{itemize}
  \item \emph{Diagonal, $j=k$.} Here $\chi_{jj}=\chi_0$ and the height is $0$. By \cref{inp:X31},
    \[
      X_{jj}=\int_0^\infty g(t)e^{(s_1-2\delta_j)t}\,dt+o(1).
    \]
  \item \emph{Distinct characters.} Here $\chi_{jk}$ is nonprincipal. By \cref{inp:X32},
    \[
      X_{jk}\le-\sum_{\rho}\Re G\bigl((w_{jk}-\rho)\LL\bigr)+\frac{\varphi(\chi_{jk})}2\,g(0)+o(1),
    \]
    the sum running over the zeros of $L(s,\chi_{jk})$ in a disc of radius $\delta<1$ about
    $1+i(\gamma_j-\gamma_k)$. Such a zero has $|\Im\rho|\le2\Gamma+1$, so by $\mathrm{SA}(s_1,2\Gamma+1)$ it
    has $\Re\rho\le1-s_1/\LL\le1-\sigma_{jk}/\LL=\Re w_{jk}$. Hence $\Re(w_{jk}-\rho)\LL\ge0$, and each term
    $\Re G(\cdot)$ is $\ge0$ by Condition~2. Dropping them and using $\varphi\le\frac13$, $g(0)\ge0$,
    we get $X_{jk}\le d_1+o(1)$. When $\sigma_{jk}<0$ the point $w_{jk}$ lies to the right of $1$, and
    the same argument applies.
  \item \emph{The same character, $j\ne k$.} Here $\chi_{jk}=\chi_0$ at the height
    $\gamma_j-\gamma_k$. By \cref{inp:X31},
    $X_{jk}=\Re G(-\sigma_{jk}+i(\gamma_j-\gamma_k)\LL)+o(1)$.
\end{itemize}
For $a\ge0$ the off-diagonal part is therefore at most
\[
  \sum_{j\ne k}a_ja_k\bigl(d_1+\epsilon_{jk}\bigr)+o(1)\Bigl(\sum_ja_j\Bigr)^2,\qquad
  \epsilon_{jk}=\1_{\chi_j=\chi_k}\bigl(\Re G(-\sigma_{jk}+i(\gamma_j-\gamma_k)\LL)-d_1\bigr)_+ .
\]
Since $\epsilon_{jk}=\epsilon_{kj}$ ($\Re G$ is even in the imaginary part, $g$ being real) and
$2a_ja_k\le a_j^2+a_k^2$, we have $\sum_{j\ne k}a_ja_k\epsilon_{jk}\le\sum_je_ja_j^2$. Also
$\sum_{j\ne k}a_ja_k=(\sum_ja_j)^2-\sum_ja_j^2$. Hence
\begin{equation}\label{eq:Q1}
  Q_1\le\sum_j\Bigl(\int_0^\infty g(t)e^{(s_1-2\delta_j)t}\,dt-d_1+e_j\Bigr)a_j^2
  +\bigl(d_1+o(1)\bigr)\Bigl(\sum_ja_j\Bigr)^2 .
\end{equation}

\emph{The sieve part $Q_2$.} It is present only when $\Xi\not\equiv0$, so the characters are distinct
and all pair excesses vanish. Fix a cell $k$, let $z_k=q^{\vartheta_k}$ and let
\[
  \xi_d=\mu(d)\frac{\log(z_k/d)}{\log z_k}\quad(d\le z_k),\qquad
  \nu_k(n)=\Bigl(\sum_{d\mid n,\ d\le z_k}\xi_d\Bigr)^2\ge0 .
\]
\begin{itemize}
  \item \emph{Majorant.} If $p$ is prime with $t_p\ge t_k$ then $p\ge q^{t_k}>z_k$, so the only
    divisor of $p$ up to $z_k$ is $1$ and $\nu_k(p)=\xi_1^2=1$. Hence
    $\Lambda(p)=(\log p)\,\nu_k(p)$. For $n$ not a prime power, $\Lambda(n)=0\le(\log n)\nu_k(n)$.
    Prime powers $p^e$ with $e\ge2$ and $t_{p^e}\ge t_k$ contribute
    $\ll\LL q^{-t_k/2}(\sum_ja_j)^2=o(1)(\sum_ja_j)^2$. So the contribution of cell $k$ to $Q_2$ is
    at most
    \[
      \frac{h_k}\LL\sum_{t_n\in[t_k,t_{k+1})}\frac{\log n}n\,\nu_k(n)\,|Y_n|^2+o(1)\Bigl(\sum_ja_j\Bigr)^2,
    \]
    where now $|Y_n|^2$ is expanded without the coprimality condition, with $\chi_j(n)=0$ for
    $(n,q)>1$.
  \item \emph{Diagonal.} The terms $j=k'$ of $|Y_n|^2$ are at most $a_j^2e^{-2\delta_jt_n}$ with
    $\chi_j\bchi_j\le\1$. By \cref{inp:graham}, $\sum_{n\le x}\nu_k(n)=x/\log z_k+O(x/\log^2z_k)$
    for $x\ge z_k$. Partial summation gives
    \[
      \frac1\LL\sum_{t_n\in[t_k,t_{k+1})}\frac{\log n}n\,\nu_k(n)e^{-2\delta_jt_n}
      =\int_{t_k}^{t_{k+1}}\frac t{\vartheta_k}e^{-2\delta_jt}\,dt+O\Bigl(\frac1\LL\Bigr)
      \le e^{-2\delta_jt_k}\frac{t_{k+1}^2-t_k^2}{2\vartheta_k}+O\Bigl(\frac1\LL\Bigr).
    \]
  \item \emph{Off-diagonal.} For $j\ne k'$ the character $\chi_{jk'}=\chi_j\bchi_{k'}$ is nonprincipal, since the
    characters are distinct.
    Put $W(n)=(\log n)n^{-1}e^{-(\delta_j+\delta_{k'})t_n}n^{-i(\gamma_j-\gamma_{k'})}$. Opening
    $\nu_k$, the term is
    \[
      \frac1\LL\sum_{d,e\le z_k}\xi_d\xi_e\,\chi_{jk'}([d,e])\sum_{m}\chi_{jk'}(m)\,W([d,e]m),
    \]
    where $m$ runs over an interval $[M_0,M_1)$ with $M_0=q^{t_k}/[d,e]\ge q^{t_k-2\vartheta_k}=
    q^{1/3+2\eps'}$. By \cref{lem:burgessall},
    $|\sum_{M_0\le m\le u}\chi_{jk'}(m)|\ll u^{2/3}q^{1/9+\eps}$ for every real $u\ge M_0$. We also have
    $|W(u[d,e])|\ll\LL(u[d,e])^{-1}$ and
    $|\partial_uW(u[d,e])|\ll\LL(1+|\gamma_j-\gamma_{k'}|)u^{-2}[d,e]^{-1}\ll\LL^2u^{-2}[d,e]^{-1}$.
    Abel summation, with the bound for the partial sums $\sum_{M_0\le m\le u}\chi_{jk'}(m)$ at every
    $u$, therefore bounds the inner sum by
    \[
      \ll\LL^2[d,e]^{-1}q^{1/9+\eps}M_0^{-1/3}=\LL^2q^{1/9+\eps-t_k/3}[d,e]^{-2/3}.
    \]
    Since $\sum_{d,e\le z_k}[d,e]^{-2/3}\le\sum_{g\le z_k}g^{-2/3}\bigl(\sum_{d'\le z_k/g}d'^{-2/3}\bigr)^2\le
    9\zeta(\frac43)z_k^{2/3}$ and $|\xi_d|\le1$, the term is
    \[
      \ll\LL\,q^{1/9+\eps-t_k/3+2\vartheta_k/3}=\LL\,q^{\eps-2\eps'/3},
    \]
    because $2\vartheta_k=t_k-\frac13-2\eps'$. With $\eps<\eps'/2$ this is $o(1)$ uniformly, and
    the sum over $j,k'$ is $o(1)(\sum_ja_j)^2$. (A bound for the partial sums only at the end
    point of the interval, combined with the total variation of $W$, would lose a factor
    $(M_1/M_0)^{2/3}$, which is too large for the cells used.)
\end{itemize}
Summing over the cells,
\begin{equation}\label{eq:Q2}
  Q_2\le\sum_ja_j^2\sum_{k<m}h_ke^{-2\delta_jt_k}\frac{t_{k+1}^2-t_k^2}{2\vartheta_k}
  +o(1)\Bigl(\sum_ja_j\Bigr)^2 .
\end{equation}

\emph{Conclusion.} By \eqref{eq:Q1} and \eqref{eq:Q2} the second factor is at most
$B+o(1)(\sum_ja_j)^2$, where
$B=\sum_j(D(\delta_j)-d_1+e_j)a_j^2+d_1(\sum_ja_j)^2\ge\sum_j(D(\delta_j)+e_j)a_j^2\ge0$. With the first
factor $I_B+o(1)$ and $I_B>0$, the right side of \eqref{eq:CS} is at most
$(1+\eta')I_B[B+\eta'(\sum_ja_j)^2]$ for $q$ large. \qed

\subsection{Responses and the threshold form}

\begin{lemma}[Uniform decay in the imaginary direction]\label{lem:decaytransform}
If $f$ satisfies Condition~1, then $\Re F(x+iY)\to0$ as $|Y|\to\infty$, uniformly for $x$ in
compact sets.
\end{lemma}

\begin{proof}
Integrating by parts twice, $F(z)=f(0)/z+O_K(|z|^{-2})$ for $\Re z$ in a compact set $K$. Moreover
$\Re(f(0)/z)=f(0)\Re z/|z|^2=O_K(|z|^{-2})$.
\end{proof}

\begin{lemma}[Response lemma]\label{lem:response}
Assume \cref{inp:X32}. Let $f$ satisfy Conditions~1 and~2, and fix $\sigma_{\max}>0$, an integer $K_0\ge0$, and
$\eta'>0$. There are $M>0$, $\delta\in(0,1)$ and $q_0$ such that the following holds for $q\ge q_0$,
and it remains true when $\delta$ is replaced by any smaller radius (with a suitable $q_0$).
Let $\chi$ be nonprincipal, $|\sigma|\le\sigma_{\max}$ and $|\gamma|\le\LL$. Let $\cS$ be a finite set
of distinct zeros of $L(s,\chi)$ in the disc $|1+i\gamma-\rho|\le\delta$, each retained with its
full multiplicity $m_\rho$. Suppose that every other zero of
$L(s,\chi)$ in the disc with $\lambda_\rho<\sigma$ satisfies $|\gamma\LL-\mu_\rho|\ge M$, and that there
are at most $K_0$ such zeros, counted with multiplicity. Then the response of $(\chi,\gamma,\sigma)$,
that is $R=-\LL^{-1}\sum_n\Lambda(n)\Re(\chi(n)n^{-(1-\sigma/\LL+i\gamma)})f(t_n)$, satisfies
\[
  R\ge\sum_{\rho\in\cS}m_\rho\,\Re F\bigl((\lambda_\rho-\sigma)+i(\gamma\LL-\mu_\rho)\bigr)
  -\frac{f(0)}6-\eta' .
\]
\end{lemma}

\begin{proof}
By \cref{lem:decaytransform} choose $M$ with $|\Re F(x+iY)|\le\eta'/(2K_0+2)$ for $|Y|\ge M$ and
$x\in[-\sigma_{\max},0]$. Apply \cref{inp:X32} with $\eps=\eta'/2$ at the point
$w=1-\sigma/\LL+i\gamma$, and let $\delta$ be its disc radius, or any smaller radius (see the remark after
\cref{inp:X32}). We have $(w-\rho)\LL=
(\lambda_\rho-\sigma)+i(\gamma\LL-\mu_\rho)$, so
\[
  R\ge\sum_{|1+i\gamma-\rho|\le\delta}\Re F\bigl((w-\rho)\LL\bigr)-\frac{\varphi(\chi)}2f(0)-\frac{\eta'}2 ,
\]
the sum counting multiplicity. Zeros in $\cS$ are kept. A zero outside $\cS$ with
$\lambda_\rho\ge\sigma$ has $\Re(w-\rho)\LL\ge0$ and a term $\ge0$ (Condition~2), which we drop. The
remaining zeros have $\lambda_\rho\in[0,\sigma)$ and $|\gamma\LL-\mu_\rho|\ge M$; there are at most $K_0$
of them, and each term is at least $-\eta'/(2K_0+2)$. Finally $\varphi(\chi)\le\frac13$ and
$f(0)\ge0$.
\end{proof}

\begin{proposition}[Threshold form]\label{prop:threshold}
Let $J$ be a finite index set, let $D_j>0$ and $v_j\in\RR$ for $j\in J$, and let $d>0$.
Suppose that
\[
  \Bigl(\sum_j a_jv_j\Bigr)_+^2\le\sum_jD_ja_j^2+d\Bigl(\sum_j a_j\Bigr)^2
  \qquad\text{for every }(a_j)_{j\in J}\in[0,\infty)^J.
\]
Then there is
$\tau\in[0,\sqrt d]$ with
\begin{equation}\label{eq:T}\tag{T}
  \sum_j\frac{(v_j-\tau)_+^2}{D_j}\le1-\frac{\tau^2}d .
\end{equation}
\end{proposition}

We call the $v_j$ the \emph{features}, the $D_j$ the \emph{diagonals} and $d$ the \emph{correlation term}
of the inequality, and $\tau$ its \emph{threshold}. A feature measures how strongly an entry is
detected, a diagonal what the entry costs on its own, and $d$ bounds the correlation between
distinct entries.

\begin{example}[One level and two levels]\label{ex:twolevels}
(i) Suppose that $N$ entries have the same feature $v>0$ and the same diagonal $D$, with $v^2>d$.
Taking every $a_j=1$ in the hypothesis of \cref{prop:threshold} gives $(Nv)^2\le ND+dN^2$, that is,
\[
  N\le\frac{D}{v^2-d}.
\]
This is the kind of bound that Heath-Brown's Lemma~12.1 provides: a bound for the number
$N(\lambda)$ of characters with a zero at distance at most $\lambda$, all counted alike.

(ii) Now take $D=1$ and $d=\frac14$, one entry with feature $v_1=1$ (a zero close to $1$) and two
entries with feature $v_2=\frac34$ (zeros further away). The count of (i), applied to each level
separately, allows this: it permits up to $D/(v_1^2-d)=\frac43$ entries with feature at least $v_1$,
and up to $D/(v_2^2-d)=\frac{16}5$ with feature at least $v_2$. But \eqref{eq:T} fails for every $\tau$:
for $0\le\tau\le\frac12$ its left side minus its right side is
\[
  (1-\tau)^2+2\bigl(\tfrac34-\tau\bigr)^2-\bigl(1-4\tau^2\bigr)=7\tau^2-5\tau+\tfrac98,
\]
whose discriminant $25-\frac{63}2$ is negative. (Equivalently, taking all three $a_j=1$ violates the
hypothesis: $(1+\frac32)^2=\frac{25}4>3+\frac94$.) So this configuration is excluded. A single
threshold couples the near entry and the farther ones, which separate counts cannot do. In
the leaf programs this coupling acts on all the bins of characters at once.
\end{example}

\begin{proof}
The function $\mathcal J(\tau)=\tau^2/d+\sum_j(v_j-\tau)_+^2/D_j$ is convex and $C^1$ on $[0,\infty)$, and
tends to $\infty$; let $\tau^*$ minimize it. If $\tau^*=0$ then $\mathcal J'(0)\ge0$, so all $v_j\le0$ and
$\mathcal J(0)=0$. Otherwise $\mathcal J'(\tau^*)=0$, that is $\tau^*=d\sum_ja_j$ with $a_j=(v_j-\tau^*)_+/D_j\ge0$.
Then $\sum_ja_jv_j=\sum_ja_j(v_j-\tau^*)+\tau^*\sum_ja_j=\sum_jD_ja_j^2+d(\sum_ja_j)^2=\mathcal J(\tau^*)$.
The hypothesis gives $\mathcal J(\tau^*)^2\le\mathcal J(\tau^*)$, so $\mathcal J(\tau^*)\le1$. Finally \eqref{eq:T} forces
$\tau^2\le d$.
\end{proof}

\begin{corollary}[Zero form]\label{cor:zeroform}
Fix a near datum whose detector $f$ also satisfies Condition~2. Fix a common bound
$K_0$ for the omitted zeros in \cref{lem:response}, and choose
$\sigma_{\max}\ge\max_{\delta\in\Delta}|s-\delta|$ with $\sigma_{\max}>0$.
Choose $\eta'>0$, $I_u\ge I_B$, and $D_u>0$, and put
\[
  \eta''=\eta'\min\{1,\sqrt{I_BD_u},D_u\}.
\]
For sufficiently large $q$, consider entries satisfying \cref{thm:graded} with tolerance
$\eta''$. For each entry $j$, suppose the retained set $\cS_j$ satisfies
\cref{lem:response} at anchor $s_j$, with the separation and disc radius chosen for
tolerance $\eta''$. Define
\[
  r_j=\sum_{\rho\in\cS_j}m_\rho\,\Re F\bigl((\lambda_\rho-s_j)+i(\gamma_j\LL-\mu_\rho)\bigr)-\frac{f(0)}6 ,
\]
and choose $D_j^+>0$ with $D_j^+\ge D(\delta_j)-d_1+e_j$ for every $j$.
Then there is $\tau\in[0,\sqrt d]$ satisfying \eqref{eq:T} with
\[
  v_j=\frac{r_j}{\sqrt{I_uD_u}}-\eta',\qquad D_j=\frac{(1+\eta')D_j^+}{D_u},\qquad
  d=\frac{(1+\eta')(d_1+\eta'D_u)}{D_u}\ \bigl(=(1+\eta')(d_1/D_u+\eta')\bigr).
\]
The same threshold remains valid if features $v_j$ are decreased, diagonals $D_j$ or
the correlation term $d$ are increased, or nonpositive features are replaced by $0$.
An entry with nonpositive feature may also be omitted.
\end{corollary}

\begin{proof}
Apply \cref{thm:graded,lem:response} with the tolerance $\eta''$ fixed in the statement.
The response lemma gives $R_j\ge k_j:=r_j-\eta''$, so for $a\ge0$
we have $(\sum_ja_jk_j)_+\le(\sum_ja_jR_j)_+$. Dividing the inequality of \cref{thm:graded} by
$I_BD_u$, we get for all $a\ge0$
\[
  \Bigl(\sum_ja_j\frac{k_j}{\sqrt{I_BD_u}}\Bigr)_+^2\le
  \sum_j\frac{(1+\eta'')D_j^+}{D_u}\,a_j^2+\frac{(1+\eta'')(d_1+\eta'')}{D_u}
  \Bigl(\sum_ja_j\Bigr)^2 ,
\]
where we also replaced each coefficient $D(\delta_j)-d_1+e_j$ by the larger $D_j^+$.
\Cref{prop:threshold} gives $\tau$ satisfying \eqref{eq:T} with these features, diagonals and
correlation term; only the positivity of the $D_j^+$ is needed, not that of $D(\delta_j)-d_1+e_j$. We compare with the displayed quantities.
\begin{itemize}
  \item If $r_j>0$ then $v_j\le k_j/\sqrt{I_BD_u}$, because $I_u\ge I_B$ and
    $\eta''/\sqrt{I_BD_u}\le\eta'$. If $r_j\le0$ then $v_j<0$.
  \item $(1+\eta')\ge(1+\eta'')$, so the displayed diagonals are at least the ones above.
  \item $(1+\eta')(d_1+\eta'D_u)\ge(1+\eta'')(d_1+\eta'')$, since $\eta''\le\eta'D_u$.
\end{itemize}
Finally, if \eqref{eq:T} holds at $\tau$ for features $v_j$, diagonals $D_j$ and correlation term $d$, it
also holds at $\tau$ for features $v'_j\le\max(v_j,0)$ with $v'_j\le v_j$ or $v'_j\le0$, diagonals
$D'_j\ge D_j$ and correlation term $d'\ge d$. Indeed each term $(v'_j-\tau)_+^2/D'_j$ is at most
$(v_j-\tau)_+^2/D_j$ (the former is $0$ when $v'_j\le0\le\tau$), $1-\tau^2/d\le1-\tau^2/d'$, and
$\tau\le\sqrt d\le\sqrt{d'}$.
\end{proof}

% !TEX root = ../linnik399.tex
\section{The near rows}\label{sec:rows}

A \emph{near row} is a constraint obtained by applying the threshold inequality
\eqref{eq:T} to selected characters and zeros. In a leaf program (\cref{subsec:vocabulary}),
characters are grouped by parameter into bins, each with its own column. Known entries from the
first or reserved family may instead be recorded separately as \emph{family terms}. There are
four kinds of rows, the family, shifted, graded and two-test rows; \cref{tab:rowkinds} in
\cref{subsec:rowkinds} lists their anchors, sieve weights and entries.

We first specify the test functions and normalization. We then justify each permitted
kind of entry and describe how to combine them into rows. Most rows follow from
\cref{cor:zeroform}; the two-test row in \cref{subsec:twotest} has a separate proof.
\Cref{prop:rowsvalid} is the conclusion needed later: every actual zero configuration
in a leaf satisfies each of its near rows at some threshold.

\subsection{The parabolic test function}
All detectors and Gram tests are normalized parabolic autocorrelations. For $c>0$ put
\[
  P(u)=(1-u)^3(1+3u+u^2)=1-5u^2+5u^3-u^5,\qquad
  \hat f_c(t)=\begin{cases}P(t/c),&0\le t\le c,\\0,&t\ge c.\end{cases}
\]
Thus $\hat f_{2\gamma}=f_\gamma/f_\gamma(0)$ with $f_\gamma$ as in \eqref{eq:ftest}. Write $\hat F_c$ for its
transform.

\begin{lemma}[Properties of the parabolic test functions]\label{lem:parabolic}
Let $c>0$.
\begin{enumerate}[(i)]
  \item $\hat f_c\ge0$, $\hat f_c(0)=1$, $\hat f_c$ is decreasing on $[0,c]$, and $\hat f_c$ extended by
    $0$ is twice continuously differentiable on $[0,\infty)$ with $|\hat f_c''|\le10/c^2$. So
    Condition~1 holds.
  \item $\Re\hat F_c(iy)=60c\,(cy\cos\frac{cy}2-2\sin\frac{cy}2)^2/(cy)^6\ge0$ for $y\ne0$, and
    Condition~2 holds.
  \item $x\mapsto\hat F_c(x)$ is positive and strictly decreasing on $\RR$.
  \item For $z\ne0$, $\hat F_c(z)=c\,\Phi(cz)$, where
    \[
      \Phi(w)=\frac1w-\frac{10}{w^3}+\frac{30(1+e^{-w})}{w^4}
        +\frac{120e^{-w}}{w^5}-\frac{120(1-e^{-w})}{w^6}\qquad(w\ne0).
    \]
    The singularity at $0$ is removable. Consequently, for $d\ge0$ and $|y|\ge U>0$,
    \[
      \Re\hat F_c(-d+iy)\le cR(cU,e^{cd}),\qquad-\Re\hat F_c(-d+iy)\le\frac d{U^2}+cR(cU,e^{cd}),
    \]
    where $R(u,E)=\frac{10}{u^3}+\frac{30}{u^4}+\frac{120}{u^6}+E\bigl(\frac{30}{u^4}+\frac{120}{u^5}+\frac{120}{u^6}\bigr)$.
  \item For $d\ge0$ put $\mathcal C_c(d)=\sup_{\Re z\ge-d}\bigl(-\Re\hat F_c(z)\bigr)$. It is finite,
    \[
      \mathcal C_c(d)=\max\Bigl(0,\ \sup_{y\in\RR}\bigl(-\Re\hat F_c(-d+iy)\bigr)\Bigr),
    \]
    and $\mathcal C_c(0)=0$.
\end{enumerate}
\end{lemma}

\begin{proof}
(i) is elementary: $P'(u)=-5u(2+u)(1-u)^2$, and $P(1)=P'(1)=P''(1)=0$. (ii) The function
$\hat f_c$ is a positive multiple of the autocorrelation of $x\mapsto(\gamma^2-x^2)_+$ with $c=2\gamma$, so
its transform on the imaginary axis is a positive multiple of $(\int_0^\infty(\gamma^2-x^2)_+\cos(xy)\,dx)^2$
\cite[Lemma~3.6]{Xylouris2011nullstellen}, which gives the closed form. Condition~2 then follows from
\cref{inp:comparison} with $F_2=0$, since $\hat F_c(z)\to0$ uniformly in $\Re z\ge0$. (iii)
$\hat F_c'(x)=-\int t\hat f_c(t)e^{-xt}\,dt<0$. (iv) Five integrations by parts, using the values of
$P$ and its derivatives at $0$ and $1$. For $z=-d+iy$ we have $c\Re\frac1{cz}=\frac{-d}{d^2+y^2}$, which is
$\le0$ and $\ge-d/U^2$, and each other term is bounded by the corresponding term of $cR(c|z|,e^{cd})$,
which decreases in $|z|\ge U$. (v) $-\Re\hat F_c(z-d)$ is harmonic in $\Re z>0$, continuous up to the
boundary and tends to $0$ uniformly at infinity, so its supremum over $\Re z\ge0$ is attained on
the boundary or is $0$; $\mathcal C_c(0)=0$ by Condition~2.
\end{proof}

\subsection{Row data}\label{subsec:rowdata}
Every family, shifted and graded row fixes rational numbers $\gamma_f>0$ (detector), $\gamma_g>0$ (Gram test),
$s_1\le s$ (safe and reference anchors), $t_0$ and $\mu_{\mathrm{row}}$, and uses the near datum of \cref{subsec:neardata} with
$f=\hat f_{2\gamma_f}$, $g=\hat f_{2\gamma_g}$, $\eps'=\frac1{2000}$, and $2000$ equal cells
$t_0<t_1<\dots<t_{2000}=2\gamma_f$, where $t_0>\frac13+2\eps'$. The heights $h_k\ge0$ are rational numbers.
They are chosen by a heuristic that brings $\omega$ close to the Cauchy--Schwarz optimum
$\omega\propto e^{st}f$: the height of cell $k$ is the positive part of
$e^{sm_k}f(m_k)/\sqrt{\mu_{\mathrm{row}}\kappa_k}-g(m_k)e^{s_1m_k}$ at the midpoint $m_k$, where
$\kappa_k=(t_{k+1}^2-t_k^2)/(2\vartheta_k(t_{k+1}-t_k))$ is the mean sieve cost per unit length of the cell,
computed in binary64 arithmetic and converted to a rational number with denominator at most
$10^{12}$. Any nonnegative heights are admissible, so this choice plays no role in the proofs. In
every family and shifted row $\mu_{\mathrm{row}}=10^9$ and all the computed heights are $0$, so $\Xi\equiv0$. The
offsets are the finitely many numbers $s-\sigma$ that occur for the anchors $\sigma$ of the entries. In
every row the lower bound of $\omega$ on each cell of the support of $f$ is checked to be positive; in
the rows with $\Xi\equiv0$ this holds because $\gamma_g>\gamma_f$.

We use the normalization of \cref{cor:zeroform}: $I_u$ is an upper Riemann
sum for $I_B$ over $4000$ cells, computed with outward interval arithmetic; $D_u\ge D(0)$; each entry
$j$ receives the feature $v_j=r_j/\sqrt{I_uD_u}-\eta'$ (rounded down, and replaced by $0$ if negative)
and the diagonal $(1+\eta')D_j^+/D_u$ (rounded up), where $D_j^+$ is an upper bound for
$D(\delta_j)-d_1+e_j$; and the row has the correlation term $d=(1+\eta')(d_1/D_u+\eta')$ (rounded up). Here
$\eta'=10^{-6}$. For an entry with offset $\delta$ the value $D(\lfloor\delta\rfloor_{1/200})\ge D(\delta)$ is used
(here $\lfloor\delta\rfloor_{1/200}$ is $\delta$ rounded down to a multiple of $\frac1{200}$, and $D$ is decreasing), and an
entry is omitted if its diagonal numerator does not meet the positivity margin used
in the construction. An omitted entry has feature $0$ in the row and contributes
nothing to \eqref{eq:T}. The rounding directions are justified by the last assertion of
\cref{cor:zeroform}.

\subsection{Facts available in a leaf}\label{subsec:leaffacts}
A leaf specifies intervals and lower bounds for the first few zeros. The full definition
of a specification $\mathfrak s$ and its configuration set $C(\mathfrak s)$ is in
\cref{def:specification}; the information needed for this section is listed below.
Write $\nu(\mathcal F)$ for the least parameter in a family among zeros of physical
height at most $1$, and $\nu_{T^*}(\chi)$ for the corresponding minimum at height at
most $T^*$. A reserved family minimizes $\nu$ outside the first family. It need not
be the family attaining the global parameter $\lambda_2$.

Fix a leaf with first-zero cell $[a,b]$ and a configuration in $C(\mathfrak s)$.
For all sufficiently large $q$, the following facts are available.
\begin{enumerate}[(Z1)]
  \item $\lambda_1\in[a,b]$. The \emph{safe anchor} of the leaf is $s_1=\lfloor50a\rfloor/50\le\lambda_1$.
  \item $\lambda'\ge p^*$, where $p^*=\lo'$ if there is a gap (an interval $[\lo',\hi']$ known to
    contain $\lambda'$) and $p^*=p$ otherwise; if there is a finite gap, $\lambda'\le\hi'$.
  \item Every zero in $R(l)$ of a character outside the first family has parameter at least
    $\lambda_2\ge l_2$.
  \item Every family $\mathcal F\ne \mathcal F_1$ has $\nu(\mathcal F)\ge r$, and ordinary characters have $\nu_{T^*}\ge\nu\ge r$.
  \item If $\mathfrak s$ is reserved, the reserved family has $n_2$ characters with common height-one
    parameter $\nu_*\in[\lo_2,\hi_2]$.
  \item Inside: $|\mu_1|\le C_B$. Outside: $|\mu_1|>C_B$.
\end{enumerate}

\begin{lemma}[No zero to the right]\label{lem:noright}
Let $\chi$ be nonprincipal, let $|\gamma|\le2l$ and $0<\delta\le1$, and fix a real anchor
$\sigma$. For sufficiently large $q$, every zero $\rho$ of $L(s,\chi)$ in the disc
$|1+i\gamma-\rho|\le\delta$ has $\lambda_\rho\ge\sigma$ in either of the following cases:
\begin{enumerate}[(a)]
  \item $\sigma\le\lambda_1$;
  \item $\chi\notin\mathcal F_1$ and $\sigma\le l_2$, where $l_2\le\lambda_2$ is the leaf's
    global second-family lower bound.
\end{enumerate}
Equivalently, this disc contains no zero to the right of the line
$\Re s=1-\sigma/\LL$.
\end{lemma}

\begin{proof}
Such a zero has $|\Im\rho|\le2l+\delta<10l$. By \cref{inp:rectangle} it lies in $R(l)$ or has
$\lambda_\rho>\frac13\log\log\LL$, which exceeds every fixed $\sigma$. In $R(l)$ its parameter is at least
$\lambda_1$, and at least $\lambda_2\ge l_2$ if $\chi\notin \mathcal F_1$ (Z3).
\end{proof}

\begin{lemma}[Separation of zeros beyond the representative strip]\label{lem:strip}
Let $\chi$ be an ordinary character with $T^*$-representative $\rho_\chi$, of height $\gamma_\chi$ and
parameter $\lambda_\chi$, and let $\sigma\le\min(\lambda_\chi,s_{\max})$. Then every zero $\rho$ of $L(s,\chi)$ with
$|1+i\gamma_\chi-\rho|\le\delta_0$ and $\lambda_\rho<\sigma$ satisfies $|\gamma_\chi\LL-\mu_\rho|\ge M$, and there are at
most $K_0$ such zeros, counted with multiplicity.
\end{lemma}

\begin{proof}
Such a zero has $|\Im\rho|\le T^*+\delta_0<2\le10l$ and $\lambda_\rho<\sigma\le s_{\max}<\frac13\log\log\LL$, so it
lies in $R(10l)$ and hence, by \cref{inp:rectangle}, in $R(l)$. As $\lambda_\rho<\sigma\le\lambda_\chi$, the minimality
of the representative gives $|\Im\rho|>T^*$; so the zero is counted in $K_0$. By \cref{lem:Tstar},
$|\Im\rho|\ge T^*+M/\LL$, so $\LL|\Im\rho-\gamma_\chi|\ge\LL(|\Im\rho|-T^*)\ge M$.
\end{proof}

\subsection{Entries}\label{subsec:entries}
We now list the entries that occur in the rows, with their kept zeros. In each case the
hypotheses of \cref{lem:response} hold with the stated kept set $\cS$ and anchor $\sigma$; we write
$r$ for the resulting lower bound for the response, before normalization, so that the entry's
feature is $r/\sqrt{I_uD_u}-\eta'$. \Cref{tab:entries} summarizes the entries. Recall that $F(x)$ is
positive and decreasing on $\RR$, and that a kept zero with $\lambda_\rho\ge\sigma$ has
$m_\rho\Re F(\cdot)\ge\Re F(\cdot)\ge0$.

\begin{table}[ht]
\centering
\footnotesize
\setlength{\tabcolsep}{3.5pt}
\renewcommand{\arraystretch}{1.3}
\begin{tabular}{@{}>{\raggedright\arraybackslash}p{3.45cm}>{\centering\arraybackslash}p{1.35cm}
  >{\centering\arraybackslash}p{1.75cm}>{\centering\arraybackslash}p{3.2cm}
  >{\centering\arraybackslash}p{2.35cm}>{\raggedright\arraybackslash}p{1.65cm}@{}}
\toprule
Entry & Anchor $\sigma$ & Kept set $\cS$ & Response bound $r$ & Diagonal numerator & Rows\\
\midrule
(a) ordinary $\chi$, bin $[\lo,\hi)$, at its $T^*$-representative & $\min(\lo,s)$ & $\{\rho_\chi\}$ &
  $F(\hi-\sigma)-\frac16f(0)$ & $D(s-\sigma)-d_1$ & all three\\
(b) inside first family: $\chi_1$ at $\gamma_1$, and $\bchi_1$ at $-\gamma_1$ & $\min(a,s)$ &
  $\{\rho_1\}$, resp.\ $\{\overline{\rho_1}\}$ & $F(b-\sigma)-\frac16f(0)$ & $D(s-\sigma)-d_1$ & family, graded\\
(c) type $\mathrm{rc}$: $\chi_1$ at $\pm\gamma_1$ & $s_1$ & $\{\rho_1,\overline{\rho_1}\}$ &
  $F(b-s_1)+R_{\lo}-\frac16f(0)$ & $D(0)-d_1+(c^{\hi}-d_1)_+$ & family\\
(d) second zero, type $\mathrm{complex}$: $\chi_1$ at $\gamma_1,\gamma'$, $\bchi_1$ at $-\gamma_1,-\gamma'$ &
  $s_1$ & $\{\rho_1,\rho'\}$ and conjugates & $F(b-s_1)+R_p-\frac16f(0)$ at $\gamma_1$;
  $F(\hi'-s_1)+R_q-\frac16f(0)$ at $\gamma'$ & $D(0)-d_1+(c^{\hi}-d_1)_+$ & family\\
(e) shifted first family: $\chi_1$ at $\gamma_1$, and $\bchi_1$ at $-\gamma_1$ & $s$ &
  $\{\rho_1\}$; $\{\rho_1,\overline{\rho_1}\}$ for $\mathrm{rc}$ & $F(b-s)-\1_{\mathrm{rc}}C_Z-\frac16f(0)$ &
  $D(0)-d_1$ & shifted\\
(f) reserved family, at height-one representatives & $\sigma_2$ & the representative &
  $F(\hi_2-\sigma_2)-\frac16f(0)$ & $D(s-\sigma_2)-d_1$ & all three\\
(g) outside first family: $\chi_1$ at $\gamma_1$ (and $\bchi_1$), and one hidden entry per character &
  $s_1$ & $\{\rho_1\}$; $\{\rho_h\}$ & $F(b-s_1)-\frac16f(0)$; $F(\hi_h-s_1)-\frac16f(0)$ & $D(0)-d_1$ &
  family\\
\bottomrule
\end{tabular}
\caption{The entries of the near rows, justified in the list below. The anchor $\sigma$ enters
through the offset $s-\sigma$, where $s$ is the reference anchor of the row; $\sigma_2$ and the
constants $R_{\lo}$, $R_p$, $R_q$, $c^{\hi}$ and $C_Z$ are defined in the list. ``All three'' means
the family, shifted and graded rows of \cref{subsec:rowkinds}.}
\label{tab:entries}
\end{table}

\begin{enumerate}[(a)]
  \item \emph{Ordinary character in the bin $[\lo,\hi)$.} Entry $(\chi,\gamma_\chi,s-\sigma)$ at its
    $T^*$-representative, with $\sigma=\min(\lo,s)$; $\cS=\{\rho_\chi\}$. If $\sigma\le l_2$ there is no
    zero to the right (\cref{lem:noright}); otherwise \cref{lem:strip} gives the hypotheses. Then
    $r=F(\hi-\sigma)-\frac16f(0)$, and the diagonal numerator is $D(s-\sigma)-d_1$.
  \item \emph{Inside first family, one entry per character.} Entries $(\chi_1,\gamma_1,s-\sigma)$ and, for
    type $\mathrm{complex}$, $(\bchi_1,-\gamma_1,s-\sigma)$, with $\sigma=\min(a,s)\le\lambda_1$; $\cS=\{\rho_1\}$
    (resp.\ $\{\overline{\rho_1}\}$). No zero to the right. $r=F(b-\sigma)-\frac16f(0)$, and the diagonal
    numerator is $D(s-\sigma)-d_1$.
  \item \emph{Conjugate pair, type $\mathrm{rc}$} (rows with $s=s_1$ and $\Xi=0$). Entries
    $(\chi_1,\pm\gamma_1,0)$, a pair of entries of the same character with height difference
    $2\gamma_1$; $\cS=\{\rho_1,\overline{\rho_1}\}$ for both. No zero to the right, and
    \[
      r=F(b-s_1)+R_{\lo}-\tfrac16f(0),\qquad R_{\lo}\le\inf\Re F(x+2i\mu)
    \]
    over $x\in[a-s_1,b-s_1]$ and $\mu$ in the height range of the leaf ($[0,1]$ or $[1,\infty)$). Pair
    excess $e\le(c^{\hi}-d_1)_+$ with $c^{\hi}\ge\sup\Re G(-s_1+2i\mu)$ over the range, so the diagonal numerator
    is $D(0)-d_1+(c^{\hi}-d_1)_+$.
  \item \emph{Second zero of the first family} (rows with $s=s_1$ and $\Xi=0$; type $\mathrm{complex}$
    with a finite gap $[p^*,\hi']$). Let $\rho'$ be an admissible zero of $\chi_1$ realizing $\lambda'$
    (\cref{sec:notation}), of height $\gamma'$, namely the one that witnesses the piece of the second-zero split
    (\cref{subsec:subdivision}) in which the configuration lies, and $y=(\gamma_1-\gamma')\LL$, restricted to that
    piece: $|y|\in[0,1],[1,\frac32],[\frac32,2],[2,3],[3,5]$ or $[5,\infty)$.
    Entries $(\chi_1,\gamma_1,0)$, $(\chi_1,\gamma',0)$, $(\bchi_1,-\gamma_1,0)$, $(\bchi_1,-\gamma',0)$, with kept sets
    $\{\rho_1,\rho'\}$, $\{\rho',\rho_1\}$ and their conjugates (a zero is kept only if it lies in the
    disc of the entry). No zero to the right. Then
    \[
      r_{\gamma_1}=F(b-s_1)+R_p-\tfrac16f(0),\qquad r_{\gamma'}=F(\hi'-s_1)+R_q-\tfrac16f(0),
    \]
    where $R_p\le\inf\Re F(x+iy)$ over $x\in[p^*-s_1,\hi'-s_1]$ and $R_q\le\inf\Re F(x+iy)$ over
    $x\in[a-s_1,b-s_1]$, with $|y|$ in the piece. Here $x\ge0$ (as $s_1\le a\le p^*$), so $\Re F(x+iy)\ge0$ by
    Condition~2 and we take $R_p,R_q\ge0$; the multiplicities of the kept zeros can then only increase
    the sums. On the last piece $R_p=R_q=0$, which also covers a zero outside the disc (its term is
    then absent). Each same-character pair is a controlled pair with excess
    $(c^{\hi}-d_1)_+$, $c^{\hi}\ge\sup\Re G(-s_1+iy)$, so every diagonal numerator is $D(0)-d_1+(c^{\hi}-d_1)_+$;
    an entry of $\chi_1$ and an entry of $\bchi_1$ form a pair of distinct characters, with the nonprincipal
    quotient $\chi_1^2$, and carry no mutual excess. If $\rho'=\rho_1$ is a double zero, the two
    entries coincide and $r\ge m_{\rho_1}F(\lambda_1-s_1)-\frac16f(0)$ with $m_{\rho_1}\ge2$, which is at least
    both $r_{\gamma_1}$ and $r_{\gamma'}$ (here $y=0$ and $\lambda'=\lambda_1$).
  \item \emph{Shifted first family} (the shifted row, $s=\min(1.9,p^*,l_2)$, used only if $s>s_1$ and
    $s\ge a$). Entries $(\chi_1,\gamma_1,0)$ and, for type $\mathrm{complex}$, $(\bchi_1,-\gamma_1,0)$, with
    anchor $s$. The only zeros of $\chi_1$ with parameter below $s$ are $\rho_1$ and, for type
    $\mathrm{rc}$, $\overline{\rho_1}$, since every other occurrence has parameter at least $\lambda'\ge p^*\ge s$.
    With $\cS=\{\rho_1\}$ (and $\overline{\rho_1}$ for type $\mathrm{rc}$) nothing remains in the hypotheses of
    \cref{lem:response}. As conjugate zeros of a real character have the same multiplicity $m$,
    \[
      \sum_{\rho\in\cS}m_\rho\Re F(\cdot)=m\bigl(F(\lambda_1-s)+\1_{\mathrm{rc}}\Re F(\lambda_1-s+2i\mu_1)\bigr)
      \ge m\bigl(F(b-s)-\1_{\mathrm{rc}}C_Z\bigr),
    \]
    with $C_Z=\mathcal C_c(s-a)$ from \cref{lem:parabolic}(v) ($0$ if $s\le a$). If $F(b-s)-\1_{\mathrm{rc}}C_Z\ge0$
    this is at least $F(b-s)-\1_{\mathrm{rc}}C_Z$, and we take $r=F(b-s)-\1_{\mathrm{rc}}C_Z-\frac16f(0)$.
    Otherwise the proposed bound $F(b-s)-C_Z-\frac16f(0)$ is negative and the entry receives the feature
    $0$; this is admissible whatever the true response bound $r_j$ is, by the last clause of
    \cref{cor:zeroform} (with $v'_j=0$). The diagonal numerator is $D(0)-d_1$.
  \item \emph{Reserved family.} Entries $(\chi,\gamma_\chi,s-\sigma_2)$ for its $n_2$ characters at their
    height-one representatives, with $\sigma_2=\min(\max(a,\min(l_2,\lo_2)),s)$. Every zero of these
    characters in $R(l)$ has parameter at least $\max(\lambda_1,\lambda_2)\ge\max(a,l_2)\ge\sigma_2$, so there is no
    zero to the right, $r=F(\hi_2-\sigma_2)-\frac16f(0)$, and the diagonal numerator is $D(s-\sigma_2)-d_1$. With
    second-family columns, the column $[\lo_{2,j},\hi_{2,j}]$ uses $\sigma=\min(\max(a,\min(l_2,\lo_{2,j})),s)$,
    $r=F(\hi_{2,j}-\sigma)-\frac16f(0)$ and the diagonal numerator $D(s-\sigma)-d_1$.
  \item \emph{Outside first family} (family row, $s=s_1$, $\Xi=0$). Global entries $(\chi_1,\gamma_1,0)$ with
    $\cS=\{\rho_1\}$ and, for type $\mathrm{complex}$, $(\bchi_1,-\gamma_1,0)$ with $\cS=\{\overline{\rho_1}\}$ (for
    type $\mathrm{rc}$ there is the single global entry $(\chi_1,\gamma_1,0)$); $r=F(b-s_1)-\frac16f(0)$. And a
    hidden entry $(\chi,\gamma_h,0)$ for each first-family character $\chi$ with a zero in $R_P$, at its zero
    $\rho_h$ of least parameter $t_\chi\in[\lo_h,\hi_h)$, with $\cS=\{\rho_h\}$ and $r=F(\hi_h-s_1)-\frac16f(0)$.
    No zero to the right. Entries of the same character have normalized height differences at least
    $C_B-C_P>M$, so their pair excess vanishes, and every diagonal numerator is $D(0)-d_1$.
\end{enumerate}

\subsection{The rows}\label{subsec:rowkinds}
Each leaf program has one or more rows, of the kinds listed in \cref{tab:rowkinds}. In the family
row of an inside leaf the first family is entered by (b), by (c) for type $\mathrm{rc}$ (with the
height split of \cref{subsec:subdivision}), or by (d) in a piece of the second-zero split. In the
shifted row every ordinary anchor is at most $s\le l_2$. In the graded rows all characters are
distinct.

\begin{table}[ht]
\centering
\small
\begin{tabular}{@{}lllll@{}}
\toprule
Row & Reference anchor $s$ & Sieve weights & Entries, inside & Entries, outside\\
\midrule
family & $s_1$ & none ($\Xi=0$) & (a), (f) and (b), (c) or (d) & (a), (f), (g)\\
shifted & $\min(1.9,p^*,l_2)$ & none ($\Xi=0$) & (a), (e), (f) & not used\\
graded & $1.5$, $1.6$, $1.9$, $2.1$ or $2.3$ & present & (a), (b), (f) & (a), (f)\\
two-test & the shift of the leaf & two test functions & \cref{thm:twotest}, on $37$ roots & not used\\
\bottomrule
\end{tabular}
\caption{The kinds of near rows. The shifted row is used only if $s>s_1$ and $s\ge a$.}
\label{tab:rowkinds}
\end{table}

In every family, shifted and graded row the safe-anchor hypothesis holds by (Z1) and
\cref{inp:rectangle}, as explained after \cref{def:safeanchor}, and the characters are pairwise distinct
whenever $\Xi\not\equiv0$. All the
entries have heights at most $l$. The first family is entered once per character in a row with
sieve weights (for type $\mathrm{rc}$ only $\rho_1$ is kept); the pair and second-zero entries occur only
in rows with $\Xi=0$.

\subsection{The two-test row}\label{subsec:twotest}
On $37$ inside roots the graded rows alone do not certify, and the leaf program contains in
addition the following near row. It is not an instance of \cref{thm:graded}: its Gram form is at
the shift $s$ of the leaf, which may exceed $\lambda_1$, and it pays explicit corrections for the
quotient characters $\chi_1,\bchi_1$. Its proof follows Xylouris's proof of his Lemma~5.3
\cite[(5.34)--(5.40), pp.~72--76]{Xylouris2011nullstellen}, with two test functions.

Fix $\gamma_G>\gamma_Z>0$ and let $f=f_{\gamma_Z}$ and $g=f_{\gamma_G}$ as in \eqref{eq:ftest}, so that $g>0$ on the
support of $f$. Let $s\le\min(p^*,l_2)$ be the shift of the leaf, and put
\begin{align*}
  R_G&=G(-s),\quad R_Z=\int_0^{2\gamma_Z}\frac{f(t)^2}{g(t)}e^{st}\,dt,\quad \mathcal N=\sqrt{R_GR_Z},\quad
  d_0=\frac{g(0)}{6R_G},\\
  C_G&=\sup_{\Re z\ge a-s}(-\Re G(z)),\quad C_F=\sup_{\Re z\ge a-s}(-\Re F(z)),\quad c_G=\frac{C_G}{R_G},\quad k_\eta=1+\eta .
\end{align*}
The function $f^2/g$, extended by $0$, satisfies Condition~1, since it vanishes to order $6$ at
$2\gamma_Z$. Define the features
\begin{gather*}
  v(\lambda)=\frac{F(\lambda-s)-\frac16f(0)}{\mathcal N}-\eta,\qquad
  v_2=\frac{F(\hi_2-s)-\frac{\phi_2}2f(0)}{\mathcal N}-\eta,\\
  v_f=\frac{F(b-s)-\frac{\phi_1}2f(0)-\1_{\mathrm{rc}}C_F}{\mathcal N}-\eta,
\end{gather*}
with $\phi_1=\frac14$ for a real $\chi_1$ and $\phi_1=\frac13$ otherwise, $\phi_2=\frac14$ if $n_2=1$ and
$\phi_2=\frac13$ otherwise, and $\1_{\mathrm{rc}}$ the indicator of type $\mathrm{rc}$; the correlation term $d=k_\eta(d_0+\eta)$; and the diagonals
\[
  \begin{array}{lll}
    \text{type} & D_o/k_\eta & D_f/k_\eta\\
    \mathrm{rr} & 1-d_0+c_G & 1-d_0\\
    \mathrm{rc} & 1-d_0+2c_G & 1-d_0\\
    \mathrm{complex} & 1-d_0+2c_G & 1-d_0+c_G
  \end{array}
\]

\begin{theorem}[Two-test inequality]\label{thm:twotest}
Use the two-test data defined above, with $\gamma_G>\gamma_Z>0$ and
$s\le\min(p^*,l_2)$. For every sufficiently large $q$ and every configuration of an
inside leaf, form the following entries:
\begin{enumerate}[(i)]
  \item any finite set of distinct ordinary characters, each at a zero
    $\rho_\chi\in R(l)$ with $|\Im\rho_\chi|\le1$;
  \item the first character at $\rho_1$, and also $\bchi_1$ at $\overline{\rho_1}$ when
    $\chi_1$ is nonreal;
  \item the $n_2$ characters of the reserved family at their height-one representatives,
    if a family is reserved.
\end{enumerate}
For type $\mathrm{rc}$, the single first-family entry retains both $\rho_1$ and
$\overline{\rho_1}$. Assign features $v(\lambda_{\rho_\chi})$, $v_f$, and $v_2$ to
groups (i)--(iii), and diagonals $D_o$, $D_f$, and $D_o$, respectively. Then
for every choice of real coefficients $a_j\ge0$,
\[
  \Bigl(\sum_ja_jv_j\Bigr)_+^2\le\sum_jD_ja_j^2+d\Bigl(\sum_ja_j\Bigr)^2 .
\]
\end{theorem}

\begin{proof}
Put $\sigma_s=1-s/\LL$ and, in the Hilbert space of sequences indexed by $n$,
\[
  A_j(n)=\sqrt{\Lambda(n)n^{-\sigma_s}g(t_n)}\,\chi_j(n)n^{-i\gamma_j},\qquad
  B(n)=\chi_0(n)\sqrt{\Lambda(n)n^{-\sigma_s}}\,\frac{f(t_n)}{\sqrt{g(t_n)}} .
\]
Then $-\Re\langle A_j,B\rangle=\LL R_j$, where $R_j$ is the response of the entry at the anchor $s$, and
$\LL\sum_ja_jR_j\le\|\sum_ja_jA_j\|\,\|B\|$.

\emph{Responses.} We use \cref{lem:response} (or directly \cref{inp:X32}). For an ordinary or reserved
character every zero in the disc has parameter at least $s$ (\cref{lem:noright}(b), as $s\le l_2$), so
$R_j\ge F(\lambda_{\rho_\chi}-s)-\frac{\varphi}2f(0)-\eps$. For the first family, the zeros of $\chi_1$ with
parameter below $s$ are $\rho_1$ and, for type $\mathrm{rc}$, $\overline{\rho_1}$ (the others have parameter at
least $\lambda'\ge p^*\ge s$); keeping them gives $F(\lambda_1-s)\ge F(b-s)$ and, for type $\mathrm{rc}$,
$\Re F(\lambda_1-s+2i\mu_1)\ge-C_F$ (\cref{lem:parabolic}(v)). Conjugate zeros of a real character have the
same multiplicity $m\ge1$, so the kept terms total at least $m(F(b-s)-\1_{\mathrm{rc}}C_F)$, which is at least
$F(b-s)-\1_{\mathrm{rc}}C_F$ when this is nonnegative (otherwise $v_f<0$, see below). Since
$\varphi(\chi)=\frac14$ for real $\chi$ and $\varphi(\chi)\le\frac13$ always (\cref{inp:X32}), the terms $\frac{\varphi}2f(0)$ are at
most $\frac16f(0)$, $\frac{\phi_1}2f(0)$ and $\frac{\phi_2}2f(0)$ respectively. Thus
$R_j\ge \mathcal N(v_j+\eta)-\eps$ for entries with positive feature. Entries with nonpositive
feature may be omitted in estimating the positive part on the left; we justify
restoring their coefficients to the right-hand side below.

\emph{Norms.} By \cref{inp:X31} applied to $f^2/g$, $\|B\|^2=\LL R_Z(1+o(1))$, and $\|A_j\|^2=\LL R_G(1+o(1))$.

\emph{Off-diagonal terms.} For $j\ne k$ the characters are distinct and $\langle A_j,A_k\rangle$ is the
$g$-weighted sum for the nonprincipal character $\chi_{jk}=\chi_j\bchi_k$ at $\sigma_s+i(\gamma_j-\gamma_k)$. By \cref{inp:X32}
its real part is at most $\LL(d_0R_G+\eps)$ plus $\LL$ times the sum of $-\Re G(\cdot)$ over the zeros of $\chi_{jk}$ in
the disc with parameter below $s$. If $\chi_{jk}\notin\{\chi_1,\bchi_1\}$ there are none, by \cref{lem:noright}(b),
which applies because $|\gamma_j-\gamma_k|\le2\le2l$.
If $\chi_{jk}=\chi_1$ (or $\bchi_1$), they are $\rho_1$ (resp.\ $\overline{\rho_1}$), and for type $\mathrm{rc}$ both, each
contributing at most $C_G$; these zeros are simple, since a multiple $\rho_1$ would give $\lambda_1=\lambda'\ge s$. For a given $j$, the $k$ with $\chi_j\bchi_k\in\{\chi_1,\bchi_1\}$ are
$\chi_k\in\{\chi_j\bchi_1,\chi_j\chi_1\}$: at most two, and at most one if $\chi_1$ is real; a first-family
character has none if $\chi_1$ is real ($\chi_1\bchi_1=\chi_0$) and at most one ($\chi_1^2$) otherwise. With
$2a_ja_k\le a_j^2+a_k^2$ for each such pair, these corrections give the diagonals of the table.

\emph{Assembly.} For $a\ge0$,
$\|\sum_ja_jA_j\|^2\le\LL R_G[\sum_ja_j^2(1-d_0+\mathrm{exc}_j)+(d_0+\eps)(\sum_ja_j)^2]$, and
$\|B\|^2\le(1+\eps)\LL R_Z$. Dividing by $\LL^2\mathcal N^2$ and absorbing the errors in $\eta$ gives the claim
for the positive-feature entries. To restore any omitted entries, expand the
right-hand side as
$\sum_j(D_j+d)a_j^2+2d\sum_{j<k}a_ja_k$. Its coefficients are nonnegative:
the table gives $D_j+d\ge k_\eta(1+\eta)>0$ and $d>0$.
Restoring nonnegative $a_j$ therefore increases this bound, while a nonpositive feature
can only decrease the positive part on the left.
\end{proof}

\begin{proposition}[Detector mixture]\label{prop:mixture}
\Cref{thm:twotest} remains true when $f$ is replaced by $f_\epsilon=f+\epsilon g$ with $\epsilon\ge0$, provided
$F$, $f(0)$ and $C_F$ are replaced by $F+\epsilon G$, $f(0)+\epsilon g(0)$ and $C_F+\epsilon C_G$, and $\mathcal N$ by
$\mathcal N_\epsilon=\sqrt{R_G(R_Z+2\epsilon F(-s)+\epsilon^2R_G)}$; the Gram quantities are unchanged.
\end{proposition}

\begin{proof}
$f_\epsilon$ satisfies Conditions~1 and~2, and $f_\epsilon^2/g=f^2/g+2\epsilon f+\epsilon^2g$ is a sum of functions
satisfying Condition~1, so $\|B_\epsilon\|^2=\LL(R_Z+2\epsilon F(-s)+\epsilon^2R_G)(1+o(1))$. The rest of the proof
is unchanged, using $-\Re(F+\epsilon G)\le C_F+\epsilon C_G$ on $\Re z\ge a-s$.
\end{proof}

\begin{proposition}[Paired first family]\label{prop:paired}
Use the data and inside-leaf hypotheses of \cref{thm:twotest}. Suppose $\chi_1$ is
nonreal and the specification $\mathfrak s$ has a finite gap
$\lambda'\in[p^*,h]$, where $h=\hi'$. Fix $\zeta>0$ and
a number $\mathcal B$ with
\begin{equation}\label{eq:pairB}
  \frac{\Re G(-s+i\Delta)}{R_G}-\zeta\,\frac{\Re F(\lambda_1-s+i\Delta)+\Re F(\lambda'-s+i\Delta)}{2\mathcal N}\le\mathcal B
\end{equation}
for all real $\Delta$, $\lambda_1\in[a,b]$ and $\lambda'\in[p^*,h]$. Let $e\ge C_F/(2\mathcal N)$, and put
\[
  m=\frac{F(b-s)+F(h-s)}{2\mathcal N}-\frac{f(0)}{6\mathcal N}-\eta-e,\qquad
  D_*=k_\eta\Bigl\{\frac{1+\mathcal B}2-d_0+c_G-\frac{\zeta e}2\Bigr\},\qquad\zeta_*=\frac{k_\eta\zeta}2 .
\]
Let $m'\le m$ and $D'\ge D_*$ with $D'>0$ and $2D'\ge\zeta_*\max(m',0)$. Then, for every configuration
of the leaf, there is $\tau\in[0,\sqrt d]$ satisfying \eqref{eq:T} for the entries of \cref{thm:twotest},
with the first-family feature $m'$ and diagonal $D'$ in place of $v_f$ and $D_f$.
\end{proposition}

\begin{proof}
For $\chi_1$ use the averaged vector $A_f=\frac12(A_{\chi_1,\gamma_1}+A_{\chi_1,\gamma'})$, with an admissible $\rho'$ of height
$\gamma'$ (\cref{sec:notation}; so $\rho'$ is a zero of $\chi_1$ and $|\gamma'|\le l$), and its conjugate for $\bchi_1$. Its
response is the average of the two responses, which we bound by \eqref{eq:X34} at the points
$1-s/\LL+i\gamma_1$ and $1-s/\LL+i\gamma'$; they lie in $R(9l)$, and \eqref{eq:X34} does not require the kept zeros
to lie in a disc. The set $A_1$ of zeros of $\chi_1$ in $R(l)$ with parameter below $s$ is contained in
$\{\rho_1\}$, since every occurrence in $R(l)$ other than the distinguished one has parameter at least
$\lambda'\ge p^*\ge s$. We put $\rho'$ into $A_2$, and also $\rho_1$ if $\lambda_1\ge s$, with their multiplicities (if
$\rho'=\rho_1$ this is $\rho_1$ with multiplicity at least $2$); by \cref{inp:logfree} their number is bounded
independently of $q$. The multiplicities can only increase the response: the terms of $\rho'$ are
$\ge0$ by Condition~2 because $\lambda'\ge s$, the term $F(\lambda_1-s)$ is positive, and if $\rho_1$ is a multiple
zero then $\lambda_1=\lambda'\ge s$, so its terms are $\ge0$ as well. With $\Delta=\LL(\gamma_1-\gamma')$, $P=\Re G(-s+i\Delta)/R_G$ and $Q$ the second
fraction in \eqref{eq:pairB}, the feature of $A_f$ is therefore at least
$\frac{F(b-s)+F(h-s)}{2\mathcal N}-\frac{f(0)}{6\mathcal N}+Q-\eta$, and $\|A_f\|^2=\LL R_G\frac{1+P}2(1+o(1))$ by \cref{inp:X31}. Each
off-diagonal term $\langle A_f,A_k\rangle$ is the average of two terms of the kind estimated in the proof of
\cref{thm:twotest}, at the heights $\gamma_1-\gamma_k$ and $\gamma'-\gamma_k$, whose absolute values are at most
$l+1\le2l$; the term between $A_f$ and its conjugate is an average of four terms, at heights
$2\gamma_1$, $\gamma_1+\gamma'$ and $2\gamma'$, which are also at most $2l$ in absolute value. So \cref{lem:noright}
applies to all of them, averaging creates no new quotient characters, and the off-diagonal estimates
are unchanged (for a cubic $\chi_1$ the quotient $\chi_1^2=\bchi_1$ of the conjugate pair is paid by the $c_G$
in $D_*$). Put $U=Q+e$; then $U\ge0$, since $\lambda'\ge s$ gives
$Q\ge-C_F/(2\mathcal N)$.

By \eqref{eq:pairB}, $P\le\mathcal B+\zeta Q$, so the actual first-family feature is at least $m+U$ and the
actual diagonal $k_\eta\{\frac{1+P}2-d_0+c_G\}$ is at most $D_*+\zeta_*U$. Put $U'=U+(m-m')\ge U\ge0$. Then the actual
feature is at least $m'+U'$, and the actual diagonal is at most $D'+\zeta_*U'$, which is positive.
Replacing the first-family features and diagonals by these bounds preserves the quadratic
inequality of \cref{thm:twotest}, and \cref{prop:threshold} gives $\tau\in[0,\sqrt d]$ satisfying \eqref{eq:T}
with the first-family data $(m'+U',D'+\zeta_*U')$. Finally, for $\tau\ge0$ and $U'\ge0$,
\[
  \frac{(m'+U'-\tau)_+^2}{D'+\zeta_*U'}\ge\frac{(m'-\tau)_+^2}{D'} .
\]
Indeed the right side is $0$ if $m'\le\tau$. Otherwise put $z=m'-\tau\in(0,m']$. The derivative of
$(z+u)^2/(D'+\zeta_*u)$ in $u\ge0$ is $(z+u)(2D'-\zeta_*z+\zeta_*u)/(D'+\zeta_*u)^2\ge0$, because
$2D'\ge\zeta_*m'\ge\zeta_*z$. So \eqref{eq:T} holds at the same $\tau$ with $(m',D')$.
\end{proof}

In the leaf program the two-test row is the threshold form \eqref{eq:T} of \cref{thm:twotest} (or of one of
its two variants, \cref{prop:mixture,prop:paired}):
\[
  \sum_ix_i\frac{(V_i-\tau)_+^2}{D_o}+n\frac{(v_f-\tau)_+^2}{D_f}+n_2\frac{(v_2-\tau)_+^2}{D_o}\le1-\frac{\tau^2}d,
\]
with $V_i=v(\hi_i)$ the feature of the right end of the ordinary bin $i$ (the tail and the
reserved columns have feature $0$ in this row), and negative features replaced by $0$. The
$T^*$-representatives are admissible entries, since \cref{thm:twotest} allows any zero of height at
most $1$, and all of them have parameter at least $l_2\ge s$ (Z3). The reserved family appears as a
fixed term with the feature $v_2$ at $\hi_2$, which bounds its feature over the whole reserved range,
and with the diagonal $D_o$. The row is used on $37$ roots, once on each: with a single test on $31$
of them, with a mixture on $2$, and in the paired form on $4$ (all with a nonreal $\chi_1$, as
\cref{prop:paired} requires). In the paired form the certified pair $(m',D')$ consists of a lower
bound for $m$ and an upper bound for $D_*$, and the condition $2D'\ge\zeta_*\max(m',0)$ is checked for
this pair; the margins $2D'-\zeta_*m'$ are between $0.18$ and $0.30$. The condition \eqref{eq:pairB} is
verified with interval arithmetic. Both of its terms are even in $\Delta$, so $\Delta\ge0$ suffices. On the grid
$\Delta\in\frac1{200}\ZZ\cap[0,30]$ the left side is enclosed at the four corners of the parameter box
$[a,b]\times[p^*,h]$. As the left side is a sum of a function of $\lambda_1$ and a function of $\lambda'$, its
maximum over the box exceeds the maximum over the corners by at most
$\frac\zeta{2\mathcal N}\bigl(\frac{(b-a)^2}8\int t^2f(t)e^{(s-a)_+t}\,dt+\frac{(h-p^*)^2}8\int t^2f(t)\,dt\bigr)$, the linear
interpolation error in each variable, and this is added. Between grid points a
second-derivative bound in $\Delta$ is used, and for $\Delta\ge30$ a closed-form bound from
\cref{lem:parabolic}(iv). The value $\mathcal B$ used is at least $0$.

\subsection{Validity of the rows}

\begin{proposition}[Every configuration satisfies the near rows]\label{prop:rowsvalid}
Let $\mathfrak s$ specify a leaf or one of its subcases. For sufficiently large $q$,
take any configuration in $C(\mathfrak s)$ and assign its characters to columns as in
\cref{prop:realization}, obtaining column values $x$. Then, for every near row $k$,
there is a threshold $\tau_k$ such that
\[
  0\le\tau_k\le\sqrt{d_k/S},\qquad
  \Phi_k(x,\tau_k)\le1-\frac{\tau_k^2}{d_k/S}.
\]
Here $d_k$ is the stored correlation term, $S=10^{16}$, and $\Phi_k$ is the row expression of
\cref{def:leafprogram}. Each row may have its own threshold.
\end{proposition}

\begin{proof}
For a graded, family or shifted row, take as entries the characters of the configuration that
are placed in the columns and family terms of the row: the ordinary characters of $O$ at their
$T^*$-representatives, the reserved characters, and the first family, with the entries of
\cref{subsec:entries}. These form an admissible family of entries for \cref{thm:graded}, with the
kept sets verified in \cref{subsec:entries}, so \cref{cor:zeroform} gives a threshold $\tau$ at
which \eqref{eq:T} holds with the normalized features, diagonals and correlation term. Each column's feature is
at most the feature of every character placed in it (features are taken at the right end of the
column, and $F$ is decreasing), or is $0$; each column's diagonal is at least theirs, or the column
has feature $0$ (an entry whose diagonal numerator is not safely positive is omitted, and its column
is given feature $0$), in which case it contributes nothing for $\tau\ge0$; each family
term $(n,v,D)$ receives $n$ entries with feature at least $v$ and diagonal at most $D$; and the
characters in the tail column and in the hidden tail are simply omitted from the row. Since each term $(v-\tau)_+^2/D$
increases with $v$ and decreases with $D$ for $\tau\ge0$, the row of the leaf holds at $\tau$. For the
two-test row the same argument applies to \cref{thm:twotest} and its variants, through
\cref{prop:threshold}.
\end{proof}

The numbers $I_u$, $D_u$, $D(\delta)$, the transforms and the special terms $R_{\lo}$, $c^{\hi}$, $R_p$, $R_q$,
$C_Z$ and $C_F$ are computed from the rational parameters of each row with outward interval arithmetic. The
transforms are enclosed with the closed form of \cref{lem:parabolic}(iv) or with its power series near
$0$. The infima and suprema over heights in $R_{\lo}$, $c^{\hi}$, $R_p$ and $R_q$ are enclosed on grids of step
at most $\frac1{50}$, with second-derivative bounds $\int t^2\hat f_c(t)e^{dt}\,dt$ between grid points and
Lipschitz bounds in the real parts, and by the tail bounds of \cref{lem:parabolic}(iv) beyond the
grid. The constants $C_Z$, $C_F$ and $C_G$ are enclosed on the line $\Re z=-d$ (by
\cref{lem:parabolic}(v)), for heights in $[0,20]$ by adaptive linear interpolation, starting from
$200$ cells of width $\frac1{10}$ and bisecting cells until the largest cell bound is within a small
fixed tolerance of the best value found, with the same second-derivative bound, and beyond height
$20$ by a closed-form tail bound. The features, diagonals and correlation terms of all rows other than the $37$ two-test rows have been
recomputed from their rational parameters by an independent program whose soundness is formally
verified (\cref{sec:verification}).

% !TEX root = ../linnik399.tex
\section{Leaf programs and exact certificates}\label{sec:lp}

For each case of the zero analysis, we construct a family of linear programs indexed by
the thresholds in the near rows. Their common objective bounds the zero sum. We cover
the threshold domain by finitely many boxes and give a linear relaxation on each box.
An exact certificate then bounds every feasible point, without having to find an
optimal one.

The argument uses weak linear-programming duality
\cite{Chvatal1983linear,Schrijver1986theory} with outward rounding of the data; compare
\cite{Jansson2004rigorous,NeumaierShcherbina2004safe,ApplegateCookDashEspinoza2007exact}.
This section is independent of the analytic number theory. Its conclusion is the
certificate soundness theorem, \cref{thm:certificate}, which is formally verified
(\cref{subsec:lean}).

\subsection{Leaf data}\label{subsec:leafdata}
The coefficients, budgets, and objective constants of a leaf are stored as integers
at scale
\[
  S=10^{16}.
\]
Thus a stored coefficient $v$ represents $v/S$. Character counts such as $n$, $n_g$,
and $n_E$ are ordinary unscaled integers. The variables $x_c$ and thresholds $\tau_k$
are also unscaled real numbers.

A leaf has a finite set of \emph{columns} $c$ and a finite list of \emph{near rows}
$k=1,\dots,K$. A column $c$ carries
\begin{itemize}
  \item an objective coefficient $G_c$ and a far weight $W_c$;
  \item a count coefficient $C_c$, a hidden-count coefficient $N_c$ and a second-family indicator
    $E_c$ (stored at scale $S$: $C_c,E_c\in\{0,S\}$, and $N_c$ is $0$, $S$ or, for the hidden tail,
    $\lfloor S/w(\mathrm{end})\rfloor$, with $\mathrm{end}$ as in \cref{subsec:leafmodel});
  \item for each row $k$, a feature $v_{ck}$ and a diagonal $D_{ck}>0$.
\end{itemize}
A row $k$ carries a correlation term $d_k>0$ and finitely many \emph{family terms} $(n,v,D)$, where
$n\ge0$ is a count, $v$ a feature and $D>0$ a diagonal. The leaf also has a far budget $F$, counts $n_g$ (the number $n$ of first-family characters, in
outside leaves) and $n_E$ (the number of reserved characters carried by second-family columns: the
size of the reserved family if it is charged through columns, and $0$ otherwise), and two
constants: $\mathrm{first}$, which bounds the cost of the first family of an inside leaf ($0$ for an
outside leaf) plus any fixed charge for the reserved family, and $\mathrm{final}$, the error
allowance, with $\mathrm{final}/S=5\eta$ (\cref{subsec:leafmodel}).

\begin{definition}[Leaf program]\label{def:leafprogram}
For real column values $x=(x_c)$ and real thresholds $\tau=(\tau_k)$, put
\[
  \Phi_k(x,\tau_k)=\sum_c x_c\frac{(v_{ck}/S-\tau_k)_+^2}{D_{ck}/S}
  +\sum_{(n,v,D)\in\mathrm{fam}(k)} n\,\frac{(v/S-\tau_k)_+^2}{D/S}.
\]
The pair $(x,\tau)$ is \emph{feasible} if all of the following hold:
\begin{itemize}
  \item $x_c\ge0$ for every $c$;
  \item (far) $\sum_cW_cx_c\le F$;
  \item (count) $\sum_cC_cx_c\le2S$;
  \item (hidden count) $\sum_cN_cx_c\le n_gS$;
  \item (second family) $\sum_cE_cx_c=n_ES$;
  \item (near rows) for every $k$: $0\le\tau_k$, $\tau_k^2\le d_k/S$ and
    $\Phi_k(x,\tau_k)\le1-\tau_k^2/(d_k/S)$.
\end{itemize}
The \emph{value} of $x$ is
\[
  \mathcal{V}(x)=\frac1S\Bigl(\mathrm{first}+\mathrm{final}+\sum_cG_cx_c\Bigr).
\]
\end{definition}

For fixed $\tau$, these are linear constraints on $x$. Allowing $\tau$ to vary gives
the family of programs to be certified. The variables $x_c$ are real and nonnegative:
integer character counts are feasible values, but the relaxation also permits
fractional counts and the weighted tail masses defined in \cref{subsec:leafmodel}.

In \cref{sec:casetree} we show that for every configuration of zeros the actual characters give a
feasible pair whose value bounds the zero sum $W$ of \cref{sec:detection}. It is therefore enough
to show that every feasible pair has value below $1$.

\subsection{Relaxations of a near row}
Fix a row $k$ and write $d=d_k/S$. The near-row constraint says
$Q_k(\tau)\le1$, where
\[
  Q_k(\tau)=\sum_jc_j(u_j-\tau)_+^2+\frac{\tau^2}{d}
\]
with $c_j\ge0$ collecting the coefficients $x_c/(D_{ck}/S)$ and $n/(D/S)$ and $u_j$ the features.
The function $Q_k$ is convex and continuously differentiable, with
$Q_k'(\tau)=-2\sum_jc_j(u_j-\tau)_++2\tau/d$. In the next two lemmas, $d>0$, all
$c_j\ge0$, and $0\le a\le b$.

\begin{lemma}[First-order relaxation]\label{lem:firstorder}
If $a\le\tau\le b$ and $Q_k(\tau)\le1$, then
\[
  \sum_jc_j(u_j-b)_+^2\le1-\frac{a^2}{d}.
\]
\end{lemma}

\begin{proof}
Each $(u_j-\tau)_+^2\ge(u_j-b)_+^2$ since $\tau\le b$, and $\tau^2\ge a^2$ since $0\le a\le\tau$.
\end{proof}

\begin{lemma}[Tangent relaxation]\label{lem:tangent}
Let $m=(a+b)/2$ and $h=(b-a)/2$. If $a\le\tau\le b$ and $Q_k(\tau)\le1$, then at least one of
the following holds:
\begin{align*}
  \text{(case 0)}\quad&\sum_jc_j\Bigl(\bigl((u_j-m)_++h\bigr)^2-h^2\Bigr)\le
    1-\frac{(m-h)^2-h^2}{d},\\
  \text{(case 1)}\quad&\sum_jc_j\,\tilde\theta_j\le1-\frac{(m+h)^2-h^2}{d},\qquad
    \tilde\theta_j=\begin{cases}\bigl((u_j-m)-h\bigr)^2-h^2,&u_j>m,\\0,&u_j\le m.\end{cases}
\end{align*}
\end{lemma}

\begin{proof}
By convexity, $1\ge Q_k(\tau)\ge Q_k(m)+Q_k'(m)(\tau-m)\ge Q_k(m)-h|Q_k'(m)|$, so
$Q_k(m)-hQ_k'(m)\le1$ or $Q_k(m)+hQ_k'(m)\le1$. Expanding with
$(u-m)_+^2\pm2h(u-m)_+=((u-m)_+\pm h)^2-h^2$ and $m^2\mp2hm=(m\mp h)^2-h^2$ gives the two cases.
In case~1 a term with $u_j\le m$ is $0$.
\end{proof}

The left-hand sides are linear in the $c_j$, hence in the column values $x$. In case~1 the
coefficients $\tilde\theta_j$ can be negative.

\subsection{Certificates}
Thresholds are measured in ticks of $1/T_S$, where $T_S=10^6$, and dual multipliers are scaled
by $D_S=10^{12}$. In this subsection $[a,b]$ denotes an interval of thresholds in ticks, not a
first-zero cell, and $m$, $h$, $p$ and $\beta$ are local to the relaxation formulas.

\begin{definition}[Root threshold box]\label{def:rootbox}
For row $k$, put
\[
  e_k=\left\lceil\sqrt{\lfloor d_kT_S^2/S\rfloor+1}\,\right\rceil.
\]
The \emph{root box} is $\prod_k[0,e_k]$ in integer-tick coordinates. Its corresponding
domain for real thresholds is $\prod_k[0,e_k/T_S]$.
\end{definition}

Since $e_k^2S\ge d_kT_S^2$, every threshold with $\tau_k^2\le d_k/S$ lies in
$[0,e_k/T_S]$.

\begin{definition}[Integer costs and budgets]\label{def:costs}
Let $[a,b]$ be an interval of integer ticks with $0\le a<b$, and let $v$ and $D>0$
be a stored feature and diagonal. The following costs are coefficients of column
variables; the family terms are subtracted from the row budget.
\begin{enumerate}[(a)]
  \item (First order.) The cost is $\lfloor(v-bS/T_S)_+^2/D\rfloor$. The budget of row $k$ is
    \[
      S-\Bigl\lfloor\frac{S^2a^2}{T_S^2d_k}\Bigr\rfloor-\sum_{(n,v,D)}\Bigl\lfloor\frac{n\,(v-bS/T_S)_+^2}{D}\Bigr\rfloor .
    \]
  \item (Tangent, case $\iota\in\{0,1\}$.) Put $m=(a+b)/(2T_S)$, $h=(b-a)/(2T_S)$ and $\bar v=v-mS$
    ($mS$ and $2hS$ are integers, since $2T_S$ divides $S$).
    The cost is $\lfloor S\,\Theta/D\rfloor$, where $\Theta=0$ if $\bar v\le0$, and otherwise
    $\Theta=\lfloor\bar v(\bar v+2hS)/S\rfloor$ for $\iota=0$ and $\Theta=\lfloor\bar v(\bar v-2hS)/S\rfloor$ for $\iota=1$. The
    budget of row $k$ is $\lceil S\beta\rceil$, where $\beta$ is the exact rational number
    \[
      \beta=1-\frac{p^2-h^2}{d_k/S}-\sum_{(n,v,D)}n\,\frac{\Theta_\iota(v/S)}{D/S},
    \]
    with $p=a/T_S$ for $\iota=0$ and $p=b/T_S$ for $\iota=1$, and $\Theta_0(u)=((u-m)_++h)^2-h^2$,
    $\Theta_1(u)=((u-m)-h)^2-h^2$ if $u>m$ and $0$ otherwise.
\end{enumerate}
\end{definition}

The costs are the exact relaxation coefficients, multiplied by $S$ and rounded \emph{down}. The
budgets are at least the exact right-hand sides multiplied by $S$: the tangent budget is rounded up,
and in the first-order budget each subtracted term is rounded down.

\begin{definition}[Certificate]\label{def:certificate}
A certificate for a leaf is a finite binary tree starting from the root box of
\cref{def:rootbox}. An internal node splits one coordinate $k$ of
its box $[a,b]$ at an integer tick $a<\mathrm{mid}<b$, into $[a,\mathrm{mid}]$ and
$[\mathrm{mid},b]$. At a terminal box, each row uses either the first-order relaxation
or the two alternatives of the tangent relaxation. A \emph{relaxation case} chooses one
alternative for each tangent row. All combinations must be certified, because the
alternative supplied by \cref{lem:tangent} can depend on the feasible point.

For consistency with the stored data, the mode is encoded by $o_k=2$ for a first-order
row and $o_k=1$ for a tangent row. The case vector has $\iota_k=2$ when $o_k=2$, and
$\iota_k\in\{0,1\}$ when $o_k=1$. In the stored trees the halves of a split are listed in the
order $[a,\mathrm{mid}]$, $[\mathrm{mid},b]$, the rows are numbered from $0$, and the relaxation
cases of a terminal box are listed in lexicographic order of $(\iota_1,\dots,\iota_K)$.
For every relaxation case the node records either
\begin{enumerate}[(i)]
  \item an \emph{exclusion}: a row $k$ whose costs are all $\ge0$ and whose budget is $<0$; or
  \item \emph{integer duals} $Y_F,Y_C,Y_N,Y_E^+,Y_E^-\ge0$ and $Z_1,\dots,Z_K\ge0$ such that for every
    column $c$
    \begin{equation}\label{eq:dualfeasible}
      D_SG_c\le Y_FW_c+Y_CC_c+Y_NN_c+(Y_E^+-Y_E^-)E_c+\sum_kZ_k\,\mathrm{cost}_{ck}.
    \end{equation}
    The value of the relaxation case is then
    \[
      \Bigl\lceil\frac{Y_FF+2SY_C+n_gSY_N+(Y_E^+-Y_E^-)n_ES+\sum_kZ_k\,\mathrm{budget}_k}{D_S}\Bigr\rceil
      +\mathrm{first}+\mathrm{final}.
    \]
\end{enumerate}
The value of a terminal box is the maximum of $-1$ and the values of its relaxation cases that
carry duals (so it is $-1$ if every relaxation case is excluded), and
the certificate is \emph{accepted} if every terminal box has value $<S$. The data must be
\emph{well formed}: all diagonals and correlation terms are positive, and all family counts are
nonnegative.
\end{definition}

\begin{theorem}[Soundness of the integer certificates]\label{thm:certificate}
Let a leaf have positive diagonals and correlation terms and nonnegative family counts. If it has
an accepted certificate in the sense of \cref{def:certificate}, then every feasible
pair $(x,\tau)$ of \cref{def:leafprogram} satisfies $\mathcal V(x)<1$.
This conclusion holds for all real nonnegative column values, including fractional ones.
\end{theorem}

\begin{proof}
Let $(x,\tau)$ be feasible and put $t_k=\tau_kT_S$. Since $\tau_k^2\le d_k/S$ we have $t_k\le e_k$,
so $t$ lies in the root box. At every internal node $t$ lies in one of the two halves, so there is
a terminal box $\prod_k[a_k,b_k]$ that contains $t$.

For each row, \cref{lem:firstorder} (if $o_k=2$) or \cref{lem:tangent} (if $o_k=1$) gives a case
$\iota_k$ in which the exact relaxed inequality holds at $x$. Take the relaxation case $\iota$ so
obtained. Multiply the relaxed inequality of row $k$ by $S$. Each coefficient of $x_c$ is then at
least $\mathrm{cost}_{ck}$ and the right side is at most $\mathrm{budget}_k$, by the rounding
directions in \cref{def:costs} and since $x_c\ge0$ and $n\ge0$. Therefore
\begin{equation}\label{eq:rowsatisfied}
  \sum_cx_c\,\mathrm{cost}_{ck}\le\mathrm{budget}_k\qquad(1\le k\le K).
\end{equation}
Here we used $S\cdot(v/S-b/T_S)_+^2/(D/S)=(v-bS/T_S)_+^2/D$ and the corresponding identities for
the tangent terms.

If the relaxation case $\iota$ is excluded through row $k$, then the left side of \eqref{eq:rowsatisfied} is $\ge0$
and the right side is $<0$, a contradiction. So $\iota$ carries duals. Multiply
\eqref{eq:dualfeasible} by $x_c\ge0$ and sum over $c$. Then use the far, count and hidden-count
constraints, with multipliers $Y_F,Y_C,Y_N\ge0$, the second-family equality, with the multiplier
$Y_E^+-Y_E^-$ of either sign, and \eqref{eq:rowsatisfied}, with multipliers $Z_k\ge0$. This gives
\[
  D_S\sum_cG_cx_c\le Y_FF+2SY_C+n_gSY_N+(Y_E^+-Y_E^-)n_ES+\sum_kZ_k\,\mathrm{budget}_k .
\]
Hence $\mathrm{first}+\mathrm{final}+\sum_cG_cx_c$ is at most the value of $\iota$, which is at most
the value of the terminal box, which is $<S$.
\end{proof}

\begin{remark}
The formal proof of \cref{thm:certificate} brought out that the family counts $n$ must be
nonnegative. With a negative count, the first-order relaxation can fail.
The leaf builders already ensure $n\ge0$, and the checker requires it.
\end{remark}

% !TEX root = ../linnik399.tex
\section{The case tree}\label{sec:casetree}

This section describes the finite case analysis of the middle range $0.1\le\lambda_1<1.5$ and
proves that it covers every configuration of zeros (\cref{thm:cover}). The analytic facts it
uses are the zero-location results of \cref{sec:location}, the costs of \cref{sec:costs}, the far
budget of \cref{sec:far} and the near rows of \cref{sec:rows}. There are two tasks:
showing that the cases cover every possible configuration, and showing that a
configuration in a case gives a feasible point of its program. These are proved
separately before being combined in \cref{thm:cover}. \Cref{tab:census}, after \cref{subsec:refinement}, lists the levels of the
analysis, with their numbers and the places where they are defined and checked.

\subsection{Configurations}\label{subsec:configurations}
\begin{definition}[Zero configuration]\label{def:configuration}
Fix a sufficiently large $q$. Its \emph{zero configuration} is the multiset $\mathcal Z$ of pairs $(\chi,\rho)$ with
$\chi\ne\chi_0$ a character modulo $q$, $\rho\in R(l)$ and $L(\rho,\chi)=0$, each pair repeated
according to the multiplicity of $\rho$.
\end{definition}

Minima over empty sets are $+\infty$. In the middle range the first zero exists.
The quantities of \cref{sec:notation} are read off from $\mathcal Z$: the first zero $(\chi_1,\rho_1)$,
the first family $\mathcal F_1$ with $n=|\mathcal F_1|$, the parameters $\lambda_2,\lambda_3$ and the family
$\mathcal F_2$ (\cref{def:families}), and the distinguished occurrences, $\lambda'$ and an admissible
$\rho'$ (\cref{def:additionalzero}); for type $\mathrm{rc}$, $\mu_1>0$. In addition we use the
following.
\begin{itemize}
  \item The \emph{height-one parameter} of $\chi$ is $\nu(\chi)=\min\{\lambda_\rho:(\chi,\rho)\in\mathcal
    Z,\ |\gamma_\rho|\le1\}$, and the \emph{$T^*$-parameter} is $\nu_{T^*}(\chi)$, defined in the same
    way with $|\gamma_\rho|\le T^*$ (\cref{def:representative}). Since $L(\rho,\chi)=0$ if and only if
    $L(\overline\rho,\bchi)=0$, both are functions of the family. We put
    $\nu_*=\min\{\nu(\mathcal F):\mathcal F\ne\mathcal F_1\}$.
  \item The configuration is \emph{inside} if $\rho_1\in R_B$ and \emph{outside} otherwise. Since
    $\lambda_1<1.5<C_B$ in the middle range, outside means $|\mu_1|>C_B$. Type $\mathrm{rr}$ is
    always inside.
\end{itemize}
The definitions give $\lambda_1\le\lambda'$, $\lambda_1\le\lambda_2\le\lambda_3$, and
$\nu_{T^*}(\chi)\ge\nu(\chi)\ge\lambda_2$ for $\chi\notin \mathcal F_1$.

\subsection{Specifications}\label{subsec:specifications}

A specification records interval information about a configuration. It may describe
no actual configuration; in that event the corresponding case can be excluded.
Its two lower bounds $l_2$ and $r$ have different scopes: $l_2$ applies throughout
$R(l)$, whereas $r$ applies only to the height-one representatives.

\begin{definition}[Specification]\label{def:specification}
A \emph{specification} is a tuple
\[
  \mathfrak s=(\mathfrak t,[a,b],p,\mathrm{gap},l_2,r,\mathrm{res},\mathrm{loc})
\]
with the following data:
\begin{itemize}
  \item a type $\mathfrak t\in\{\mathrm{rr},\mathrm{rc},\mathrm{complex}\}$ and a
    first-zero cell $[a,b]$, where $a<b$;
  \item a lower bound $p$ for $\lambda'$, and optionally an interval
    $\mathrm{gap}=[\lo',\hi']$ with $p\le\lo'<\hi'\le\infty$, called a \emph{gap}
    (used only for type $\mathrm{complex}$);
  \item a global second-family lower bound $l_2$ and a height-one lower bound $r$,
    with $l_2\le r\le3$;
  \item either no reservation, or a pair $\mathrm{res}=(\hi_2,n_2)$ with
    $r<\hi_2\le2$ and $n_2\in\{1,2\}$; in the latter case put $\lo_2=r$;
  \item a height class $\mathrm{loc}\in\{\mathrm{inside},\mathrm{outside}\}$.
\end{itemize}
The \emph{configuration set} $C(\mathfrak s)$ consists of the configurations for which
some admissible choices of the first zero and the minimizing families satisfy all of
the following:
\begin{enumerate}[(C1)]
  \item the type is $\mathfrak t$ and $\lambda_1\in[a,b]$;
  \item $\lambda'\ge p$, and $\lambda'\in[\lo',\hi']$ if a gap is present (with $\lambda'=\infty$ allowed
    when $\hi'=\infty$);
  \item $\lambda_2\ge l_2$;
  \item if $\mathrm{res}$ is none, then $\nu(\mathcal F)\ge r$ for every family $\mathcal F\ne\mathcal F_1$; if
    $\mathrm{res}=(\hi_2,n_2)$, then $\nu_*\in[r,\hi_2]$ and some family $\mathcal F_{\mathrm{res}}\ne\mathcal F_1$ with
    $\nu(\mathcal F_{\mathrm{res}})=\nu_*$ has exactly $n_2$ characters (so it is a real character if $n_2=1$ and
    a pair of conjugate nonreal characters if $n_2=2$);
  \item the configuration is inside if $\mathrm{loc}=\mathrm{inside}$ and outside if $\mathrm{loc}=\mathrm{outside}$.
\end{enumerate}
In the reserved case, (C4) implies $\nu(\mathcal F)\ge r$ for every $\mathcal F\ne\mathcal F_1$. We call $\mathcal F_{\mathrm{res}}$ the
\emph{reserved family} and its characters \emph{reserved}; the characters outside $\mathcal F_1$ and
$\mathcal F_{\mathrm{res}}$ are \emph{ordinary}.
\end{definition}

In a formula containing $\lo'$, that argument is omitted if no gap is present.
Likewise, reserved-family terms are absent when there is no reservation.

The reserved family is a family with the least height-one parameter among the families other than
the first; it is not assumed to be the family $\mathcal F_2$ that attains $\lambda_2$, which may have its
nearest zero at a height above $1$. The following three facts are used in every leaf.

\begin{lemma}[Consequences of reserving a family]\label{lem:familyfacts}
Let $\mathfrak s$ be a specification and fix a configuration and admissible choices
witnessing its membership in $C(\mathfrak s)$.
\begin{enumerate}[(i)]
  \item \emph{(Count.)} At most two characters $\chi\notin \mathcal F_1$ have $\nu_{T^*}(\chi)<\lambda_3$.
  \item \emph{(Drop.)} If $\mathfrak s$ is reserved, then $\nu_{T^*}(\chi)\ge\nu(\chi)\ge\lambda_3$ for every
    ordinary character $\chi$.
  \item \emph{(Cap.)} If $\mathfrak s$ is reserved, then $\lambda_2\le\nu_*\le\hi_2$.
\end{enumerate}
\end{lemma}

\begin{proof}
(i) Such a character has a zero in $R(l)$ with parameter below $\lambda_3$, so it lies in $\mathcal F_2$, and
$|\mathcal F_2|\le2$. (ii) Let $\mathcal F$ be the family of an ordinary character. If $\mathcal F\ne\mathcal F_2$, every zero of $\mathcal F$
in $R(l)$ has parameter at least $\lambda_3$. If $\mathcal F=\mathcal F_2$, then $\mathcal F_{\mathrm{res}}\notin\{\mathcal F_1,\mathcal F_2\}$, so
$\nu(\mathcal F_{\mathrm{res}})\ge\lambda_3$, and $\nu(\mathcal F)\ge\nu_*=\nu(\mathcal F_{\mathrm{res}})$ by minimality. (iii) The
height-one representative of $\mathcal F_{\mathrm{res}}$ is a zero in $R(l)$ of a character outside $\mathcal F_1$.
\end{proof}

\subsection{The source cover}\label{subsec:sourcecover}
The roots of the case tree are $2768$ specifications, each used with the height class
$\mathrm{inside}$, and the $1685$ of them whose type is not $\mathrm{rr}$ used also with the height
class $\mathrm{outside}$: $4453$ roots in all. They are constructed from the zero-location results of
\cref{sec:location} in four steps.
\begin{enumerate}[(1)]
  \item \emph{Parent rows.} There are $58$ \emph{parent rows} $(\mathfrak t,A_j,B_j,p_j,r_j)$: $33$ of type
    $\mathrm{rr}$ tiling $[0.1,1.5]$, $13$ of type $\mathrm{rc}$ tiling $[0.628,1.5]$, and $12$ of type
    $\mathrm{complex}$ tiling $[0.44,1.5]$. Each is the statement
    \begin{equation}\tag{T2}\label{eq:T2}
      \text{type }\mathfrak t,\ \lambda_1\in[A_j,B_j]\ \Longrightarrow\ \lambda'\ge p_j\ \text{and}\ \lambda_2\ge r_j,
    \end{equation}
    which is proved in \cref{prop:parentrows} from the tables of Heath-Brown and Xylouris.
  \item \emph{Base cells.} Each parent interval is cut into cells of width $0.01$, giving $340$
    \emph{base cells}.
  \item \emph{Cells and gaps.} The base cells are tiled by $478$ first-zero cells in all. A cell
    $[\lo,\hi]$ of the parent $j$ receives numbers $p\le\max(p_j,\lo)$ and $l_2\le\max(r_j,\lo)$; this is
    valid since $\lambda',\lambda_2\ge\lambda_1\ge\lo$. A cell of type
    $\mathrm{complex}$ may be \emph{partitioned} by gaps $[p,m_1],[m_1,m_2],\dots,[m_k,\infty]$ for
    $\lambda'$. The cells that are not partitioned, and the gaps of those that are, are the $684$
    \emph{gap cases}.
  \item \emph{Records.} Each gap case is either \emph{unreserved}, with one specification with
    $r=l_2$ and no reservation, or \emph{reserved} with a chain $l_2=u_0<u_1<\dots<u_m=e\le2$,
    with specifications $(\mathrm{res}=(u_{j+1},n_2),\,r=u_j)$ for $n_2=1,2$ and $0\le j<m$,
    followed by an unreserved \emph{tail} specification with $r=e$.
\end{enumerate}
By type there are $1083$ specifications of type $\mathrm{rr}$ ($195$ unreserved, $444$ reserved
with $n_2=1$ and $444$ with $n_2=2$), $115$ of type $\mathrm{rc}$ (all unreserved and without gap)
and $1570$ of type $\mathrm{complex}$.

\begin{theorem}[Coverage by the root cases]\label{thm:sourcecover}
Every configuration with $0.1\le\lambda_1<1.5$ lies in $C(\mathfrak s)$ for at least one of the $4453$
root specifications $\mathfrak s$.
\end{theorem}

\begin{proof}
Let $\mathfrak t$ be the type. By \cref{prop:firstlower}, $\lambda_1>0.628$ if $\mathfrak t=\mathrm{rc}$ and
$\lambda_1>0.44$ if $\mathfrak t=\mathrm{complex}$, so $\lambda_1$ lies in some cell of type $\mathfrak t$ (the
cells are closed, and a boundary value lies in two cells). By \eqref{eq:T2} for the parent of the
cell, $\lambda'\ge p$ and $\lambda_2\ge l_2$. If the cell is partitioned, some gap contains $\lambda'$.
In the corresponding gap case, if it is unreserved, every family $\mathcal F\ne\mathcal F_1$ has $\nu(\mathcal F)\ge\lambda_2\ge
l_2=r$. If it is reserved with chain $u_0<\dots<u_m$, then $\nu_*\ge\lambda_2\ge u_0$; if $\nu_*<e$
we choose $j$ with $\nu_*\in[u_j,u_{j+1}]$ and a minimizing family, which is a real character or a
nonreal pair, and the configuration lies in the corresponding record; otherwise it lies in the
tail record. Finally the height class is determined by $|\mu_1|$.
\end{proof}

\subsection{Refinement trees}\label{subsec:refinement}
Each root is the root of a finite \emph{refinement tree}. Every internal node carries a
specification (computed from the root, never read from the data) and is of one of the following
kinds. In each case every configuration of the node lies in the configuration set of one of its
children, or the node has no configuration at all.
\begin{itemize}
  \item \emph{First-zero split} at $a<m<b$: children with cells $[a,m]$ and $[m,b]$, all other data
    kept. Exhaustive, since $\lambda_1\le m$ or $\lambda_1\ge m$.
  \item \emph{Second-family split} at $r<m<\hi_2$ (reserved nodes): children $(\hi_2:=m)$ and
    $(r:=m,\ \lo_2:=m)$. Exhaustive, since $\nu_*\le m$ or $\nu_*\ge m$, and in the second case every
    family $\mathcal F\ne\mathcal F_1$ has $\nu(\mathcal F)\ge\nu_*\ge m$.
  \item \emph{Gap split} at $\lo'<m<\hi'$ (type $\mathrm{complex}$): children with gaps $[\lo',m]$ and
    $[m,\hi']$. Exhaustive.
  \item \emph{Location update}: a lower bound $\lambda_2>h$ valid on the whole cell
    (\cref{prop:realrows,prop:complexrows}; \cref{lem:updates}) replaces $l_2$ and $r$ by $\max(l_2,h)$ and
    $\max(r,h)$, and excludes the node if it is reserved with $\hi_2\le h$. Sound, since every family
    $\mathcal F\ne\mathcal F_1$ has $\nu(\mathcal F)\ge\lambda_2$.
  \item \emph{Exclusion} of a reserved interval or a gap by a lower bound for $\lambda_2$ or
    $\lambda'$ (\cref{prop:polyrows,prop:realfirstsecond}), or of a whole node by the positivity
    argument of \cref{prop:positivity}.
\end{itemize}
Before a leaf model is built, all applicable location updates are applied once more to its
specification; this only adds valid implications. None of these nodes depends on $L$. The terminal
nodes are the \emph{leaves}. The inside trees have $3146$ leaves. Of these, $233$ (in $199$ roots) are
stored in a separate file and are called \emph{supplementary leaves}; the reason is historical
(\cref{rem:history}). Every leaf, supplementary or not, receives the same model. The
outside trees have $1642$ leaves. In all there are $4788$ leaves; $343$ inside roots and $70$
outside roots have none, because all their configurations are excluded.

\begin{table}[ht]
\centering
\small
\begin{tabular}{lrr}
\toprule
Node & Inside & Outside\\
\midrule
first-zero / second-family / gap splits & 1 / 408 / 593 & 0 / 27 / 0\\
location updates, real (restrict / exclude) & 225 / 75 & --\\
location updates, complex (restrict / exclude) & 11 / 0 & --\\
exclusions by lower bounds for $\lambda_2$ or $\lambda'$ & 352 & 70\\
exclusions by positivity & 197 & 0\\
leaves & 3146 & 1642\\
\bottomrule
\end{tabular}
\caption{Nodes of the refinement trees.}
\label{tab:nodes}
\end{table}

\begin{table}[tp]
\centering
\footnotesize
\setlength{\tabcolsep}{4pt}
\begin{tabular}{@{}>{\raggedright\arraybackslash}p{\dimexpr0.145\linewidth-6.4pt\relax}
                >{\raggedleft\arraybackslash}p{\dimexpr0.115\linewidth-6.4pt\relax}
                >{\raggedright\arraybackslash}p{\dimexpr0.37\linewidth-6.4pt\relax}
                >{\raggedright\arraybackslash}p{\dimexpr0.165\linewidth-6.4pt\relax}
                >{\raggedright\arraybackslash}p{\dimexpr0.205\linewidth-6.4pt\relax}@{}}
\toprule
Level & Number & What it partitions or represents & Where defined & How it is checked\\
\midrule
Parent rows & $58$ &
  Published implications: type $\mathfrak t$ and $\lambda_1\in[A_j,B_j]$ give $\lambda'\ge p_j$,
  $\lambda_2\ge r_j$ ($33$ $\mathrm{rr}$, $13$ $\mathrm{rc}$, $12$ $\mathrm{complex}$) &
  \cref{prop:parentrows} &
  Printed tables; Heath-Brown's Tables~4 and~7 recomputed\\
Base cells & $340$ &
  Parent intervals cut into cells of width $0.01$ ($145$, $89$, $106$) &
  \cref{subsec:sourcecover} &
  Validated in the replay\\
First-zero cells & $478$ &
  Tiling of the base cells ($195$, $115$, $168$), with bounds $p\le\lambda'$ and $l_2\le\lambda_2$ &
  \cref{subsec:sourcecover} &
  Validated in the replay\\
Gap cases & $684$ &
  $422$ cells without a gap; $262$ gaps for $\lambda'$ partitioning $56$ $\mathrm{complex}$ cells &
  \cref{subsec:sourcecover} &
  Validated in the replay\\
Specifications & $2768$ &
  Gap cases with a second-family record: $1083$ $\mathrm{rr}$ ($195$ unreserved, $444+444$
  reserved), $115$ $\mathrm{rc}$, $1570$ $\mathrm{complex}$ &
  \cref{def:specification} &
  Validated in the replay\\
Root cases & $4453$ &
  $2768$ specifications inside, and the $1685$ not of type $\mathrm{rr}$ also outside &
  \cref{thm:sourcecover} &
  Hand proof; replay\\
Refinement nodes & $1959$ &
  $1029$ splits, $236$ restricting location updates, $694$ exclusions ($75$ by real location updates, $422$ by
  lower bounds for $\lambda_2$ or $\lambda'$, $197$ by positivity); $4453+1029=4788+694$ &
  \cref{subsec:refinement}, \cref{tab:nodes} &
  Replay; some location rows by separate programs\\
Leaves & $4788$ &
  $3146$ inside ($233$ supplementary) and $1642$ outside; $343$ inside and $70$ outside roots have none &
  \cref{subsec:refinement} &
  Replay; Lean numeric checker\\
Certificate trees & $5591$ &
  One per leaf or subcase of a leaf: $3949$ inside ($349$ for supplementary leaves), $1642$ outside &
  \cref{subsec:subdivision} &
  Replay; independent checker; Lean checker\\
Threshold boxes & $4{,}196{,}879$ &
  Terminal boxes of the trees: $4{,}109{,}455$ inside ($1{,}042{,}363$ for supplementary leaves), $87{,}424$
  outside &
  \cref{def:certificate} &
  As above\\
Relaxation cases & $29{,}397{,}336$ &
  $2^k$ in a box with $k$ tangent rows; $9430$ closed by an exclusion, the rest by integer duals &
  \cref{def:certificate} &
  As above, in exact arithmetic\\
\midrule
\multicolumn{5}{@{}l}{\emph{Rows and exclusions}}\\
\midrule
Near rows & $16{,}842$ &
  All family, shifted and graded rows ($13{,}528$ inside, $3314$ outside), with $11{,}135{,}323$
  entries &
  \cref{subsec:rowkinds} &
  Interval arithmetic; Lean row checker\\
Two-test rows & $37$ &
  On $37$ inside roots: $31$ single, $2$ mixture, $4$ paired &
  \cref{subsec:twotest} &
  Replay; exact comparison\\
Complex location rows & $94$ &
  $25$ for $\lambda_2$ and $69$ for $\lambda'$ (\cref{prop:complexrows,prop:polyrows}) &
  Appendix~\ref{app:rows} &
  Interval arithmetic; separate program\\
Degree-$5$ real rows & $22$ &
  Cells covering $[0.700,0.755]$, used in $300$ location nodes &
  \cref{prop:realrows}, \cref{tab:realrows} &
  Interval arithmetic; replay\\
Positivity exclusions & $197$ &
  Reserved nodes of type $\mathrm{rr}$; least normalized margin $7.8\cdot10^{-5}$ &
  \cref{prop:positivity} &
  Interval arithmetic; replay\\
\bottomrule
\end{tabular}

\medskip
\begin{minipage}{\linewidth}
\footnotesize
Triples of numbers are counts for the types $\mathrm{rr}$, $\mathrm{rc}$, $\mathrm{complex}$;
``supplementary'' refers to the supplementary leaves of \cref{subsec:refinement}. Each split adds one
terminal node; hence the identity in the refinement row. The printed tables were compared with
the pages of \cite{HeathBrown1992zero,Xylouris2011nullstellen}, and Heath-Brown's Tables~4 and~7
were recomputed in floating point (\cref{subsec:locconventions}). The replay does not regenerate
the $94$ complex rows or the conditions \cite[(4.29), (4.34)]{Xylouris2011nullstellen}; separate
interval programs check them. The Lean row checker covers all near rows except the $37$
two-test rows.
\end{minipage}
\caption{The finite analysis of the middle range $0.1\le\lambda_1<1.5$ at a glance. The case
tree, from the parent rows to the leaves, partitions configurations of zeros; the boxes and
relaxation cases partition the thresholds of the near rows. The checks are described in
\cref{sec:verification}.}
\label{tab:census}
\end{table}

\subsection{The leaf model}\label{subsec:leafmodel}
This subsection gives the real coefficients before integer scaling. An interval
$[\lo,\hi)$ is charged at its left endpoint in the objective, where the cost is
largest, and at its right endpoint in density constraints, where the available
weight is smallest. These choices enlarge the feasible set and give an upper bound
for every actual distribution of characters in the interval.

Let $\mathfrak s$ be the specification of a leaf, after the location updates. Let $s$ denote the
\emph{shift} $\min(1.9,p^*,l_2)$, where $p^*=\lo'$ if $\mathfrak s$ has a gap and $p^*=p$ otherwise, and let
$\lambda_3^{\lo}$ be the lower bound for $\lambda_3$ of \cref{prop:third}, computed from $[a,b]$, the
type and, for a reserved $\mathfrak s$, the cap $\hi_2$. The leaf program (\cref{def:leafprogram}) has the
columns of \cref{tab:leafmodel}. Every number is rounded in the safe direction at scale $S$: far weights,
features and the hidden-tail count coefficient $1/w(\mathrm{end})$ down; objective coefficients,
diagonals, correlation terms, the far budget and the constant $\mathrm{first}$ up. Let $O$ be the set
of ordinary characters with a zero in $R_P$, and $\mathcal H$ the set of first-family characters with a
zero in $R_P$.
\begin{enumerate}[(a)]
  \item \emph{Ordinary bins.} Let $R=\max(3,r)$ and let $r=\xi_0<\xi_1<\dots<\xi_N=R$ consist of $r$,
    $R$ and the multiples of $1/\mathrm{den}$ between them, where $\mathrm{den}\in\{200,400,500,800,2000\}$
    is fixed per leaf. The bins are the intervals $[\xi_i,\xi_{i+1})$, except that in a reserved leaf the
    bins with $\xi_{i+1}\le\lambda_3^{\lo}$ are omitted. The ordinary characters beyond $R$ form the tail.
  \item \emph{Reserved family.} Either (fixed terms) the family is charged $n_2\mathcal G_{\phi_2}(\lo_2)$ in
    $\mathrm{first}$ and $n_2w(\hi_2)$ in the far budget, and it appears in the near rows as a family
    term with its feature at $\hi_2$ (there are then no second-family columns, and the count $n_E$ of the
    second-family equality of \cref{def:leafprogram} is $0$); or (columns) the interval $[\lo_2,\hi_2]$ is cut
    at the multiples of $1/400$ into reserved columns $[\lo_{2,j},\hi_{2,j}]$, with values $z_j$ and
    $\sum_jz_j=n_2=n_E$.
  \item \emph{Hidden columns} (outside leaves only). Put $p^\circ=\max(a,p,\lo')$ and
    $\mathrm{end}=\max(3,p^\circ)$. The hidden bins $[\lo_h,\hi_h)$ divide $[p^\circ,\mathrm{end}]$ on the grid
    $1/\mathrm{den}$, and the hidden tail collects the values $t_\chi\ge\mathrm{end}$. The hidden count is at
    most $n$.
\end{enumerate}

\begin{table}[ht]
\centering
\footnotesize
\setlength{\tabcolsep}{3.5pt}
\renewcommand{\arraystretch}{1.35}
\begin{tabular}{@{}>{\raggedright\arraybackslash}p{2.25cm}>{\raggedright\arraybackslash}p{3.3cm}cccccl@{}}
\toprule
Column & Value $x_c$ & $G_c$ & $W_c$ & $C_c$ & $N_c$ & $E_c$ & Near rows\\
\midrule
ordinary bin $[\xi_i,\xi_{i+1})$ & number of $\chi\in O$ with $\nu_{T^*}(\chi)$ in the bin & $\mathcal G(\xi_i)$ &
  $w(\xi_{i+1})$ & $0$ or $1$\textsuperscript{a} & $0$ & $0$ & entry (a)\\
tail & $\sum w(\nu_{T^*}(\chi))$ over $\chi\in O$ with $\nu_{T^*}(\chi)\ge R$ & $\mathcal G(R)/w(R)$ & $1$ & $0$ &
  $0$ & $0$ & omitted\\
reserved column\textsuperscript{b} $[\lo_{2,j},\hi_{2,j}]$ & $z_j=n_2$ if $\nu_*$ is in the column, else $0$ &
  $\mathcal G_{\phi_2}(\lo_{2,j})$ & $w(\hi_{2,j})$ & $0$ & $0$ & $1$ & entry (f)\\
hidden bin\textsuperscript{c} $[\lo_h,\hi_h)$ & number of $\chi\in\mathcal H$ with $t_\chi$ in the bin &
  $e^{-A\lo_h}B_{\phi_1}(p^\circ)$ & $w(\hi_h)$ & $0$ & $1$ & $0$ & entry (g)\\
hidden tail\textsuperscript{c} & $\sum w(t_\chi)$ over $\chi\in\mathcal H$ with $t_\chi\ge\mathrm{end}$ &
  $\dfrac{e^{-A\,\mathrm{end}}B_{\phi_1}(p^\circ)}{w(\mathrm{end})}$ & $1$ & $0$ & $\dfrac1{w(\mathrm{end})}$ & $0$ &
  omitted\\
\bottomrule
\end{tabular}
\caption{The columns of a leaf program and their real coefficients before scaling by $S$:
objective $G_c$, far weight $W_c$, count $C_c$, hidden count $N_c$ and second-family indicator
$E_c$ (\cref{def:leafprogram}). The last column gives the entry of \cref{tab:entries} that supplies
the feature and diagonal of the column in each near row. (a) $1$ if $\mathfrak s$ is unreserved
and $\xi_{i+1}\le\lambda_3^{\lo}$, and $0$ otherwise. (b) Only if the reserved family is charged
through columns. (c) Outside leaves only.}
\label{tab:leafmodel}
\end{table}
Before scaling and rounding, the far budget is $(1+\eta)V-\1_{\mathrm{inside}}\,n\,w(b)$,
with a further subtraction of $n_2w(\hi_2)$ for a fixed reserved-family term.
The stored integer budget $F$ is formed by rounding the initial budget up and each
subtraction down, as in \cref{sec:far}. The count budget is $2$, and the stored
allowance satisfies $\mathrm{final}/S=5\eta$.
For an inside leaf, $\mathrm{first}/S$ is rounded up from
$J_1=\min(J_{\mathrm{old}},J_{\mathrm{new}})$ with anchor $p^*$, omitting
$J_{\mathrm{new}}$ when $p^*<b$. For an outside leaf it starts at $0$.
The fixed reserved-family charge, if used, is added with upward rounding.
The near rows are those of \cref{sec:rows}.

\subsection{Subdivision of a leaf}\label{subsec:subdivision}
A leaf may be certified through finitely many subcases, each of which is certified with its own
leaf model. Besides the second-family split above, three kinds are used. The first replaces the
specification by three specifications. The other two restrict one further quantity $Q$ of the
configuration to one of finitely many intervals $I$ whose union is the range of $Q$; the
configuration set of the subcase $(\mathfrak s,I)$ is the set $C(\mathfrak s,I)$ of configurations for which, for
some admissible choice of $(\chi_1,\rho_1)$, of the minimizers and of $\rho'$, (C1)--(C5) hold and
$Q\in I$. Then $C(\mathfrak s)$ is the union of the sets $C(\mathfrak s,I)$, and all the results of
\cref{sec:rows,sec:casetree} that are stated for $C(\mathfrak s)$ hold for $C(\mathfrak s,I)$ with the same proofs.

\begin{lemma}[Reservation split]\label{lem:reservationsplit}
Let $\mathfrak s$ be unreserved with ordinary lower bound $r$, and let $r<h\le2$. Let $\mathfrak s_1$ and
$\mathfrak s_2$ be $\mathfrak s$ with the reservation $(h,1)$ and $(h,2)$ (so $\lo_2=r$ and $\hi_2=h$), and let $\mathfrak s_0$ be $\mathfrak s$ with
$r$ replaced by $h$. Then $C(\mathfrak s)\subseteq C(\mathfrak s_0)\cup C(\mathfrak s_1)\cup C(\mathfrak s_2)$.
\end{lemma}

\begin{proof}
If $\nu_*<h$, a family $\mathcal F\ne\mathcal F_1$ attaining $\nu_*$ exists, and it is a real character or a pair of
conjugate nonreal characters; the configuration lies in $C(\mathfrak s_1)$ or $C(\mathfrak s_2)$, since
$\nu_*\ge r$. Otherwise every family $\mathcal F\ne\mathcal F_1$ has $\nu(\mathcal F)\ge h$, and the configuration lies in
$C(\mathfrak s_0)$.
\end{proof}

In a reserved child the model uses only the minimality of the reserved family
(\cref{lem:familyfacts}(ii)) and the cap $\lambda_2\le h$ (\cref{lem:familyfacts}(iii)); it never uses
that the reserved family is $\mathcal F_2$. The split of \cref{lem:reservationsplit} is used with $h=\min(\lambda_3^{\lo},2)$ on four inside
roots and on five roots that contain supplementary leaves.

\emph{Height split for type $\mathrm{rc}$.} An inside leaf of type $\mathrm{rc}$ is split by $Q=\mu_1$
(recall that $\mu_1>0$ for this type) into $\mu_1\in[0,1]$ and $\mu_1\in[1,\infty)$, which changes only the
family row (\cref{sec:rows}).

\emph{Second-zero height split.} An inside leaf of type $\mathrm{complex}$ with a finite gap is
split by $Q=|y|$, where $y=\mu_1-\mu_{\rho'}=(\gamma_1-\gamma')\LL$ is the normalized height difference between
$\rho_1$ and the non-distinguished zero $\rho'$ of $\chi_1$ realizing $\lambda'$ (a zero of $\chi_1$, not of $\bchi_1$;
\cref{sec:notation}), into six pieces, according as $|y|$ lies in $[0,1]$, $[1,\frac32]$, $[\frac32,2]$, $[2,3]$,
$[3,5]$ or $[5,\infty)$. This changes only the family row. A piece may be certified by an exclusion (\cref{def:certificate}(i)) rather than by
duals: the near row alone is then infeasible.

\subsection{Realization and the cover theorem}\label{subsec:realization}

\begin{proposition}[Realization]\label{prop:realization}
Let $\mathfrak s$ specify a leaf or one of its subcases, with the program constructed
in \cref{subsec:leafmodel}. For every sufficiently large $q$ whose configuration lies
in $C(\mathfrak s)$, there is a feasible pair $(x,\tau)$ such that
\[
  W(q)+2\eta\le\mathcal V(x).
\]
The column values $x$ are obtained from the actual character counts and weighted tail
masses. The threshold $q_0$ may be chosen uniformly over the finite set of leaves and
subcases.
\end{proposition}

\begin{proof}
Let $O$ and $\mathcal H$ be as in \cref{subsec:leafmodel} ($\mathcal H$ is used only if the
configuration is outside). For
$\chi\in O$ we have $\nu_{T^*}(\chi)\le C_P$ and $\nu_{T^*}(\chi)\ge\nu(\chi)\ge r$ by (C4). Put $x_i$
equal to the number of $\chi\in O$ with $\nu_{T^*}(\chi)\in[\xi_i,\xi_{i+1})$, $x_{\mathrm{tail}}=\sum w(\nu_{T^*}(\chi))$
over $\chi\in O$ with $\nu_{T^*}(\chi)\ge R$, $z_j$ equal to $n_2$ for the reserved column containing
$\nu_*$ and $0$ otherwise, $x^{\mathrm h}_i$ equal to the number of $\chi\in\mathcal H$ whose value $t_\chi$
(\cref{subsec:firstoutside}) lies in the hidden bin $i$, and $x^{\mathrm h}_{\mathrm{tail}}=\sum w(t_\chi)$ over the
$\chi\in\mathcal H$ in the hidden tail. In a reserved leaf no $\chi\in O$ lies in an omitted bin, by \cref{lem:familyfacts}(ii) and
$\lambda_3\ge\lambda_3^{\lo}$.

\emph{Far.} The characters of $O$, the reserved characters, the first-family characters of $\mathcal
H$ and, for an inside configuration, $\chi_1$ and $\bchi_1$ are distinct, and to each we assign one
zero with $|\gamma|\le1$ and parameter at most $C_P$: the $T^*$-representative, the height-one
representative, the zero realizing $t_\chi$, and $\rho_1$ or $\overline{\rho_1}$, which has
$|\gamma_1|\le C_B/\LL\le1$. By \eqref{eq:farbudget}, and since $w$ is decreasing, the far
constraint holds.

\emph{Count.} By \cref{lem:familyfacts}(i), since a character in a counted bin has
$\nu_{T^*}(\chi)<\xi_{i+1}\le\lambda_3^{\lo}\le\lambda_3$.

\emph{Hidden count and second family.} We have $|\mathcal H|\le n$, a character in the hidden
tail contributes $w(t_\chi)/w(\mathrm{end})\le1$ to the hidden count, and $\sum_jz_j=n_2=n_E$ when the reserved
family is charged through columns (otherwise there are no second-family columns and $n_E=0$).

\emph{Near rows.} By \cref{prop:rowsvalid}, each near row holds for some threshold, with features at
least and diagonals at most those of the columns.

\emph{Value.} We use \eqref{eq:Wsplit}. An ordinary character in a bin costs at most the
objective $\mathcal G(\xi_i)$ at the left end of the bin, a reserved character in a column at most $\mathcal G_{\phi_2}$ at
the left end (with fixed terms, at most $\mathcal G_{\phi_2}(\lo_2)$), and the tail characters at most $\mathcal G(R)/w(R)$ per unit of far mass, by the monotonicity
of $\mathcal G$, $\mathcal G_{\phi_2}$ and $\mathcal G/w$ (\cref{lem:envelope}(v)). For an outside configuration we apply
\cref{prop:firstoutside} with the anchor $p^\circ=\max(a,p,\lo')$, which is admissible because every
non-distinguished occurrence of the first family has parameter at least $\lambda'\ge p^\circ$; then
$t_\chi\ge p^\circ$, a character in a hidden bin $[\lo_h,\hi_h)$ costs at most the objective $e^{-A\lo_h}B_{\phi_1}(p^\circ)$
since $e^{-A\lambda}$ is decreasing, and a character in the hidden tail costs at most
$e^{-A\cdot\mathrm{end}}B_{\phi_1}(p^\circ)/w(\mathrm{end})$ per unit of far mass, since $e^{-A\lambda}/w(\lambda)$ is decreasing
(\cref{lem:decay}). Including the directions of integer rounding, we obtain
\[
  W\le\mathcal V(x)-\frac{\mathrm{final}}S+E
    \le\mathcal V(x)-5\eta+2\eta
    \le\mathcal V(x)-2\eta.
\]
\end{proof}

\begin{theorem}[The zero-sum bound in the middle range]\label{thm:cover}
For the case tree and certificates described in this paper, there is a threshold
$q_0$ such that every $q\ge q_0$ with $0.1\le\lambda_1<1.5$ satisfies
\[
  W(q)<1-2\eta.
\]
More precisely, its configuration belongs to a leaf or subcase whose program has a
feasible pair $(x,\tau)$ with $W(q)+2\eta\le\mathcal V(x)<1$.
\end{theorem}

\begin{proof}
By \cref{thm:sourcecover} the configuration lies in the configuration set of a root. Following the
refinement tree and the subdivisions of \cref{subsec:subdivision}, each of which is exhaustive,
we reach a leaf or subcase whose configuration set contains the configuration; it is not an
excluded node, because excluded nodes have no configurations. \Cref{prop:realization} gives the
feasible pair, and \cref{thm:certificate} and the certificates of \cref{sec:verification} give
$\mathcal V(x)<1$.
\end{proof}

% !TEX root = ../linnik399.tex
\section{The exterior regimes}\label{sec:exteriors}

The case tree covers $0.1\le\lambda_1<1.5$. We now treat the remaining possibilities:
an empty zero rectangle, a first zero with $\lambda_1<0.1$, and a first zero with
$\lambda_1\ge1.5$. The small-parameter case is different from the others. Its first
zero is real: it is an exceptional zero in the sense of Landau and Page
\cite{Landau1918imaginar,Page1935number}. By the theorems of Siegel and Tatuzawa
\cite{Siegel1935classenzahl,Tatuzawa1951theorem} (see also
\cite{Pintz1977elementaryV,Pintz1977elementaryVIII}) it cannot be extremely close to $1$, but these
bounds do not keep $\lambda_1$ away from $0$. Although such a zero has a large cost, the
Deuring--Heilbronn phenomenon \cite{Deuring1933imaginare,Heilbronn1934class,Linnik1944deuring}
repels the remaining zeros. We use Heath-Brown's quantitative estimates for this repulsion (see
also \cite{BenliGoelTwissZaman2026explicit} for explicit versions) and, for a sufficiently close
zero, his separate theorem on Siegel zeros \cite{HeathBrown1990siegel}, by which a very strong
exceptional zero is even helpful.

\subsection{No zero in \texorpdfstring{$R(l)$}{R(l)}}
If no nonprincipal $L(s,\chi)$ vanishes in $R(l)$, then for large $q$ none vanishes in $R_P$, because
$C_P\le\frac13\log\log\LL$ and $C_P/\LL\le l$. So $W=0$ and \cref{prop:detect} gives a prime $p\equiv
a\pmod q$ with $q^A<p<q^L$. Xylouris notes this case himself \cite[p.~36]{Xylouris2011nullstellen}.

\subsection{A tiny exceptional zero}
Let $\eta_{\mathrm{HB}}(\delta)$ be the effectively computable constant of \cref{inp:siegel} and put
\[
  u_{\mathrm{exc}}=\min\Bigl(\frac1{\eta_{\mathrm{HB}}(1/2)},\,0.08\Bigr).
\]

\begin{proposition}[Primes in the presence of a very close real zero]\label{prop:tiny}
Let $u_{\mathrm{exc}}$ be the positive cutoff defined above. For all sufficiently large
$q$, if $\lambda_1<u_{\mathrm{exc}}$, then $P(a,q)\le q^{3.5}$ for every $(a,q)=1$.
\end{proposition}

\begin{proof}
By \cref{prop:firstlower}, $\chi_1$ and $\rho_1$ are real, so $\beta_0=\rho_1=1-\lambda_1/\LL$ is a real zero of the
real nonprincipal $L(s,\chi_1)$. Since $\lambda_1<u_{\mathrm{exc}}\le\frac13$ we have $\beta_0\ge1-1/(3\log q)$, and
$((1-\beta_0)\log q)^{-1}=1/\lambda_1>1/u_{\mathrm{exc}}\ge\eta_{\mathrm{HB}}(\frac12)$. \Cref{inp:siegel} with $\delta=\frac12$ gives the claim.
\end{proof}

\subsection{A small exceptional zero}
For $u_{\mathrm{exc}}\le\lambda_1\le0.1$ we use a different prime weight, a single triangle, and Heath-Brown's own
estimates. The inputs are the following.

\begin{inp}[Prime detection with a single triangle {\cite[Lemma~13.2]{HeathBrown1992zero}}]\label{inp:hb132}
Let $L'>2K+3$, let $f=h_{L',K}$ and $F(z)=e^{-(L'-2K)z}F_2(z)$ with $F_2(z)=\bigl(\frac{1-e^{-Kz}}z\bigr)^2$. For
every $\eps>0$ there is $R=R(\eps)$ such that for large $q$
\[
  \sum_{p\equiv a\,(q)}\frac{\log p}pf\Bigl(\frac{\log p}\LL\Bigr)\ge\frac\LL{\varphi(q)}\Bigl[K^2-\eps-\sum_{\chi\ne\chi_0}{\sum_\rho}'
  \bigl|F\bigl((1-\rho)\LL\bigr)\bigr|\Bigr],
\]
where $\sum'$ runs over the zeros with $1-R/\LL\le\beta\le1$, $|\gamma|\le R/\LL$.
\end{inp}

\begin{inp}[Per-character bound for the triangular kernel {\cite[Lemma~13.3]{HeathBrown1992zero}}]\label{inp:hb133}
Let $\chi\ne\chi_0$, $0<\lambda\le R$, and suppose that $L(s,\chi)$ has no zero with $1-\lambda/\LL<\beta\le1$ and
$|\gamma|\le1$. Then for every $\eta_1>0$ and $q\ge q(\eta_1,\lambda)$
\[
  {\sum_\rho}'\bigl|F_2\bigl((1-\rho)\LL\bigr)\bigr|\le\mathcal B_{\varphi(\chi)}(\lambda)+\eta_1,\qquad
  \mathcal B_\phi(\lambda)=\frac\phi2\,\frac{1-e^{-2K\lambda}}\lambda+\frac{2K\lambda-1+e^{-2K\lambda}}{2\lambda^2}.
\]
\end{inp}

\begin{inp}[Heath-Brown's weighted density estimate]\label{inp:hb111}
In the notation of \cite[Lemma~11.1]{HeathBrown1992zero} and
\cite[(5.18), p.~66]{Xylouris2011nullstellen}, the following holds.
Let $\eps,c_1,c_2>0$ and $\phi=\max_\chi\varphi(\chi)$. For $q\ge q(\eps,c_1,c_2)$ and any choice, for each
nonprincipal $\chi$ with a zero in $\beta\ge1-\lambda_0/\LL$, $|\gamma|\le1$ ($\lambda_0=\frac13\log\log\LL$), of one such zero
with parameter $\lambda(\chi)$,
\[
  \sum_\chi\frac{\lambda(\chi)}{e^{(4\phi+6c_1+2c_2)\lambda(\chi)}-e^{(2\phi+4c_1)\lambda(\chi)}}\le\frac{2\phi+2c_1+c_2}{4c_1c_2}+\eps .
\]
\end{inp}

The left side decreases and the right side increases with $\phi$, so the inequality holds with
$\phi=\frac13$. We also use \cref{inp:rrtables}(b),(d) and the following two rows of Heath-Brown's
tables, which hold for all $\lambda_1\le0.10$ bounded below by a fixed positive constant
\cite[\S8, Table~2 and Table~5]{HeathBrown1992zero} (see the remark after \cref{inp:rrtables}; we use them
for $\lambda_1\ge0.08$): if $\chi_1$ and
$\rho_1$ are real and $\lambda_1\le0.10$, then $\lambda'\ge4.96-\eps$ and $\lambda_2\ge2.83-\eps$ for $q\ge q(\eps)$. We have
recomputed these rows from Heath-Brown's printed parameters, obtaining $4.96355$ and $2.84577$; we use
them with $\eps=0.001$. (Table~5 is stated under $\lambda_2\le\lambda'$, which is redundant here since
$\lambda'\ge4.959>2.83$.)

\emph{Fixed data.} $L=3.99$, $K=0.1821$, $A_s=L-2K=3.6258$; $c_1=0.057$, $c_2=0.1554$, $\phi=\frac13$;
\[
  a_d=\tfrac43+6c_1+2c_2=\tfrac{3724}{1875},\quad b_d=\tfrac23+4c_1=\tfrac{671}{750},\quad
  V_d=\frac{2/3+2c_1+c_2}{4c_1c_2}=\frac{184750}{6993};
\]
$\alpha=1.09<\frac{12}{11}$, $B_1=K^2+\frac K4$, $m^*=\alpha\log12.5=2.7530442\ldots$, $m_2=2.829$ and $m_2'=4.959$.
Let $H_s(u)=e^{-A_su}F_2(u)$ and $\mathcal M(u)=(K^2-H_s(u))/u$, and
\[
  Q(m)=\frac{e^{-(A_s-a_d)m}-e^{-(A_s-b_d)m}}m=\int_{A_s-a_d}^{A_s-b_d}e^{-tm}\,dt .
\]

\begin{theorem}[The small exceptional-zero range]\label{thm:small}
For all sufficiently large $q$ with $u_{\mathrm{exc}}\le\lambda_1\le0.1$, and every
integer $a$ with $(a,q)=1$, there is a prime $p$ such that
\[
  p\equiv a\pmod q,\qquad q^{3.6258}<p<q^{3.99}.
\]
The threshold is uniform over the fixed interval $[u_{\mathrm{exc}},0.1]$.
\end{theorem}

\begin{proof}
Write $u=\lambda_1$. By \cref{prop:firstlower}, $\chi_1$ and $\rho_1$ are real. The zero $\rho_1$ is simple, since a
double zero would give $\lambda'=u$, whereas $\lambda'\ge(2-\eps)\log(1/u)>u$ (\cref{inp:rrtables}(b)) on
$[u_{\mathrm{exc}},0.08]$ and $\lambda'\ge4.959>u$ on $[0.08,0.1]$.

\emph{The criterion.} Apply \cref{inp:hb132} with $L'=L$, since $3.99>2K+3$; the zero $\rho_1$ contributes
$H_s(u)$. Enlarging $R$ only adds nonnegative terms to $\sum'$, so the inequality of \cref{inp:hb132}
remains true with $R$ replaced by $\max(R,3)$; we use it with $R\ge3$, so that \cref{inp:hb133} applies
with the values $\lambda\le m_2<3$ below. It remains to bound the other terms. Put $m(u)=\alpha\log(1/u)$ on $\mathcal P_1=[u_{\mathrm{exc}},0.08]$ and
$m(u)=m_2$ on $\mathcal P_2=[0.08,0.1]$, and $m_1(u)=m(u)$ on $\mathcal P_1$, $m_1(u)=m_2'$ on $\mathcal P_2$. By
\cref{inp:rrtables}(b),(d) (with $\eps=\frac{12}{11}-\alpha$) and the two table rows, every zero $\rho\ne\rho_1$ of
$L(s,\chi_1)$ in $R(l)$ has $\lambda_\rho\ge m_1(u)$, and every zero of every other nonprincipal $L(s,\chi)$ in
$R(l)$ has $\lambda_\rho\ge m(u)$.

\emph{The first character.} $L(s,\chi_1)$ has no zero with $\lambda<u_{\mathrm{exc}}/2$ and $|\gamma|\le1$, so
\cref{inp:hb133} with $\lambda=u_{\mathrm{exc}}/2$ and $\varphi(\chi_1)=\frac14$ gives
$\sum'_\rho|F_2((1-\rho)\LL)|\le\mathcal B_{1/4}(u_{\mathrm{exc}}/2)+\eta_1\le B_1+\eta_1$, as $\mathcal B_\phi$ is decreasing in $\lambda$, with limit $\phi K+K^2$ as $\lambda\to0$.
Hence the zeros of $\chi_1$ other than $\rho_1$ contribute
\[
  T_1\le e^{-A_sm_1(u)}(B_1+\eta_1).
\]

\emph{The other characters.} For $\chi\ne\chi_0,\chi_1$ with a zero in the rectangle of \cref{inp:hb132},
let $\lambda(\chi)$ be the parameter of its zero of least parameter with $|\gamma|\le1$; then $\lambda(\chi)\ge m(u)$.
Apply \cref{inp:hb133} with the fixed value $\lambda=m^*$ on $\mathcal P_1$ (note $m^*\le m(u)$ there) and $\lambda=m_2$
on $\mathcal P_2$, and with $\varphi(\chi)\le\frac13$. Since every zero in the rectangle has parameter at least
$\lambda(\chi)$, the zeros of $\chi$ contribute at most $e^{-A_s\lambda(\chi)}(\mathcal B_{1/3}(\bar m)+\eta_1)$, with $\bar m=m^*$
or $m_2$. Since $A_s>a_d$, the function $e^{-A_st}(e^{a_dt}-e^{b_dt})/t=Q(t)$ is decreasing, so
$e^{-A_s\lambda(\chi)}\le Q(m(u))\,\lambda(\chi)/(e^{a_d\lambda(\chi)}-e^{b_d\lambda(\chi)})$. By \cref{inp:hb111}, the characters
other than $\chi_1$ contribute
\[
  T_2\le(\mathcal B_{1/3}(\bar m)+\eta_1)(V_d+\eps_2)\,Q(m(u)).
\]

\emph{Monotonicity.} We have $K^2-H_s(u)=\int h_{L,K}(t)(1-e^{-ut})\,dt$, so
$\mathcal M(u)=\int h_{L,K}(t)\,\frac{1-e^{-ut}}u\,dt$, which decreases in $u$ because $h_{L,K}\ge0$ and
$(1-e^{-ut})/u$ decreases for each $t>0$. On $\mathcal P_1$ we have
$u=e^{-m/\alpha}$, and
\[
  \frac{Q(m(u))}u=\int_{A_s-a_d-1/\alpha}^{A_s-b_d-1/\alpha}e^{-tm(u)}\,dt,\qquad\frac{e^{-A_sm(u)}}u=e^{-(A_s-1/\alpha)m(u)},
\]
with $A_s-a_d-1/\alpha=\frac{236171}{327000}>0$; both increase with $u$. On $\mathcal P_2$ the bounds for $T_1$ and
$T_2$ do not depend on $u$, and $u\ge0.08$. Hence on $\mathcal P_i$
\[
  K^2-H_s(u)-T_1-T_2\ge\varpi_iu-\mathcal E,
\]
where
\begin{align*}
  \varpi_1&=\mathcal M(0.08)-V_d\,\mathcal B_{1/3}(m^*)\frac{Q(m^*)}{0.08}-B_1\frac{e^{-A_sm^*}}{0.08},\\
  \varpi_2&=\mathcal M(0.1)-V_d\,\mathcal B_{1/3}(m_2)\frac{Q(m_2)}{0.08}-B_1\frac{e^{-A_sm_2'}}{0.08},
\end{align*}
and $\mathcal E=\eta_1(1+(V_d+\eps_2)(a_d-b_d))+\eps_2\mathcal B_{1/3}(0)(a_d-b_d)$, using $Q\le a_d-b_d$.

\emph{Numerical margins.} The interval computations give the following values,
displayed here to ten decimal places:
\[
  \begin{array}{lcccc}
    & \mathcal M & \text{other characters} & \text{first character} & \varpi_i\\
    \mathcal P_1 & 0.1088458429 & 0.0783116746 & 0.0000454655 & 0.0304887029\\
    \mathcal P_2 & 0.1050056698 & 0.0668874692 & 0.0000000153 & 0.0381181853
  \end{array}
\]
The certified lower bounds give $\varpi_i>0.0304$. Choose $\eta_1$, $\eps_2$ and the $\eps$ of \cref{inp:hb132}, after $u_{\mathrm{exc}}$, with
$\mathcal E+\eps<0.03u_{\mathrm{exc}}$. Then the bracket in \cref{inp:hb132} is at least $(0.0304-0.03)u_{\mathrm{exc}}>0$, and there is a
prime $p\equiv a\pmod q$ with $\log p/\LL\in(A_s,L)$.
\end{proof}

At $L=3.99$, the margins are approximately $28\%$ and $36\%$ of the main term.
Exploratory evaluations with the same $K$, $c_1$, and $c_2$ remain positive at $L=3.90$;
the certified exterior result used here is the one at $3.99$.

\subsection{A large first zero}\label{subsec:large}

\begin{theorem}[The large first-zero range]\label{thm:large}
For all sufficiently large $q$ with a nonempty zero rectangle and $\lambda_1\ge1.5$,
\[
  W(q)\le0.9966814-2\eta.
\]
Consequently, every reduced residue class contains a prime $p$ with
$q^{3.156342}<p<q^{3.99}$.
\end{theorem}

\begin{proof}
Every nonprincipal character with a zero in $R_P$ is treated as ordinary, with its height-one
representative, of parameter $\lambda_\chi\ge\lambda_1\ge1.5$; by \cref{cor:ordinarycost} it contributes at most
$\mathcal G(\lambda_\chi)+\eps_1e^{-A\lambda_\chi}$. Consider the leaf program with the bins
$[1.5+\frac j{100},1.5+\frac{j+1}{100})$, $0\le j<150$, and the tail $[3,\infty)$; the far budget
$F=(1+\eta)V$ with profile~I (\cref{sec:far}); $\mathrm{first}=0$ and $\mathrm{final}=5\eta$; no count, hidden or second-family
constraints; and one near row. The near row is \cref{cor:zeroform} with $\Xi\equiv0$, $s=s_1=\frac32$,
detector $\hat f_2$ ($\gamma_f=1$), Gram test $\hat f_3$ ($\gamma_g=\frac32$) and no family entries, normalized as in
\cref{subsec:rowdata} with $t_0=1$ (so $I_u$ is the upper Riemann sum over $2000$ cells of each of
$[0,1)$ and $[1,2)$): every zero of every
nonprincipal $L$-function in $R(l)$ has parameter at least $1.5=s_1$, so the safe-anchor hypothesis
holds, and every entry has no zero to the right of its anchor $\min(\lo,s)=s$. The actual characters
give a feasible point as in \cref{prop:realization}.

The certificate consists of $36$ contiguous threshold intervals covering $[0,e]$; on each it uses the
first-order relaxation and integer duals $Y_F,Z\ge0$ satisfying \eqref{eq:dualfeasible}; there are no
exclusions and $5436$ column inequalities. The largest value is
$0.9966813270344512$, on the threshold interval $[0.1625,0.165]$, where $Y_F=5145926349$ and
$Z=1593393341176$. By \cref{thm:certificate}, $\mathcal V(x)<0.9966814$, and $W\le\mathcal V(x)-2\eta$ as in
\cref{prop:realization}. \Cref{prop:detect} gives the prime.
\end{proof}

In this regime one near row, together with the far-density budget, suffices.
All characters are treated by the same per-character bound.

% !TEX root = ../linnik399.tex
\section{Assembly of the proof}\label{sec:join}

Each preceding argument holds once $q$ exceeds a threshold depending on fixed choices.
We first verify that these choices can be made in a noncircular order. A single
threshold then works for every leaf and every exterior case.

\subsection{The order of the choices}\label{subsec:order}
The constants of the argument are chosen in the following order. Each stage depends only on the
earlier ones, and there are finitely many choices at each stage.
\begin{enumerate}[(S1)]
  \item The fixed data: the exponent $L=3.99$ and the kernel of \cref{sec:detection}; the two far
    profiles of \cref{sec:far}; the test functions, anchors, sieve weights and grids of all near
    rows (\cref{sec:rows}); the case tree, its leaf programs and their certificates
    (\cref{sec:casetree,sec:lp}); the zero-location rows of \cref{sec:location}; the data of the
    exterior regimes (\cref{sec:exteriors}); $\eta=\eta'=10^{-6}$, $\lambda_{11}=0.05$ and $s_{\max}=3$; and
    $u_{\mathrm{exc}}=\min(1/\eta_{\mathrm{HB}}(\frac12),0.08)$, with $\eta_{\mathrm{HB}}(\delta)$ from \cref{inp:siegel}; and the envelope tolerance
$\eps_1=\eta/23$ of \cref{subsec:allowance}.
  \item The prime window: $C_P\ge\max(C_0(\eta H_0),3)$, with $C_0$ from \cref{inp:criterion}.
  \item The zero count $N=N(C_P)$ of \cref{lem:zerocount}; the mesh $h_{\mathcal A}$ and the finite set
    of anchors $\mathcal A$ of \cref{subsec:zerocount}; then the smoothing parameter $\theta$ of
    \cref{subsec:smoothing}, chosen so that \eqref{eq:thetachoice} holds for $\eps_1=\eta/23$; then the
    constant $c_{\mathrm{out}}$ of \cref{prop:firstoutside}.
  \item The count $K_0$ of \cref{subsec:representatives}, a bound, uniform in $q$, for the number
    of zeros of all nonprincipal $L$-functions modulo $q$ with $\lambda\le s_{\max}$ and $|\gamma|\le2$
    (\cref{inp:logfree}). Then the separation $M$, so large that for $|Y|\ge M$ we have
    $|\Re F(x+iY)|\le\eta''_{\min}/(2K_0+2)$ for every detector and $x\in[-s_{\max},0]$, where
    $\eta''_{\min}>0$ is the least of the numbers $\eta''=\eta'\min(1,\sqrt{I_BD_u},D_u)$ of
    \cref{cor:zeroform} over the finitely many rows, and
    $|\Re G(-\sigma+iY)|\le d_1$ for every Gram test and every $\sigma=s_1-\delta-\delta'$ that occurs
    (\cref{lem:decaytransform}). Then the buffer $C_B>\max(C_P+M,\,1.5)$ with
    $4c_{\mathrm{out}}/(H_0C_B^2)\le\eta$.
  \item The tolerances of the published results and of the lemmas proved here, each below the
    finitely many positive margins it must fit: the tolerances of \cref{inp:X31,inp:X32} for the
    finitely many test functions (those of \cref{sec:costs}, the functions $f^\theta_\sigma$ with tolerance
    $H_0\eps_1/8$, and the detectors of the near rows with tolerance $\eta''_{\min}/2$), the disc radius $\delta_0\le\frac13$ of \cref{subsec:zerocount},
    which is at most the radius that \cref{inp:X32} provides for each of these, and the tolerances
    of \eqref{eq:X34} and \cref{lem:explicitphi}, of
    \cref{inp:far} (with $\eps=\eta J^2$), of \cref{inp:burgess,inp:graham} and Mertens' estimate, of
    \cref{inp:criterion} (with $\eps=\eta H_0$), of the zero tables, and of the small-branch
    estimates of \cref{thm:small} (with $\mathcal E+\eps<0.03u_{\mathrm{exc}}$).
  \item Finally $q_0$: the maximum of the finitely many thresholds of the previous stages, together
    with those of \cref{inp:rectangle}, of the inclusions $R_P\subseteq R_B\subseteq$ the discs, and
    $\LL>4M(K_0+1)$ for \cref{lem:Tstar}.
\end{enumerate}
The certificates of stage (S1) depend only on fixed data, not on $C_P$, $M$, $C_B$ or $q$. The
count $K_0$ counts zeros up to height $2$, which exceeds $1+\delta_0$ for every admissible radius;
so $M$, which depends on $K_0$, can be fixed before the radius $\delta_0$, and there is no circular
dependence. The strip height $T^*=T^*(q)\in[\frac12,1]$ of \cref{lem:Tstar} is then determined by $q$.

Envelope errors that are summed over unboundedly many characters are weighted by
$e^{-A\lambda_\chi}$, and $\sum_\chi e^{-A\lambda_\chi}<19$ by \eqref{eq:Kfar}; errors in the near rows are
uniform per entry or are multiples of $(\sum_ja_j)^2$ (\cref{thm:graded}). Thus the error
allowances remain uniform as the number of characters grows with $q$.

\subsection{The regimes}
Let $q\ge q_0$. Exactly one of the following holds.
\begin{enumerate}[(R1)]
  \item No nonprincipal $L(s,\chi)$ has a zero in $R(l)$.
  \item $\lambda_1<0.1$.
  \item $0.1\le\lambda_1<1.5$.
  \item $\lambda_1\ge1.5$.
\end{enumerate}
(The first zero $\lambda_1$ is defined when (R1) fails; since $R(l)$ is compact there are finitely
many zeros in it.)

\subsection{Proof of \texorpdfstring{\cref{thm:main}}{the main theorem}}
Let $q\ge q_0$ and let $a$ be coprime to $q$.

In case (R1), $W=0$ and \cref{prop:detect} gives a prime $p\equiv a\pmod q$ with $q^A<p<q^L$.

In case (R2), if $\lambda_1<u_{\mathrm{exc}}$ then $P(a,q)\le q^{3.5}$ by \cref{prop:tiny}; if $u_{\mathrm{exc}}\le\lambda_1<0.1$ then
\cref{thm:small} gives a prime $p\equiv a\pmod q$ with $q^{3.6258}<p<q^{3.99}$.

In case (R3), \cref{thm:cover} gives a leaf of the case tree, or a subcase of a leaf, and a
feasible point of its program with $W\le\mathcal V(x)-2\eta<1-2\eta$. In case (R4), \cref{thm:large}
gives $W<1-2\eta$. In both cases \cref{prop:detect} gives a prime $p\equiv a\pmod q$ with
$q^A<p<q^L$.

Thus $P(a,q)<q^{3.99}$ for every $q\ge q_0$ and every $(a,q)=1$.
There are only finitely many moduli $2\le q<q_0$, and finitely many reduced residue
classes for each such modulus. Dirichlet's theorem supplies a least prime in each
class. Taking the maximum of $1$ and the finitely many ratios $P(a,q)/q^{3.99}$ gives
an absolute $C$ valid for all $q\ge2$. \qed

\begin{remark}[Effectivity]
Heath-Brown's theorem on Siegel zeros \cite[Corollary~1]{HeathBrown1990siegel}, used for
the closest exceptional zeros, is effective. The other source estimates used here are
also stated with effective constants. However, we have neither audited the effectivity
of every step of the assembly nor computed a value of $q_0$ or $C$.
\Cref{thm:main} asserts their existence and does not supply an explicit numerical bound
for either constant.
\end{remark}

% !TEX root = ../linnik399.tex
\section{The computations and their verification}\label{sec:verification}

The finite computation has three logical components. The zero-location checks justify
the case analysis; enclosures of analytic quantities supply the leaf coefficients;
and integer certificates bound the resulting programs. The exterior ranges require
separate computations. A check of one component does not by itself verify the others.
This section records the evidence for each and the scope of the Lean formalization.

Computer assistance is well established in explicit analytic number theory, for instance in the
ternary Goldbach problem \cite{Helfgott2015ternary} and in verifications of the Riemann hypothesis
\cite{Platt2016numerical,PlattTrudgian2021riemann}, and linear programs with rigorously checked
bounds are central to the proof of the Kepler conjecture \cite{ObuaNipkow2009flyspeck,
SolovyevHales2011efficient,Hales2017formal}. In our computations, the real coefficients and
margins are enclosed by interval arithmetic \cite{Moore1966interval,MooreKearfottCloud2009introduction,
Rump2010verification,Tucker2011validated}. It is implemented with the interval context of mpmath
\cite{mpmath2023library}, at a working precision of at least $50$ digits, and with a vectorized
binary64 interval module that widens the result of every operation outward by one unit in the
last place and evaluates $\exp$, $\sin$ and $\cos$ by Taylor polynomials with explicit remainder
bounds, not by the platform's library; the Lean checkers of \cref{subsec:lean} use their own
verified rational interval arithmetic. Floating-point linear programming (with the HiGHS solver
\cite{HuangfuHall2018parallelizing} in SciPy \cite{Virtanen2020scipy}) is used only to propose dual
multipliers, and every certificate is checked in exact integer arithmetic. Independent
numerical audits use multiple-precision arithmetic \cite{mpmath2023library}.
The supplementary floating-point checks of certain published table entries are
identified separately in \cref{subsec:locconventions}. \Cref{tab:status} summarizes, for each
component of the proof, whether it is published or new, how it was obtained and how it was
checked.

\begin{table}[tp]
\centering
\footnotesize
\setlength{\tabcolsep}{4pt}
\begin{tabular}{@{}>{\raggedright\arraybackslash}p{\dimexpr0.25\linewidth-6.4pt\relax}
                >{\raggedright\arraybackslash}p{\dimexpr0.15\linewidth-6.4pt\relax}
                >{\raggedright\arraybackslash}p{\dimexpr0.11\linewidth-6.4pt\relax}
                >{\raggedright\arraybackslash}p{\dimexpr0.15\linewidth-6.4pt\relax}
                >{\raggedright\arraybackslash}p{\dimexpr0.34\linewidth-6.4pt\relax}@{}}
\toprule
Component & Source & Kind & Method & Checked by\\
\midrule
Positivity criterion and detection (\cref{inp:criterion,prop:detect}) &
  Published \cite{Xylouris2011nullstellen} & Analytic & -- &
  Lean, with the criterion as a hypothesis\\
Envelope and first-family bounds (\cref{sec:costs}) &
  New, after \cite{Xylouris2011nullstellen} & Analytic & Interval arithmetic &
  Hand proof; a hypothesis in Lean, which checks only the enclosures\\
Far density (\cref{inp:far}) &
  Published \cite{Xylouris2011nullstellen} & Analytic & -- &
  Lean hypothesis; profile conditions proved in Lean\\
Far constants $V$, $w$ (\cref{sec:far}) &
  New profiles & Numerical & Interval arithmetic &
  Multiple-precision audit; Lean numeric checker\\
Printed zero-location tables (\cref{prop:parentrows}) &
  Published \cite{HeathBrown1992zero,Xylouris2011nullstellen} & Analytic & As printed &
  Compared with the pages; Heath-Brown's Tables~4 and~7 recomputed in floating point\\
New location rows (\cref{sec:location}) &
  New & Both & Interval arithmetic &
  Replay; the $94$ complex rows by a separate program\\
Conditions \cite[(4.29), (4.34)]{Xylouris2011nullstellen} &
  Published; new check & Numerical & Interval arithmetic &
  Separate programs\\
Conductor refinement $\frac1{108}$ (\cref{subsec:conductor}) &
  New & Analytic & -- & Hand proof\\
Graded near lemma, response lemma, threshold form (\cref{sec:near}) &
  New & Analytic & -- & Lean (Comparator, two kernels)\\
Family, shifted and graded rows (\cref{sec:rows}) &
  New & Numerical & Interval arithmetic &
  Lean row checker ($16{,}842$ rows); multiple-precision audit\\
Two-test rows (\cref{subsec:twotest}) &
  New, after \cite{Xylouris2011nullstellen} & Both & Interval arithmetic &
  Hand proof; replay; exact comparison with a separate implementation on every certificate
  interval (constants: that implementation's interval code only)\\
Certificate soundness (\cref{thm:certificate}) &
  New & Analytic & -- & Lean (Comparator, two kernels)\\
Leaf certificates (\cref{def:certificate}) &
  New & Numerical & Exact integer &
  Replay; independent checker; Lean checker\\
Leaf data (\cref{subsec:leafmodel}) &
  New & Numerical & Interval arithmetic &
  Multiple-precision audits; Lean numeric checker (count-type integers: audit only)\\
Case tree and realization (\cref{sec:casetree}) &
  New & Analytic & -- & Hand proof; replay\\
Tiny exceptional zero (\cref{prop:tiny}); empty rectangle &
  Published \cite{HeathBrown1990siegel}; $W=0$ & Analytic & -- & As printed\\
Small exceptional zero (\cref{thm:small}) &
  New, from lemmas of \cite{HeathBrown1992zero} & Both & Interval arithmetic &
  Hand proof; exterior replay; margins recomputed independently\\
Large first zero (\cref{thm:large}) &
  New & Numerical & Exact integer &
  Replay's exact checker; independent checker; Lean certificate (native and kernel), numeric and
  row checkers\\
Assembly (\cref{sec:join}) &
  New & Analytic & -- & Hand proof\\
\bottomrule
\end{tabular}

\medskip
\begin{minipage}{\linewidth}
\footnotesize
The interval arithmetic is described at the beginning of this section. Floating point is used
only to propose parameters, dual multipliers and sieve heights, which need no justification, and
to recompute published tables.
\end{minipage}
\caption{Status of the components. Published inputs are used as printed, and those used in the
Lean theorems enter them as hypotheses; ``Both'' means an analytic argument with numerical
margins; ``hand proof'' means a conventional proof in this paper. The checks are described in
\cref{sec:verification}.}
\label{tab:status}
\end{table}

\subsection{The data}
The certificates are stored in three compressed JSON record files. The inside and
outside files are indexed by root case; the supplementary file records the $233$
supplementary leaves of $199$ inside roots (\cref{subsec:refinement}). Their sizes and SHA-256 hashes are:
\begin{center}
\small
\begin{tabular}{lr}
\toprule
File & Bytes \\
\midrule
\texttt{inside.jsonl.gz} & $298{,}536{,}973$ \\
\multicolumn{2}{l}{\scriptsize\texttt{539362a03e7b330dad03b33d6b5037871cfcbbf7b0973e7976fccf12509fe63c}}\\[2pt]
\texttt{nodes.jsonl.gz} & $159{,}624{,}716$ \\
\multicolumn{2}{l}{\scriptsize\texttt{b715ac735b6976da17a2671667ecfdf07c3f19b8ab74ae5f0533dbcf6b2a6cfc}}\\[2pt]
\texttt{outside.jsonl.gz} & $8{,}626{,}435$ \\
\multicolumn{2}{l}{\scriptsize\texttt{a53d4d5793f71c92647a7f22022110bdc8135cbda1f2575d1dc0c1f487f5ac1d}}\\
\bottomrule
\end{tabular}
\end{center}
The first holds the $2768$ inside roots with their $2913$ other leaves, the second the $233$
supplementary leaves, and the third the $1685$
outside roots with their $1642$ leaves. A leaf record contains the parameters of its near rows
and its certificate tree; it does \emph{not} contain the leaf program itself, which the verifier
regenerates from the case tree. Some leaves are split into subcases (\cref{subsec:subdivision}), each with
its own certificate tree, so the $4788$ leaves carry $5591$ certificate trees in all: $3949$
inside (including the supplementary leaves) and $1642$ outside.

\subsection{Exact verification}
All numbers entering a leaf program are enclosed with outward-rounded interval arithmetic and
stored as integers scaled by $S=10^{16}$, rounded in the safe direction (\cref{sec:lp}). The
verification of a certificate uses only integer arithmetic. The following checks have been run.
\begin{enumerate}[(1)]
  \item \emph{The replay.} A verifier traverses the case tree from its roots. It re-checks every
    split, location node and exclusion (regenerating the location rows, except that the table
    of complex polynomial rows and the side conditions \cite[(4.29), (4.34)]{Xylouris2011nullstellen} are
    checked by separate programs).
    At every leaf it regenerates the leaf program from the specification and the tree,
    independently of the stored certificate, and checks the certificate exactly. It passes on all
    $4453$ roots: $4788$ leaves, $4{,}196{,}879$ boxes, and largest certified value
    $0.9999998223\ldots<1$ (inside) and $0.9999942478\ldots$ (outside).
  \item \emph{An independent checker.} A second verifier was written from \cref{sec:lp} alone,
    without reading or using the checking code of the first, and it uses only exact integer and
    rational arithmetic. It takes the leaf data from the same generator as the replay (the leaf
    data are checked separately, in (4) and in \cref{subsec:lean}). It accepts every certificate of
    the corpus ($29{,}397{,}336$ relaxation cases in all) and that of \cref{thm:large}, and for every
    certificate tree its value and number of boxes are exactly those of the first verifier. In a
    negative test on ten leaves, including the tightest, it rejected all $308$ certificates made
    invalid by deliberate forgeries, among them duals lowered, and coefficients raised, by the
    smallest amount that breaks \eqref{eq:dualfeasible}; it accepted all $82$ controls perturbed by
    one unit less. It replaces an earlier independent checker, whose code was not preserved and
    which had also accepted every certificate of the corpus.
  \item \emph{A mutation test.} On the $19$ leaves of three of the tightest inside roots, $133$
    perturbations of the regenerated leaf data (far budget, objective, first-family charge,
    correlation term, features), of certificate multipliers and of the allowance $5\eta$, each in the
    unfavourable direction, were all rejected by the replay.
  \item \emph{Audits of the constants.} For $38$ leaves of all kinds, the $19{,}616$ objective and
    $19{,}616$ far coefficients, the tail, hidden and second-family columns, $F$, $\mathrm{first}$ and
    all $120$ family, shifted and graded near rows were recomputed from their definitions in
    multiple-precision arithmetic by a program sharing no code with the builders; all lie on the
    safe side of their true values. (The program of this audit was not preserved; the numeric and
    row checkers of \cref{subsec:lean} now check the same quantities for every leaf, except the
    two-test rows.) Separately, for $117$ roots ($100$ of them chosen at random),
    all $530{,}942$ column and budget integers of their $154$ leaves (objective, far, count,
    hidden-count and second-family coefficients, $F$, $\mathrm{first}$ and $\mathrm{final}$) were
    recomputed from their defining formulas in $50$-digit arithmetic and agree exactly with the
    stored ones.
  \item \emph{The two-test rows.} A separate program replays the $37$ roots whose leaves use the
    row of \cref{subsec:twotest} ($31$ with a single test, $2$ with a mixture, $4$ paired). It
    compares the row as the certificates use it with a separate implementation of the same
    inequality, with its own builder and its own certificate checker. The
    comparison runs on every threshold interval that the $37$ certificates use ($1079$ intervals,
    from $23{,}341$ boxes) and on $11{,}211$ seeded random intervals, in the first-order relaxation
    and in both tangent relaxations (for these, the tangent rule applied to the exact inequality of
    that implementation). The exact costs and budgets are those of that inequality divided by $D$,
    and the integers used are rounded in the safe direction. The row's integers are regenerated by
    the interval routines of that implementation and agree with the leaves; every one-unit
    perturbation of them is detected. These constants rest on those interval routines alone.
\end{enumerate}

\emph{The exterior regimes.} The margins of \cref{thm:small} were regenerated from their
definitions and checked in interval arithmetic by a separate replay of the exterior components,
and they have been recomputed independently in multiple-precision interval arithmetic. The
program of \cref{thm:large} is a leaf program in the sense of \cref{def:leafprogram}. Its
objective and far integers are regenerated and compared with the stored ones, and its near row is
built with the normalization of \cref{subsec:rowdata}. The exact checking code of the replay (1)
and the independent checker (2) accept its certificate, and so do the verified checkers of
\cref{subsec:lean}, which also accept its numeric data and its near row.

\subsection{Formal verification}\label{subsec:lean}
Several parts of the argument have been formalized in the Lean~4 proof assistant
\cite{deMouraUllrich2021lean} with the Mathlib library \cite{mathlib2020lean}; see
\cite{Avigad2024mathematics} for the role of such formalizations. Earlier formalizations in
analytic number theory include the prime number theorem
\cite{AvigadDonnellyGrayRaff2007formally,Harrison2009formalizing,KontorovichTao2024pnt} and
Dirichlet's theorem \cite{Stoll2024dirichlet}. The formalization uses Lean \texttt{v4.35.0-rc3} with
the corresponding Mathlib release, and contains about $24{,}000$ lines and about $800$ theorems and
lemmas. The $18$ statements of the near lemma and the certificate checker were checked against
their proofs with the Comparator tool \cite{LeanFRO2025comparator}, which replays the proofs in the
Lean kernel and in the independent kernel nanoda \cite{Bailey2020nanoda}; the leaf-level and near-row
theorems were checked by the Lean kernel. All proofs are complete and use only the
three standard axioms (propositional extensionality, the axiom of choice, and the soundness of
quotients). The published inputs of \cref{sec:inputs} enter as explicit hypotheses, each stated
no more strongly than its printed source. The following results are formally verified.
\begin{enumerate}[(a)]
  \item The graded near-density lemma (\cref{thm:graded}), from the local explicit formulas,
    Burgess's estimate and Graham's estimate as hypotheses. This includes the response lemma,
    the threshold form, and the entry bounds of \cref{sec:rows}.
  \item The soundness of the certificate checker of \cref{sec:lp}: an accepted certificate
    proves that every feasible point of the leaf program has value below $1$.
  \item The passage from a certified leaf with valid data to primes. From Xylouris's criterion
    and his far-density lemma, stated as printed, and the localized envelope of
    \cref{prop:envelope}, stated as a hypothesis, a leaf whose certificate and data pass the checks
    yields a prime $q^{A}<p<q^{L}$ in every reduced class, for every modulus whose zeros satisfy
    the leaf's zero-level hypotheses. The count, hidden-count and second-family integers and
    $\mathrm{final}$ enter as they are (they were recomputed in the audit of (4) above).
  \item Interval arithmetic with proved inclusion, including the exponential function, and
    enclosures of the functions entering the objective and far data of the leaves. A numeric
    checker is proved to imply that each leaf's objective and far coefficients, far budget and
    first-family charge are valid.
  \item Heath-Brown's Conditions~1 and~2 for the parabolic test functions of \cref{sec:rows}.
  \item The near rows: each family, shifted and graded near row is a concrete instance of
    \cref{thm:graded}. A row checker is proved sound, including the three family terms that keep a second zero and the constant
    $C_Z$ of \cref{sec:rows}. A checked row is a valid row of the leaf program under
    zero-level hypotheses only, and this is composed with (c). (The formal version of
    \cref{cor:zeroform} assumes that the numerators $D(\delta_j)-d_1+e_j$ themselves are positive,
    which is stronger than what the proof in \cref{sec:near} needs; the row checker verifies this
    for every checked entry.)
\end{enumerate}
The certificate checker has a Mathlib-free copy that is proved to compute the same value; the
numeric checker and the row checker are Mathlib-free themselves and are proved sound directly.
Run as compiled programs, the certificate and numeric checkers accept every certificate tree and
the objective, far and first-family data of every leaf in the corpus: $5591$ trees and
$4{,}196{,}879$ boxes.
The compiled row checker accepts every near row of the corpus except the $37$ two-test rows of
\cref{subsec:twotest}, which it does not cover: $16{,}842$ rows with $11{,}135{,}323$ entries. The same
three checkers accept the leaf program of \cref{thm:large}: its certificate ($36$ boxes), its
objective and far data, and its near row ($150$ entries). So (c) and (f) apply to it as to the leaves of the
case tree. In addition, seven certificates are checked by evaluation in the Lean kernel itself:
six of the case tree ($122$ boxes) and that of \cref{thm:large} ($36$ boxes).

The published inputs (\cref{sec:inputs,sec:location,sec:exteriors}) enter the formal
statements as hypotheses. The following parts have conventional proofs in this paper
but are \emph{not} formally verified. Their numerical checks are as described here and
in the relevant sections; in particular, a multiple-precision audit is distinct from
a formally verified enclosure.
\begin{itemize}
  \item the per-character costs of \cref{sec:costs}: the localized envelope, which enters as a
    hypothesis, and the first-family bounds (the enclosures of $J_{\mathrm{old}}$ and $J_{\mathrm{new}}$
    are verified, but not the inequalities they bound);
  \item the conductor coefficient $\frac14$ for real characters and the refinement of
    \cref{subsec:conductor};
  \item the zero-location arguments and the case tree (\cref{sec:location,sec:casetree}), that is, the
    statement that every configuration of zeros falls into a leaf whose zero-level hypotheses it
    satisfies, including the count, hidden-count and second-family constraints;
  \item the exterior regimes (\cref{sec:exteriors}): the range of \cref{prop:tiny,thm:small}, and,
    for \cref{thm:large}, the statement that every configuration with $\lambda_1\ge1.5$ satisfies the
    zero-level hypotheses of its program;
  \item the two-test row of \cref{subsec:twotest}, which is checked by the replay and compared
    exactly with a separate implementation in (5), and whose constants were computed only by the
    interval routines of that implementation;
  \item the count-type integers of the leaf programs (checked in the audit of (4));
  \item the assembly of \cref{sec:join}.
\end{itemize}
A compiled program is trusted to compute what it is proved to compute; its correctness then rests
on the Lean compiler and runtime and on the unverified routines that read the data files. The Lean
proofs were written by AI agents, like the rest of this work (see the note after the abstract).
The Lean kernel checks the proofs, but whether each formal statement says what the corresponding
statement of this paper says has so far been reviewed only by AI agents.

\subsection{Data availability}
The development repository contains the certificates, the programs that generate and
check them, the verification reports, and the Lean formalization:
\begin{center}\url{https://github.com/enaslund/linniks-constant}.\end{center}
At the date of this revision the repository is private. A public, versioned archive
of these materials is needed before submission so that readers can reproduce the
computations and identify the precise data used in the proof. This paper and the Lean
formalization are publicly available in the repository prepared for the Palomar registry of
formalized mathematics:
\begin{center}\url{https://github.com/enaslund/linniks-constant-3.99}.\end{center}

The repository includes instructions for rerunning the replay, mutation tests,
numerical audits, and formal checks. On the machine used,
the replay took $1.7$ hours of wall-clock time with parallel workers, and the compiled Lean checks
of all certificates, leaf data and near rows took about $5$ hours of wall-clock time. One caveat
concerns exact reproducibility: the sieve heights of the graded rows (\cref{subsec:rowdata}) are
regenerated in binary64 floating point and are not stored in the records. Any nonnegative heights
are admissible, so this does not affect soundness, but on a platform whose elementary functions
round differently the regenerated rows, and hence the certificates that use them, could differ and
would have to be recomputed. Thus the reproducibility requirement concerns the
particular stored certificates; admissibility of the analytic sieve weight alone does
not guarantee that a different regenerated row will pass those certificates.

% !TEX root = ../linnik399.tex
\section{Concluding remarks}\label{sec:remarks}

\emph{The difficult configurations.} The tightest certified cases have a nonreal first
character with $\lambda_1\in[0.64,0.80]$ and a nearby second family. The numerical
experiments of \cref{subsec:why} suggest that stronger location bounds in this range
would be useful. One limitation of a single near row can be seen directly: if $N$
entries have the same positive feature $v$ and diagonal $D$, then its quadratic
inequality permits $N\le D/(v^2-d)$ when $v^2>d$ (\cref{ex:twolevels}). This is a restriction of that
relaxation, not a lower bound on what other tests or arguments can achieve.

Improved far-density constants may also help. Xylouris estimates
\cite[p.~87]{Xylouris2011nullstellen} that halving the constant in his (5.19) would lower
the exponent obtained by his method to about $4.6$. A corresponding gain for the
present method would require new certificates for the complete cover.

\emph{Possible improvements.} Heath-Brown's list of possible improvements
\cite[\S16]{HeathBrown1992zero} contains several that we have not used: optimal test functions
for the combined inequalities (his item~1), positivity beyond the support of the test functions via
Brun--Titchmarsh bounds (item~3; see \cite{Titchmarsh1930divisor,MontgomeryVaughan1973large,
Iwaniec1982brun,Maynard2013brun}), and modified sieve coefficients in the density argument (item~6;
compare \cite{FriedlanderIwaniec2023selberg}). His item~7, a continuous weight in the density
argument, is what Xylouris's far-density lemma implements, and we use it. For cube-free moduli Burgess's bounds hold for every $k$, and even the
Weyl bound is known \cite{PetrowYoung2020weyl}; Heath-Brown already noted that all the arguments
improve in that case. Recent large-value estimates for Dirichlet polynomials
\cite{GuthMaynard2026new,ChenGuptaLi2025large,TaoTrudgianYang2026new} and new density theorems
\cite{Pintz2019some} improve zero-density estimates away from $\sigma=1$; it is not clear
whether they help in the range $\sigma=1-O(1/\log q)$ that matters here.

\emph{Other approaches.} For moduli with bounded cubic part the exponent $4.5$ has been claimed
\cite{Meng2010note}, and for moduli all of whose prime factors are small much better
exponents are known \cite{Chang2014short}. A different recent approach, through the exceptional set of Goldbach's problem (compare
\cite{Pintz2018new}), claims the exponent $5$ in a preprint \cite{Zhao2025exceptional}.

\emph{The role of the computation.} The graded lemma gives quadratic constraints for
different test functions and anchors. Their threshold forms can be combined with the
far-density budget in linear programs, whose bounds are certified by exact duality.
This separates the search for effective test parameters from the verification of the
resulting bound. Any improvement of the global exponent must still include the
zero-location arguments, a complete case cover, and the exterior regimes.

% !TEX root = ../linnik399.tex
\appendix
\section{The complex location rows}\label{app:rows}

The $94$ rows below make the implications of \cref{prop:complexrows,prop:polyrows}
explicit. In each row, the hypothesis is that $\chi_1$ is nonreal and
$\lambda_1\in[a,b]$. The number $h$ is the resulting strict lower bound for
$\lambda_2$ or $\lambda'$, as specified in the caption. The parameters $\gamma$ and
$t$ select the test function \eqref{eq:ftest} and the polynomial $P_t$.

Each margin is the interval-certified amount, divided by $f(0)$, by which the necessary
inequality fails at error tolerance $0$; the displayed margins are rounded for
readability. Their positivity allows a common sufficiently large $q$.
For the $\lambda_2$ rows, the real-second-character case uses the test parameter $\gamma_{\mathrm r}$ and $t'=0.85$;
the column labelled ``alias'' gives the margin in \eqref{eq:alias}.
The $\lambda'$ rows use the parent lower bound described after \cref{prop:polyrows}
and assume $\ord\chi_1\ge5$. For orders $3$ and $4$, the stronger published bound
quoted there applies.

\begin{center}
\scriptsize
\begin{longtable}{@{}cccccccc@{}}
\caption{Rows of \cref{prop:complexrows}: $\lambda_1\in[a,b]$ implies $\lambda_2>h$.}\label{tab:complexrows}\\
\toprule
$[a,b]$ & $h$ & $\gamma$ & $t$ & $\gamma_{\mathrm r}$ & generic & real & alias \\
\midrule
\endfirsthead
\toprule
$[a,b]$ & $h$ & $\gamma$ & $t$ & $\gamma_{\mathrm r}$ & generic & real & alias \\
\midrule
\endhead
\bottomrule
\endfoot
$[.66,.6625]$ & $.763$ & $1.3$ & $.95$ & $1.15$ & $0.00146$ & $0.0323$ & $0.1637$ \\
$[.6625,.665]$ & $.762$ & $1.3$ & $.95$ & $1.15$ & $0.00084$ & $0.0314$ & $0.1639$ \\
$[.665,.6675]$ & $.76$ & $1.3$ & $.9$ & $1.15$ & $0.00114$ & $0.0311$ & $0.1644$ \\
$[.6675,.67]$ & $.758$ & $1.3$ & $.9$ & $1.15$ & $0.00148$ & $0.0308$ & $0.1647$ \\
$[.67,.6725]$ & $.757$ & $1.3$ & $.9$ & $1.15$ & $0.00091$ & $0.0299$ & $0.1649$ \\
$[.6725,.675]$ & $.755$ & $1.3$ & $.9$ & $1.15$ & $0.00126$ & $0.0296$ & $0.1651$ \\
$[.675,.6775]$ & $.753$ & $1.3$ & $.9$ & $1.15$ & $0.00162$ & $0.0293$ & $0.1654$ \\
$[.6775,.68]$ & $.752$ & $1.3$ & $.9$ & $1.15$ & $0.00106$ & $0.0284$ & $0.1655$ \\
$[.68,.6825]$ & $.75$ & $1.3$ & $.9$ & $1.15$ & $0.00143$ & $0.0281$ & $0.1658$ \\
$[.6825,.685]$ & $.749$ & $1.3$ & $.9$ & $1.15$ & $0.00088$ & $0.0272$ & $0.1658$ \\
$[.685,.6875]$ & $.747$ & $1.3$ & $.9$ & $1.15$ & $0.00126$ & $0.0269$ & $0.1658$ \\
$[.6875,.69]$ & $.745$ & $1.3$ & $.9$ & $1.15$ & $0.00164$ & $0.0266$ & $0.1658$ \\
$[.69,.6925]$ & $.744$ & $1.3$ & $.9$ & $1.15$ & $0.00110$ & $0.0258$ & $0.1658$ \\
$[.6925,.695]$ & $.742$ & $1.3$ & $.9$ & $1.15$ & $0.00150$ & $0.0254$ & $0.1658$ \\
$[.695,.6975]$ & $.741$ & $1.3$ & $.9$ & $1.15$ & $0.00095$ & $0.0246$ & $0.1658$ \\
$[.6975,.7]$ & $.739$ & $1.3$ & $.9$ & $1.1$ & $0.00136$ & $0.0243$ & $0.1657$ \\
$[.7,.7025]$ & $.738$ & $1.3$ & $.9$ & $1.1$ & $0.00082$ & $0.0236$ & $0.1657$ \\
$[.7025,.705]$ & $.736$ & $1.3$ & $.9$ & $1.1$ & $0.00124$ & $0.0234$ & $0.1657$ \\
$[.705,.7075]$ & $.734$ & $1.3$ & $.9$ & $1.1$ & $0.00167$ & $0.0232$ & $0.1657$ \\
$[.7075,.71]$ & $.733$ & $1.3$ & $.9$ & $1.1$ & $0.00114$ & $0.0225$ & $0.1657$ \\
$[.71,.7125]$ & $.731$ & $1.3$ & $.9$ & $1.1$ & $0.00158$ & $0.0223$ & $0.1657$ \\
$[.7125,.715]$ & $.73$ & $1.3$ & $.9$ & $1.1$ & $0.00106$ & $0.0216$ & $0.1657$ \\
$[.715,.7175]$ & $.728$ & $1.25$ & $.95$ & $1.1$ & $0.00146$ & $0.0214$ & $0.1657$ \\
$[.7175,.72]$ & $.727$ & $1.25$ & $.95$ & $1.1$ & $0.00105$ & $0.0207$ & $0.1657$ \\
$[.72,.7225]$ & $.725$ & $1.25$ & $.95$ & $1.1$ & $0.00155$ & $0.0206$ & $0.1657$ \\
\end{longtable}
\end{center}

\begin{center}
\scriptsize
\begin{longtable}{@{}ccccc@{\qquad}ccccc@{}}
\caption{Rows of \cref{prop:polyrows}: $\lambda_1\in[a,b]$ implies $\lambda'>h$ (for $\chi_1$ of order at least $5$).}\label{tab:polyrows}\\
\toprule
$[a,b]$ & $h$ & $\gamma$ & $t$ & margin & $[a,b]$ & $h$ & $\gamma$ & $t$ & margin \\
\midrule
\endfirsthead
\toprule
$[a,b]$ & $h$ & $\gamma$ & $t$ & margin & $[a,b]$ & $h$ & $\gamma$ & $t$ & margin \\
\midrule
\endhead
\bottomrule
\endfoot
$[.68,.6825]$ & $.975$ & $1.15$ & $.95$ & $0.00179$ & $[.77,.7725]$ & $.908$ & $1.1$ & $.95$ & $0.00169$ \\
$[.6825,.685]$ & $.973$ & $1.15$ & $.95$ & $0.00175$ & $[.7725,.775]$ & $.906$ & $1.1$ & $.95$ & $0.00187$ \\
$[.685,.6875]$ & $.971$ & $1.15$ & $.95$ & $0.00171$ & $[.775,.7775]$ & $.904$ & $1.1$ & $.95$ & $0.00205$ \\
$[.6875,.69]$ & $.968$ & $1.15$ & $.95$ & $0.00228$ & $[.7775,.78]$ & $.903$ & $1.1$ & $.95$ & $0.00161$ \\
$[.69,.6925]$ & $.966$ & $1.15$ & $.95$ & $0.00225$ & $[.78,.7825]$ & $.901$ & $1.1$ & $.95$ & $0.00180$ \\
$[.6925,.695]$ & $.964$ & $1.15$ & $.95$ & $0.00222$ & $[.7825,.785]$ & $.899$ & $1.1$ & $.95$ & $0.00199$ \\
$[.695,.6975]$ & $.962$ & $1.15$ & $.95$ & $0.00220$ & $[.785,.7875]$ & $.897$ & $1.1$ & $.95$ & $0.00219$ \\
$[.7,.7025]$ & $.958$ & $1.15$ & $.95$ & $0.00216$ & $[.7875,.79]$ & $.896$ & $1.1$ & $.95$ & $0.00175$ \\
$[.7025,.705]$ & $.956$ & $1.15$ & $.95$ & $0.00215$ & $[.79,.7925]$ & $.894$ & $1.1$ & $.95$ & $0.00196$ \\
$[.705,.7075]$ & $.954$ & $1.15$ & $.95$ & $0.00214$ & $[.7925,.795]$ & $.892$ & $1.1$ & $.95$ & $0.00217$ \\
$[.7075,.71]$ & $.952$ & $1.15$ & $.95$ & $0.00213$ & $[.795,.7975]$ & $.891$ & $1.1$ & $.95$ & $0.00173$ \\
$[.71,.7125]$ & $.95$ & $1.15$ & $.9$ & $0.00214$ & $[.7975,.8]$ & $.889$ & $1.1$ & $.95$ & $0.00195$ \\
$[.7125,.715]$ & $.948$ & $1.15$ & $.9$ & $0.00218$ & $[.8,.8025]$ & $.887$ & $1.1$ & $.95$ & $0.00217$ \\
$[.715,.7175]$ & $.946$ & $1.15$ & $.9$ & $0.00222$ & $[.8025,.805]$ & $.886$ & $1.1$ & $.95$ & $0.00174$ \\
$[.7175,.72]$ & $.945$ & $1.1$ & $.95$ & $0.00170$ & $[.805,.8075]$ & $.884$ & $1.1$ & $.9$ & $0.00197$ \\
$[.72,.7225]$ & $.943$ & $1.1$ & $.95$ & $0.00181$ & $[.8075,.81]$ & $.883$ & $1.1$ & $.9$ & $0.00157$ \\
$[.7225,.725]$ & $.941$ & $1.1$ & $.95$ & $0.00193$ & $[.81,.8125]$ & $.881$ & $1.1$ & $.9$ & $0.00184$ \\
$[.725,.7275]$ & $.939$ & $1.1$ & $.95$ & $0.00205$ & $[.8125,.815]$ & $.879$ & $1.1$ & $.9$ & $0.00212$ \\
$[.7275,.73]$ & $.937$ & $1.1$ & $.95$ & $0.00217$ & $[.815,.8175]$ & $.878$ & $1.1$ & $.9$ & $0.00173$ \\
$[.73,.7325]$ & $.936$ & $1.1$ & $.95$ & $0.00170$ & $[.8175,.82]$ & $.876$ & $1.1$ & $.9$ & $0.00202$ \\
$[.7325,.735]$ & $.934$ & $1.1$ & $.95$ & $0.00183$ & $[.82,.8225]$ & $.875$ & $1.1$ & $.9$ & $0.00163$ \\
$[.735,.7375]$ & $.932$ & $1.1$ & $.95$ & $0.00196$ & $[.8225,.825]$ & $.873$ & $1.1$ & $.9$ & $0.00192$ \\
$[.7375,.74]$ & $.93$ & $1.1$ & $.95$ & $0.00210$ & $[.825,.8275]$ & $.872$ & $1.1$ & $.9$ & $0.00153$ \\
$[.74,.7425]$ & $.928$ & $1.1$ & $.95$ & $0.00223$ & $[.8275,.83]$ & $.87$ & $1.1$ & $.9$ & $0.00183$ \\
$[.7425,.745]$ & $.927$ & $1.1$ & $.95$ & $0.00177$ & $[.83,.8325]$ & $.868$ & $1.1$ & $.9$ & $0.00214$ \\
$[.745,.7475]$ & $.925$ & $1.1$ & $.95$ & $0.00192$ & $[.8325,.835]$ & $.867$ & $1.1$ & $.9$ & $0.00175$ \\
$[.7475,.75]$ & $.923$ & $1.1$ & $.95$ & $0.00207$ & $[.835,.8375]$ & $.865$ & $1.1$ & $.9$ & $0.00207$ \\
$[.75,.7525]$ & $.921$ & $1.1$ & $.95$ & $0.00222$ & $[.8375,.84]$ & $.864$ & $1.1$ & $.9$ & $0.00169$ \\
$[.7525,.755]$ & $.92$ & $1.1$ & $.95$ & $0.00176$ & $[.84,.8425]$ & $.862$ & $1.1$ & $.9$ & $0.00200$ \\
$[.755,.7575]$ & $.918$ & $1.1$ & $.95$ & $0.00192$ & $[.8425,.845]$ & $.861$ & $1.1$ & $.9$ & $0.00163$ \\
$[.7575,.76]$ & $.916$ & $1.1$ & $.95$ & $0.00208$ & $[.845,.8475]$ & $.859$ & $1.1$ & $.9$ & $0.00195$ \\
$[.76,.7625]$ & $.915$ & $1.1$ & $.95$ & $0.00163$ & $[.8475,.85]$ & $.858$ & $1.1$ & $.9$ & $0.00158$ \\
$[.7625,.765]$ & $.913$ & $1.1$ & $.95$ & $0.00179$ & $[.85,.8525]$ & $.856$ & $1.05$ & $.95$ & $0.00193$ \\
$[.765,.7675]$ & $.911$ & $1.1$ & $.95$ & $0.00196$ & $[.8525,.855]$ & $.855$ & $1.05$ & $.95$ & $0.00166$ \\
$[.7675,.77]$ & $.909$ & $1.1$ & $.95$ & $0.00214$ & & & & & \\
\end{longtable}
\end{center}

\section{The fixed constants}\label{app:constants}

This appendix collects the constants and parameters of the proof. The error budget that they
serve is summarized in \cref{subsec:budget}.

{\small
\setlength{\tabcolsep}{4pt}
\setlength{\LTcapwidth}{\linewidth}
\begin{longtable}{@{}>{\raggedright\arraybackslash}p{\dimexpr0.13\linewidth-6.4pt\relax}
                  >{\raggedright\arraybackslash}p{\dimexpr0.28\linewidth-6.4pt\relax}
                  >{\raggedright\arraybackslash}p{\dimexpr0.35\linewidth-6.4pt\relax}
                  >{\raggedright\arraybackslash}p{\dimexpr0.17\linewidth-6.4pt\relax}
                  >{\raggedright\arraybackslash}p{\dimexpr0.07\linewidth-6.4pt\relax}@{}}
\caption{Fixed constants and parameters. The last column gives the stage of the order of
choices of \cref{subsec:order} in which the constant is fixed; a dash marks a quantity computed
from earlier ones. Apart from the tolerances $\eps=\eta H_0$ and $\eps=\eta J^2$, the constants
of stages (S2)--(S6) are not computed; they exist by the arguments cited. Decimals are exact
rational numbers unless they end with an ellipsis.}
\label{tab:constants}\\
\toprule
Symbol & Value & Role & Where fixed & Stage\\
\midrule
\endfirsthead
\multicolumn{5}{@{}l}{\emph{\cref{tab:constants}, continued}}\\
\toprule
Symbol & Value & Role & Where fixed & Stage\\
\midrule
\endhead
\bottomrule
\endlastfoot
\multicolumn{5}{@{}l}{\emph{Exponent and kernel}}\\
$L$ & $3.99$ & The exponent & \cref{sec:notation} & (S1)\\
$T$ & $0.416829$ & Support of $\psi$ is $[0,T]$; $h$ is supported in $[L-2T,L]$ & \cref{subsec:kernel} & (S1)\\
$\kappa$ & $T/16=0.0260518125$ & Width of the $16$ steps of $\psi$ & \cref{subsec:kernel} & (S1)\\
$b_0,\dots,b_{15}$ & \cref{tab:beta} & Step heights of $\psi$ & \cref{tab:beta} & (S1)\\
$A$ & $L-2T=3.156342$ & Lower end of the support of $h$; primes are found in $(q^A,q^L)$; $A>3$ and $A>2x$ & \cref{subsec:kernel}, \cref{lem:decay} & --\\
$\Psi(0)$ & $\kappa\sum_ib_i=0.1936422928\ldots$ & Integral of $\psi$ & \cref{subsec:kernel} & --\\
$H_0$ & $\Psi(0)^2=0.0374973375\ldots$ & $H(0)$; normalizes the zero sum $W$ & \eqref{eq:Hkernel} & --\\
\midrule
\multicolumn{5}{@{}l}{\emph{Tolerances}}\\
$\eta$ & $10^{-6}$ & Detection needs $W\le1-2\eta$; far budget $(1+\eta)V$; allowance $5\eta$ & \cref{sec:notation}, \cref{prop:detect} & (S1)\\
$\eta'$ & $10^{-6}$ & Tolerance of the near rows: in every feature, diagonal and correlation term & \cref{subsec:rowdata} & (S1)\\
$\eps_1$ & $\eta/23$ & Envelope error $\eps_1e^{-A\lambda}$ per character; $23=19+4$, where $19>K_{\mathrm{far}}$ bounds $\sum_\chi e^{-A\lambda_\chi}$ over the ordinary characters and $4$ counts the first and reserved characters & \cref{subsec:allowance} & (S1)\\
$\mathrm{final}/S$ & $5\eta$ & Allowance included in every leaf value & \cref{subsec:leafmodel} & (S1)\\
$\eps$ in \cref{inp:criterion} & $\eta H_0$ & Tolerance of the positivity criterion & \cref{prop:detect} & (S5)\\
$\eps$ in \cref{inp:far} & $\eta J^2$ & Tolerance of the far-density lemma; gives \eqref{eq:farbudget} & \cref{sec:far} & (S5)\\
\midrule
\multicolumn{5}{@{}l}{\emph{Analytic constants}}\\
$\lambda_{11}$ & $0.05$ & Least anchor of the envelope; the anchor set $\mathcal A$ lies in $[\lambda_{11},C_P]$ & \cref{subsec:zerocount} & (S1)\\
$s_{\max}$ & $3$ & Largest anchor of an ordinary entry; $K_0$ counts zeros with $\lambda\le s_{\max}$ & \cref{subsec:representatives} & (S1)\\
$C_P$ & $\ge\max(C_0(\eta H_0),3)$ & Size of the prime window $R_P$; $C_0$ is the constant of \cref{inp:criterion} & \eqref{eq:CP} & (S2)\\
$N$ & $N(C_P)$ & Bound for the zeros of one character in $R_P$ & \cref{lem:zerocount} & (S3)\\
$h_{\mathcal A}$, $\mathcal A$ & Mesh with $|B_\phi(x)-B_\phi(y)|\le\eps_1/4$ for $|x-y|\le h_{\mathcal A}$ & Finite anchor set: the grid $\lambda_{11}+h_{\mathcal A}\ZZ$ and the anchors $a$, $p^*$, $p^\circ$ of the leaves & \cref{subsec:zerocount} & (S3)\\
$\theta$ (smoothing) & In $(0,1)$, with \eqref{eq:thetachoice} & $L^1$-distance of the smooth step function $\psi^\theta$ from $\psi$ & \cref{subsec:smoothing} & (S3)\\
$c_{\mathrm{out}}$ & $C_Pf^\theta_p(0)+|(f^\theta_p)'(0)|+e^{C_PT}\|(f^\theta_p)''\|_1$, maximized over the anchors $p$ & Cost of the distinguished terms of an outside first family & \cref{prop:firstoutside} & (S3)\\
$K_0$ & Independent of $q$ & Bound for the zeros with $\lambda\le s_{\max}$ and $|\gamma|\le2$ & \cref{subsec:representatives} & (S4)\\
$M$ & Large & Height separation for \cref{lem:Tstar,lem:response}; depends on $K_0$ and $\eta''_{\min}$ & \cref{subsec:order} & (S4)\\
$C_B$ & $>\max(C_P+M,1.5)$, with $4c_{\mathrm{out}}/(H_0C_B^2)\le\eta$ & Size of the buffer $R_B$, which separates inside from outside & \cref{sec:notation}, \cref{subsec:order} & (S4)\\
$\delta_0$ & In $(0,\frac13]$ & Common disc radius of the explicit formulas & \cref{subsec:zerocount} & (S5)\\
$q_0$ & Not computed & Maximum of the finitely many thresholds & \cref{subsec:order} & (S6)\\
\midrule
\multicolumn{5}{@{}l}{\emph{Far density} (profile~I; profile~II)}\\
$J$, $\eps_0$ & $10$, $10^{-7}$ & Number of weights $\alpha_i$; offset in $w_0$ (both profiles) & \cref{sec:far} & (S1)\\
$c_1$ & $0.09035$; $0.0821922$ & Parameter of \cref{inp:far} & \cref{sec:far} & (S1)\\
$c_2$ & $0.235968$; $0.2170903$ & Parameter of \cref{inp:far} & \cref{sec:far} & (S1)\\
$\theta_{\mathrm{far}}$ & $1.28683$; $1.4964274$ & Exponent in $w_0(t)^2$ & \cref{sec:far} & (S1)\\
$\alpha_1,\dots,\alpha_{10}$ & \cref{tab:alpha} & Weights of \cref{inp:far}, of sum $1$ & \cref{tab:alpha} & (S1)\\
$x$ & $1.1736846\ldots$; $1.1303335\ldots$ & Upper end $\frac23+3c_1+c_2$ of the profile; $2x<A$ & \cref{lem:decay} & --\\
$V$ & $175.26640331\ldots$; $243.32098105\ldots$ & Right side of \cref{inp:far} with $\eps=0$; far budget $(1+\eta)V$ & \cref{sec:far} & --\\
$w(0)$ & $10.7054\ldots$; $13.1730\ldots$ & Far weight at $\lambda=0$ & \cref{sec:far} & --\\
$K_{\mathrm{far}}$ & $(1+\eta)V/w(0)$: $<16.38$; $<18.48$ & Bound for $\sum_\chi e^{-A\lambda_\chi}$ & \eqref{eq:Kfar} & --\\
\midrule
\multicolumn{5}{@{}l}{\emph{Leaf programs and certificates}}\\
$S$ & $10^{16}$ & Scale of the stored coefficients & \cref{sec:notation}, \cref{subsec:leafdata} & (S1)\\
$T_S$ & $10^{6}$ & Thresholds are measured in ticks of $1/T_S$ & \cref{sec:lp} & (S1)\\
$D_S$ & $10^{12}$ & Scale of the dual multipliers & \cref{sec:lp} & (S1)\\
$\mathrm{den}$ & One of $200$, $400$, $500$, $800$, $2000$ per leaf & Ordinary bins between $r$ and $R=\max(3,r)$ on the grid $1/\mathrm{den}$ & \cref{subsec:leafmodel} & (S1)\\
Reserved grid & $1/400$ & Grid of the columns of a reserved second family & \cref{subsec:leafmodel} & (S1)\\
\midrule
\multicolumn{5}{@{}l}{\emph{Near rows}}\\
$\eps'$ & $1/2000$ & Sieve levels $\vartheta_k=\frac12(t_k-\frac13)-\eps'$ & \cref{subsec:rowdata} & (S1)\\
Cells & $2000$ sieve cells; $4000$ cells for $I_u$ & Equal cells $t_0<\dots<t_{2000}=2\gamma_f$; upper Riemann sum for $I_B$ & \cref{subsec:rowdata} & (S1)\\
$h_k$, $\mu_{\mathrm{row}}$ & Rational, denominators at most $10^{12}$; $\mu_{\mathrm{row}}=10^9$ in family and shifted rows & Sieve heights; they vanish in family and shifted rows & \cref{subsec:rowdata} & (S1)\\
$\lfloor\delta\rfloor_{1/200}$ & Grid $1/200$ & Offsets at which the diagonal $D(\delta)$ is evaluated & \cref{subsec:rowdata} & (S1)\\
$s_1$ & $\lfloor50a\rfloor/50$ & Safe anchor of a leaf with first-zero cell $[a,b]$ & \cref{subsec:leaffacts} & (S1)\\
$s$ (shift) & $\min(1.9,p^*,l_2)$ & Anchor of the shifted row and of the two-test row & \cref{subsec:rowkinds}, \cref{subsec:leafmodel} & (S1)\\
Graded anchors & $1.5$, $1.6$, $1.9$, $2.1$, $2.3$ & Reference anchors $s$ of the graded rows used & \cref{subsec:rowkinds} & (S1)\\
\midrule
\multicolumn{5}{@{}l}{\emph{Exterior regimes}}\\
$u_{\mathrm{exc}}$ & $\min\bigl(1/\eta_{\mathrm{HB}}(\frac12),0.08\bigr)$ & Below it, \cref{prop:tiny} ($P(a,q)\le q^{3.5}$); not computed & \cref{sec:exteriors} & (S1)\\
$K$, $A_s$ & $0.1821$; $L-2K=3.6258$ & Single triangle $h_{L,K}$ of the small branch & \cref{thm:small} & (S1)\\
$c_1$, $c_2$, $\phi$ & $0.057$, $0.1554$, $\frac13$ & Parameters of \cref{inp:hb111} & \cref{thm:small} & (S1)\\
$a_d$, $b_d$, $V_d$ & $\frac{3724}{1875}$, $\frac{671}{750}$, $\frac{184750}{6993}$ & Exponents and constant of \cref{inp:hb111} & \cref{thm:small} & --\\
$\alpha$ & $1.09<\frac{12}{11}$ & Exponent in the bounds $\lambda',\lambda_2\ge\alpha\log(1/\lambda_1)$ on $[u_{\mathrm{exc}},0.08]$ & \cref{thm:small} & (S1)\\
$m^*$ & $\alpha\log12.5=2.7530442\ldots$ & Anchor for the other characters on $[u_{\mathrm{exc}},0.08]$ & \cref{thm:small} & --\\
$m_2$, $m_2'$ & $2.829$, $4.959$ & Bounds $\lambda_2\ge2.83-\eps$, $\lambda'\ge4.96-\eps$ with $\eps=0.001$, on $[0.08,0.1]$ & \cref{thm:small} & (S1)\\
$B_1$ & $K^2+\frac K4$ & Bound for $\mathcal B_{1/4}$ for the first character & \cref{thm:small} & --\\
Large branch & $s=s_1=\frac32$, $\gamma_f=1$, $\gamma_g=\frac32$, $t_0=1$ & One near row; bins of width $\frac1{100}$ on $[1.5,3)$ and a tail & \cref{thm:large} & (S1)\\
\end{longtable}
}


\bibliographystyle{amsplain}
\bibliography{references}

\end{document}
